% Toolbox appendix for "High-Dimensional Gaussian Approximations over
% Rectangles without Covariance Restrictions" (draft of 26 October 2026).
%
% This file is the body of a stand-alone appendix. Rectangle_CLT_Toolbox.tex
% compiles it alone. To include it in Rectangle_CLT_Main-reworking.tex, add
%
%     \def\toolboxinpaper{}
%     % Toolbox appendix for "High-Dimensional Gaussian Approximations over
% Rectangles without Covariance Restrictions" (draft of 26 October 2026).
%
% This file is the body of a stand-alone appendix. Rectangle_CLT_Toolbox.tex
% compiles it alone. To include it in Rectangle_CLT_Main-reworking.tex, add
%
%     \def\toolboxinpaper{}
%     % Toolbox appendix for "High-Dimensional Gaussian Approximations over
% Rectangles without Covariance Restrictions" (draft of 26 October 2026).
%
% This file is the body of a stand-alone appendix. Rectangle_CLT_Toolbox.tex
% compiles it alone. To include it in Rectangle_CLT_Main-reworking.tex, add
%
%     \def\toolboxinpaper{}
%     % Toolbox appendix for "High-Dimensional Gaussian Approximations over
% Rectangles without Covariance Restrictions" (draft of 26 October 2026).
%
% This file is the body of a stand-alone appendix. Rectangle_CLT_Toolbox.tex
% compiles it alone. To include it in Rectangle_CLT_Main-reworking.tex, add
%
%     \def\toolboxinpaper{}
%     \input{Rectangle_CLT_Toolbox_appendix}
%
% after Appendix E and before \end{document}. The section then becomes
% Appendix F, and every reference to the paper below becomes a \ref to the
% paper's own labels. In the stand-alone version the same references print
% the numbers that the labels have in the paper on main (commit 12b6efd).
% The file defines the environment `fact` if the host does not, and defines
% its macros with \providecommand, so it leaves the host's definitions as
% they are. The list of sources at the end can be merged into the paper's
% bibliography.

\makeatletter
\providecommand{\E}{\mathbb E}
\providecommand{\Pp}{\mathbb P}
\providecommand{\R}{\mathbb R}
\providecommand{\tr}{\operatorname{tr}}
\providecommand{\op}{\mathrm{op}}
\providecommand{\norm}[1]{\left\lVert#1\right\rVert}
\providecommand{\ind}{\mathbf 1}
\providecommand{\dd}{\,d}
\providecommand{\Hh}{H_2}
\providecommand{\surf}{\mathcal J}
\providecommand{\relint}{\operatorname{relint}}
\providecommand{\aff}{\operatorname{aff}}
\providecommand{\cone}{\operatorname{cone}}
\providecommand{\dist}{\operatorname{dist}}
\providecommand{\prj}{\operatorname{pr}}
\providecommand{\Haus}{\mathcal H}
\providecommand{\TBc}[1]{[#1]}
% \PN{label}{C.5} prints a number of the paper (a \ref in the paper);
% \PR{Lemma}{label}{C.5} prints "Lemma C.5"; \PE{label}{C.21} an equation.
\providecommand{\PN}[2]{\ifdefined\toolboxinpaper\ref{#1}\else#2\fi}
\providecommand{\PR}[3]{#1~\PN{#2}{#3}}
\providecommand{\PE}[2]{(\PN{#1}{#2})}
\providecommand{\usedin}[1]{\par\smallskip\noindent{\small\emph{Used in:} #1}\par\smallskip}
\@ifundefined{c@fact}{\newtheorem{fact}{Fact}[section]}{}
\makeatother

\section{Toolbox: facts from analysis and geometry}\label{app:toolbox}

This appendix collects the facts from analysis, convex geometry, and
measure theory on which the proofs in Appendices~\PN{appx-sec:smoothing}{B}
and~\PN{appx-sec:support}{C} rely and which a mathematical statistician may
not have at hand. Every surface in those proofs is a finite union of flat
pieces of hyperplanes, so no differential geometry beyond the flat case is
needed. On a hyperplane the area formula is Fubini's theorem, the
Gauss--Green formula (the divergence theorem) is the fundamental theorem of
calculus along parallel lines, and the coarea formula is Fubini's theorem
applied piece by piece; the proofs below carry out each of these
reductions. Each fact is stated in the form in which
the proofs use it, and is followed by the places that use it and by a proof
or a reference. A proof is given when it takes a few lines, and a standard
source, listed at the end of the appendix, is cited otherwise. The facts are
standard results and not claims of the paper, so they carry their own
numbering. The symbols in the lines headed ``Used in'' are those of the
paper's appendices.

Subsection~\ref{tb:sec:calc} treats integration and calculus,
Subsection~\ref{tb:sec:gauss} the Gaussian identities, and
Subsection~\ref{tb:sec:matrix} the matrix and entropy inequalities.
Subsection~\ref{tb:sec:logconcave} covers log-concavity and
Pr\'ekopa's theorem. Subsections~\ref{tb:sec:poly}--\ref{tb:sec:coarea}
treat polyhedra, from faces and normal cones to the divergence theorem,
tubes, and the coarea formula, in the order in which
\PR{Lemma}{appx-an:lem:surface}{C.5} needs them, and
Subsection~\ref{tb:sec:index} lists where each fact is used.

\subsection{Notation}\label{tb:sec:notation}

The space \(\R^n\) carries the inner product \(x\cdot y\) and the Euclidean
norm \(\norm x\); in the proofs \(n\) is the rank \(r\) of the covariance,
or \(r-1\) when a hyperplane of \(\R^r\) is identified with
\(\R^{r-1}\). We write \(dx\) for Lebesgue measure, \(\varphi_n\) and
\(\Phi_n\) for the standard Gaussian density and law on \(\R^n\), and \(G\)
for a standard Gaussian vector. For a linear subspace \(L\), \(\varphi_L\)
is the standard Gaussian density on \(L\) with respect to Lebesgue measure
on \(L\). For a set \(S\), \(\ind_S\) is its indicator function, \(\aff S\)
is the smallest affine subspace containing it, the dimension of \(S\) is
that of \(\aff S\), and \(\relint S\) is the interior of \(S\) relative to
\(\aff S\). The distance from a point \(y\) to a set \(S\) is
\(\dist(y,S)=\inf_{x\in S}\norm{y-x}\).

For \(f\in C^q(\R^n)\), the derivative of order \(q\) in the directions
\(z_1,\ldots,z_q\in\R^n\) is
\[
 D^qf(x)[z_1,\ldots,z_q]
 =\partial_{s_1}\cdots\partial_{s_q}f(x+s_1z_1+\cdots+s_qz_q)\Big|_{s=0}
 =\sum_{i_1,\ldots,i_q}\partial_{i_1\cdots i_q}f(x)\,
   z_{1,i_1}\cdots z_{q,i_q},
\]
which is symmetric and multilinear in the directions. We write
\(D^qf(x)[z^{\otimes q}]\) when all directions equal \(z\), and
\(D^1f=\nabla f\). The Loewner order \(M\preceq N\) means that
\(N-M\) is positive semidefinite, \(\norm\Xi_\op\) is the operator norm
of a matrix \(\Xi\), and \(\norm M_*\) is the \emph{nuclear norm}, which for
a symmetric matrix \(M\) with eigenvalues \(\mu_i\) equals \(\sum_i|\mu_i|\).

\subsection{Integration and calculus}\label{tb:sec:calc}

\begin{fact}[Fubini and Tonelli]\label{tb:fubini}
Let \((S_1,\mu_1)\) and \((S_2,\mu_2)\) be \(\sigma\)-finite measure
spaces and let \(f\) be a measurable function on \(S_1\times S_2\). If
\(f\ge0\), then \(s_1\mapsto\int f(s_1,s_2)\,\mu_2(ds_2)\) is measurable
and
\[
 \int\Bigl(\int f\,d\mu_2\Bigr)d\mu_1
 =\int\Bigl(\int f\,d\mu_1\Bigr)d\mu_2
 =\int f\,d(\mu_1\otimes\mu_2)\in[0,\infty].
\]
If \(\int|f|\,d(\mu_1\otimes\mu_2)<\infty\), the same equalities hold, and
the inner integrals are finite almost everywhere.
\end{fact}

Three special cases occur in the proofs. If \(X\) and \(Y\) are
independent random elements and either \(f\ge0\) or \(f(X,Y)\) is
integrable, then
\[
 \E f(X,Y)=\E g(X),\qquad g(x)=\E f(x,Y),
\]
because the joint law is the product of the marginal laws. This is the
formal content of every argument that is carried out for a fixed value of
an independent component such as \(V_0\) and then averaged over it. With
\(\mu_2\) the counting measure, the theorem says that
\(\E\sum_k\xi_k=\sum_k\E\xi_k\) for nonnegative terms, and for terms with
\(\sum_k\E|\xi_k|<\infty\). For an orthogonal decomposition
\(\R^n=L\oplus L^\perp\), Lebesgue measure on \(\R^n\) is the product of the
Lebesgue measures on \(L\) and on \(L^\perp\), so that
\(\int_{\R^n}f\,dx=\int_{L^\perp}\int_Lf(l+w)\,dl\,dw\).

The appeal to Tonelli's theorem in Step~3 of the proof of
\PR{Lemma}{appx-an:lem:surface}{C.5}, whose notation we use in this
paragraph, is of the first kind. With \(V_0\) and
\(G\) independent, \(\Phi_r(\mathcal O_\iota(V_0))\) is the nonnegative
integral \(\int\ind\{g\in\mathcal O_\iota(V_0)\}\varphi_r(g)\,dg\), so
\(\E\sum_\iota R_\iota^3\,\Phi_r(\mathcal O_\iota(V_0))\) equals
\(\int\sum_\iota R_\iota^3\,\Pp\{g\in\mathcal O_\iota(V_0)\}\varphi_r(g)\,dg\)
with no check of integrability, since every term is nonnegative, and the
probability can then be bounded for each fixed \(g\).

\usedin{conditioning on \(V_0\) or on the \(\sigma\)-field \(\mathcal H\) in
\PR{Lemmas}{appx-an:lem:restoration}{C.2}, \PN{appx-an:lem:matrix}{C.4}
and~\PN{appx-an:lem:surface}{C.5}; conditioning on the exponentials in
\PR{Lemmas}{appx-lem:envelope}{B.1} and~\PN{appx-lem:endpoint-free}{B.2}
and on the candidate pairs in \PR{Appendix}{appx-app:concentration}{D};
interchanges of sums and expectations at the end of the proof of
\PR{Lemma}{appx-an:lem:restoration}{C.2}; Step~3 of the proof of
\PR{Lemma}{appx-an:lem:surface}{C.5}; iterated integrals over a hyperplane
and its normal in Subsections~\ref{tb:sec:surface}--\ref{tb:sec:coarea}.}
\emph{Source.} \TBc{Fol99}, Section~2.5.

\begin{fact}[Differentiation under the integral sign; Taylor's formula]
\label{tb:taylor}
(a) Let \(I\) be an open interval and let \(f(\lambda,\omega)\) be
integrable in \(\omega\) and differentiable in \(\lambda\in I\), with
\(|\partial_\lambda f(\lambda,\omega)|\le g(\omega)\) for an integrable
function \(g\). Then \(\frac d{d\lambda}\E f(\lambda,\cdot)
=\E\,\partial_\lambda f(\lambda,\cdot)\) on \(I\).

(b) For \(f\in C^q(\R^n)\) and \(x,v\in\R^n\),
\[
 f(x+v)=\sum_{k=0}^{q-1}\frac1{k!}D^kf(x)[v^{\otimes k}]
 +\int_0^1\frac{(1-s)^{q-1}}{(q-1)!}D^qf(x+sv)[v^{\otimes q}]\,ds.
\]
For \(q=1\) this is \(f(x+v)-f(x)=\int_0^1Df(x+sv)[v]\,ds\).
\end{fact}

\begin{proof}
For (a), the mean value theorem bounds the difference quotients by \(g\),
and they converge pointwise, so dominated convergence applies. For (b),
the function \(s\mapsto f(x+sv)\) is of class \(C^q\) with \(k\)th derivative
\(D^kf(x+sv)[v^{\otimes k}]\); apply the one-dimensional Taylor formula with
integral remainder at \(s=0\) and evaluate it at \(s=1\).
\end{proof}

\usedin{\PR{Lemma}{appx-an:lem:identity}{C.1}, where the derivative of
\(P_A(\lambda)=\E F(U_\lambda)\) is taken under the expectation (the
summands are finitely many and bounded) and Steps~1 and~3 use (b), to degree
three with the fourth-order integral remainder, to degree two for the
derivative, and to first order for the difference of third derivatives;
\PR{Lemma}{appx-lem:endpoint-gap}{B.3}, where (b) with \(q=1\) writes
\(\E[W\cdot\nabla\Psi(W)]\) as an integral of the Hessian along the segment
from \(W_i\) to \(W\).}
\emph{Source.} \TBc{Fol99}, Section~2.3.

\begin{fact}[Fa\`a di Bruno's formula]\label{tb:faa}
Let \(\phi\) be a smooth function on an open interval of \(\R\) and let
\(g:\R^n\to\R\) be smooth with values in that interval. For \(q\ge1\), let
\(\Pi_q\) be the set of partitions \(\mathcal P\) of \(\{1,\ldots,q\}\) into
nonempty blocks, let \(|\mathcal P|\) be the number of blocks, and for a
block \(J\) let \(z_J\) be the list of directions \(z_a\) with
\(a\in J\). Then
\[
 D^q(\phi\circ g)(x)[z_1,\ldots,z_q]
 =\sum_{\mathcal P\in\Pi_q}\phi^{(|\mathcal P|)}(g(x))
  \prod_{J\in\mathcal P}D^{|J|}g(x)[z_J].
\]
\end{fact}

\begin{proof}
For \(q=1\) this is the chain rule. Differentiate the right side for \(q\)
in the direction \(z_{q+1}\). By the product rule, the derivative either
falls on \(\phi^{(|\mathcal P|)}(g(x))\), which gives
\(\phi^{(|\mathcal P|+1)}(g(x))\,Dg(x)[z_{q+1}]\) and adds the block
\(\{q+1\}\) to \(\mathcal P\), or falls on one factor
\(D^{|J|}g(x)[z_J]\), which adds \(q+1\) to the block \(J\). Every partition
of \(\{1,\ldots,q+1\}\) arises exactly once in this way, by removing
\(q+1\) from its block.
\end{proof}

\begin{fact}[Derivatives of the logistic approximation]\label{tb:logistic}
Let \(\nu_1,\ldots,\nu_m\in\R^n\), \(a_1,\ldots,a_m\in\R\), \(h>0\), and
\[
 w_\alpha(x)=e^{(\nu_\alpha\cdot x-a_\alpha)/h},\qquad
 T=\sum_\alpha w_\alpha,\qquad
 F=\frac1{1+T},\qquad
 \pi_\alpha=\frac{w_\alpha}{1+T}=Fw_\alpha .
\]
For directions \(z_1,\ldots,z_q\), with the blocks \(J_1,\ldots,J_{|\mathcal P|}\)
of \(\mathcal P\) listed in any fixed order,
\[
 D^qF(x)[z_1,\ldots,z_q]
 =h^{-q}\sum_{\mathcal P\in\Pi_q}(-1)^{|\mathcal P|}|\mathcal P|!
 \sum_{\alpha_1,\ldots,\alpha_{|\mathcal P|}}
 F(x)\prod_{b=1}^{|\mathcal P|}\pi_{\alpha_b}(x)
 \prod_{b=1}^{|\mathcal P|}\prod_{a\in J_b}(\nu_{\alpha_b}\cdot z_a),
\]
and \(\sum_{\mathcal P\in\Pi_q}|\mathcal P|!\le q!\,2^{q-1}\). Consequently
\(|D^qF(x)[z_1,\ldots,z_q]|\le q!\,2^{q-1}h^{-q}\,F(x)\prod_{a=1}^q
\max_\alpha|\nu_\alpha\cdot z_a|\).
\end{fact}

\begin{proof}
Apply Fact~\ref{tb:faa} with \(\phi(\tau)=(1+\tau)^{-1}\) and \(g=T\). Then
\(\phi^{(k)}(\tau)=(-1)^kk!\,(1+\tau)^{-k-1}\), and since each \(w_\alpha\) is
an exponential of an affine function,
\(D^lT(x)[z_1,\ldots,z_l]=h^{-l}\sum_\alpha w_\alpha(x)\prod_{a\le l}
(\nu_\alpha\cdot z_a)\). Substituting, and using
\((1+T)^{-|\mathcal P|-1}\prod_bw_{\alpha_b}=F\prod_b\pi_{\alpha_b}\) and
\(\sum_b|J_b|=q\), gives the formula. The number
\(\sum_{\mathcal P}|\mathcal P|!\) counts ordered partitions
\((J_1,\ldots,J_k)\) of \(\{1,\ldots,q\}\). Writing the elements of
\(J_1\), then of \(J_2\), and so on, in increasing order inside each block,
gives a permutation together with the set of positions at which a new block
starts, which is a subset of the \(q-1\) gaps between consecutive entries.
Every ordered partition arises from at least one of these
\(q!\,2^{q-1}\) pairs. For the bound, the sum over \(\alpha_b\) of
\(\pi_{\alpha_b}\) is \(1-F\le1\), and each product of projections is
at most \(\prod_a\max_\alpha|\nu_\alpha\cdot z_a|\).
\end{proof}

\usedin{Subsection~\PN{appx-sec:mixture}{B.2}, where the first display
and the bound are \PE{appx-eq:partition-derivative}{B.1} and
\PE{appx-eq:all-orders}{B.2} (the text there cites the softmax cumulant formula,
\TBc{GZ}, Lemma~13); the envelope~\PE{appx-eq:local-envelope}{B.3} is
obtained by summing the same expansion over coordinate tuples;
\PR{Lemma}{appx-an:lem:restoration}{C.2}, Step~1, with \(u_\alpha\) in place
of \(\nu_\alpha\).}
\emph{Source.} \TBc{Com74}; the proof above is complete.

\subsection{Gaussian identities}\label{tb:sec:gauss}

\begin{fact}[Gaussian density in orthogonal coordinates]\label{tb:gauss-coord}
(a) Let \(\R^n=L\oplus L^\perp\) be an orthogonal decomposition. Then
\(\varphi_n(l+w)=\varphi_L(l)\varphi_{L^\perp}(w)\) for \(l\in L\) and
\(w\in L^\perp\), so the projections of a standard Gaussian vector onto
\(L\) and \(L^\perp\) are independent standard Gaussian vectors. In
particular, for a unit vector \(u\), \(a\in\R\), and \(y\perp u\),
\(\varphi_n(au+y)=\varphi_1(a)\varphi_{n-1}(y)\).

(b) For \(x,v\in\R^n\),
\(\varphi_n(x+v)=\varphi_n(x)\exp\{-x\cdot v-\tfrac12\norm v^2\}\).
\end{fact}

\begin{proof}
Both follow from \(\norm{l+w}^2=\norm l^2+\norm w^2\) and from expanding
\(\norm{x+v}^2\).
\end{proof}

\usedin{\PR{Lemma}{appx-an:lem:mean}{C.3}, where the normal and tangential
components of \(G\) relative to a face are independent;
\PR{Lemma}{appx-an:lem:surface}{C.5}, Steps~1 and~2, for the factorization of
\(\varphi_r\) on a facet and on a ridge and for the density ratio on a tube.}

\begin{fact}[Gaussian integration by parts]\label{tb:stein}
Let \(G\sim N(0,I_n)\) and let \(f\in C^2(\R^n)\) have \(f\) and its partial
derivatives of orders at most two of polynomial growth. Then
\begin{enumerate}
\item[(a)] \(\E[Gf(G)]=\E\nabla f(G)\);
\item[(b)] \(\E[(GG^\top-I_n)f(G)]=\E D^2f(G)\); equivalently, for \(z\in\R^n\),
\(\E D^2f(G)[z,z]=\E\bigl[\tr\{(GG^\top-I_n)zz^\top\}f(G)\bigr]\), where
\(\tr\{(GG^\top-I_n)zz^\top\}=(z\cdot G)^2-\norm z^2\);
\item[(c)] if \(v\in\R^n\), \(\sigma>0\), \(U=v+\sigma G\), and
\(F\in C^{s+2}(\R^n)\) has derivatives of polynomial growth, then for
\(z,z_1,\ldots,z_s\in\R^n\),
\[
 \E D^{s+2}F(U)[z,z,z_1,\ldots,z_s]
 =\sigma^{-2}\E\bigl[\tr\{(GG^\top-I_n)zz^\top\}\,
 D^sF(U)[z_1,\ldots,z_s]\bigr].
\]
\end{enumerate}
\end{fact}

\begin{proof}
For (a) in one coordinate, \(x\varphi_1(x)=-\varphi_1'(x)\), so integration by
parts gives
\[\int xf(x)\varphi_1(x)\,dx=\int f'(x)\varphi_1(x)\,dx,\]
and the boundary terms vanish because of the growth condition. The other coordinates
are handled by Fact~\ref{tb:fubini}. For (b), apply (a) to \(\partial_if\) and
to \(G_kf\): \(\E[G_iG_kf(G)]=\E\partial_i(G_kf)(G)=\delta_{ik}\E f(G)
+\E[G_k\partial_if(G)]=\delta_{ik}\E f(G)+\E\partial_{ik}f(G)\). For (c), the
function \(f(g)=D^sF(v+\sigma g)[z_1,\ldots,z_s]\) has
\(D^2f(g)[z,z]=\sigma^2D^{s+2}F(v+\sigma g)[z,z,z_1,\ldots,z_s]\); apply (b).
\end{proof}

\usedin{Subsection~\PN{sec:outline-late}{2.4} (the display after
``integration by parts against the covariance of the summands''),
\PR{Lemma}{appx-an:lem:restoration}{C.2}, Step~1, where (c) is applied
conditionally on \(\mathcal H\) with \(v=V_0+v_i\) and \(z=Z_i\);
\PR{Lemma}{appx-lem:endpoint-gap}{B.3}, for \(\E[Z\cdot\nabla\Psi(Z)]
=\E\tr(\Sigma D^2\Psi(Z))\), which is (a) with \(Z=\Gamma G\) for a factor
\(\Gamma\) of the covariance.}
\emph{Source.} Stein~\TBc{Ste81}; the proof above is complete.

\begin{fact}[Heat equation for the Gaussian component]\label{tb:heat}
Let \(Z\sim N(0,\Sigma)\) in \(\R^p\), where \(\Sigma\) may be singular, and
let \(f\in C^2(\R^p)\) be bounded with bounded first and second derivatives.
Then, for \(x\in\R^p\) and \(\lambda>0\),
\[
 \frac d{d\lambda}\E f(x+\sqrt\lambda Z)
 =\tfrac12\E\tr\{\Sigma D^2f(x+\sqrt\lambda Z)\}.
\]
\end{fact}

\begin{proof}
Write \(Z=\Gamma G\) with \(\Gamma\Gamma^\top=\Sigma\) and \(G\) standard
Gaussian. By Fact~\ref{tb:taylor}(a), which applies because the difference
quotients are bounded by \(\norm\nabla f_\infty\norm{\Gamma G}/(2\sqrt\lambda)\)
locally uniformly in \(\lambda>0\), the derivative equals
\((2\sqrt\lambda)^{-1}\E[G\cdot\Gamma^\top\nabla f(x+\sqrt\lambda\Gamma G)]\).
Fact~\ref{tb:stein}(a) applied to the function \(g\mapsto
\Gamma^\top\nabla f(x+\sqrt\lambda\Gamma g)\) turns this into
\((2\sqrt\lambda)^{-1}\sqrt\lambda\,\E\tr\{\Gamma^\top D^2f\,\Gamma\}
=\frac12\E\tr\{\Sigma D^2f\}\).
\end{proof}

\usedin{Subsection~\PN{sec:outline-path}{2.2} and the proof of
\PR{Lemma}{appx-an:lem:identity}{C.1}, Step~2, where the Gaussian part of the
path contributes \(\frac12\sum_iM_i:\E D^2F(U_\lambda)\).}

\subsection{Matrix and entropy inequalities}\label{tb:sec:matrix}

\begin{fact}[Duality between the operator and nuclear norms]\label{tb:nuclear}
(a) For every square matrix \(\Xi\) and every symmetric matrix \(M\),
\(|\tr(\Xi M)|\le\norm\Xi_\op\norm M_*\). Moreover
\(\norm M_*=\max\{\tr(\Xi M):\norm\Xi_\op\le1\}\), so that \(\norm\cdot_*\) is
a norm on symmetric matrices.

(b) If \(M=P-N\) with \(P\succeq0\) and \(N\succeq0\), then
\(\norm M_*\le\tr P+\tr N=2\tr P-\tr M\).
\end{fact}

\begin{proof}
For (a), let \(M=\sum_i\mu_ie_ie_i^\top\) be a spectral decomposition. Then
\(\tr(\Xi M)=\sum_i\mu_i\,e_i^\top\Xi e_i\) and \(|e_i^\top\Xi e_i|\le
\norm\Xi_\op\), which gives the inequality. Equality is attained at
\(\Xi=\sum_i\operatorname{sgn}(\mu_i)e_ie_i^\top\). A supremum of linear
functionals satisfies the triangle inequality. For (b), the norm of a
positive semidefinite matrix is its trace, so the triangle inequality gives
\(\norm M_*\le\tr P+\tr N\), and \(\tr N=\tr P-\tr M\).
\end{proof}

\usedin{Subsection~\PN{sec:outline-late}{2.4}; \PR{Lemma}{appx-an:lem:restoration}{C.2},
Step~4; \PR{Lemma}{appx-an:lem:matrix}{C.4}, Step~1, where
(b) is applied with \(P=p_\alpha\mu_\alpha\mu_\alpha^\top\) and
\(N=p_\alpha(I_r-\Sigma_\alpha)\), and gives
\PE{appx-an:matrix-trace-reduction}{C.17}.}

\begin{fact}[Log-sum inequality and a concave function]\label{tb:logsum}
(a) For positive numbers \(y_1,\ldots,y_l\) and \(z_1,\ldots,z_l\) with sums
\(y\) and \(z\),
\(\sum_jy_j\log(z_j/y_j)\le y\log(z/y)\).

(b) The function \((x,y)\mapsto x\log(y/x)\) is concave on \((0,\infty)^2\).
Hence, for positive integrable random variables \(X,Y\),
\(\E[X\log(Y/X)]\le\E X\,\log(\E Y/\E X)\).
\end{fact}

\begin{proof}
For (a), Jensen's inequality for the concave function \(\log\), with the
weights \(y_j/y\), gives \(\sum_j(y_j/y)\log(z_j/y_j)\le\log\sum_j(y_j/y)(z_j/y_j)
=\log(z/y)\). For (b), the function \(f(x,y)=x\log y-x\log x\) has Hessian
\(\begin{pmatrix}-1/x&1/y\\1/y&-x/y^2\end{pmatrix}\), whose determinant is
zero and whose trace is negative, so the Hessian is negative semidefinite;
the second statement is Jensen's inequality for the concave function
\(f\) and the random vector \((X,Y)\).
\end{proof}

\usedin{\PR{Lemma}{appx-an:lem:matrix}{C.4}, Step~3: (a) adds the reference
masses of the tuples that share a first face, and (b) averages over
\(V_0\) and over the faces.}
\emph{Source.} \TBc{CT06}, Theorem~2.7.1, for (a).

\subsection{Log-concavity}\label{tb:sec:logconcave}

A function \(f:\R^n\to[0,\infty)\) is \emph{log-concave} if
\(f((1-t)x+ty)\ge f(x)^{1-t}f(y)^t\) for all \(x,y\) and \(t\in[0,1]\);
when \(f>0\) this says that \(\log f\) is concave.

\begin{fact}[Pr\'ekopa--Leindler inequality and Pr\'ekopa's theorem]
\label{tb:prekopa}
(a) Let \(0<t<1\) and let \(f,g,h:\R^n\to[0,\infty)\) be bounded,
continuous, and integrable, with
\(h((1-t)x+ty)\ge f(x)^{1-t}g(y)^t\) for all \(x,y\). Then
\(\int h\ge(\int f)^{1-t}(\int g)^t\).

(b) Let \(f:\R^a\times\R^b\to[0,\infty)\) be bounded, continuous, and
log-concave, with \(f(z,\cdot)\) integrable for every \(z\). Then
\(z\mapsto\int f(z,y)\,dy\) is log-concave on \(\R^a\).
\end{fact}

\begin{proof}
(a) For \(n=1\), assume first that \(\sup f=\sup g=1\); the general case
follows by dividing \(f,g,h\) by \(\sup f\), \(\sup g\), and
\((\sup f)^{1-t}(\sup g)^t\). For \(0<s<1\) the level sets
\(\{f\ge s\}\) and \(\{g\ge s\}\) are nonempty, and the hypothesis gives
\(\{h\ge s\}\supseteq(1-t)\{f\ge s\}+t\{g\ge s\}\). In one dimension
\(|A+B|\ge|A|+|B|\) for nonempty compact sets \(A\) and \(B\), because
\(A+B\) contains the translates \(A+\min B\) and \(\max A+B\), which meet in
one point; the inequality passes to measurable sets by approximation from
inside by compact sets. Therefore
\(|\{h\ge s\}|\ge(1-t)|\{f\ge s\}|+t|\{g\ge s\}|\), and integration over
\(s\in(0,1)\) gives \(\int h\ge(1-t)\int f+t\int g\ge(\int f)^{1-t}(\int g)^t\),
by the inequality between the arithmetic and geometric means. For \(n\ge2\),
apply the case of dimension \(n-1\) to the sections of \(f,g,h\) at the last
coordinates \(s_0\), \(s_1\), and \((1-t)s_0+ts_1\), and then the case of
dimension one to the integrals of the sections, with Fubini's theorem
(Fact~\ref{tb:fubini}) connecting the steps; the routine measurability
checks are omitted here and can be found in~\TBc{Gar02}.

(b) Fix \(z_0,z_1\) and put \(z_t=(1-t)z_0+tz_1\). Log-concavity of \(f\)
gives \(f(z_t,(1-t)y_0+ty_1)\ge f(z_0,y_0)^{1-t}f(z_1,y_1)^t\), so (a) applies
to \(f(z_0,\cdot)\), \(f(z_1,\cdot)\), and \(f(z_t,\cdot)\).
\end{proof}

\emph{Source.} Pr\'ekopa~\TBc{Pre71,Pre73}, Leindler~\TBc{Lei72}; a
survey with the proof of the general case is~\TBc{Gar02}, Section~7. The
paper cites Theorem~6 of~\TBc{Pre73}.

\begin{fact}[Gaussian laws tilted by a log-concave weight]\label{tb:tilt}
Let \(w:\R^r\to[0,\infty)\) be bounded, continuous, and log-concave, with
\(p=\E w(G)>0\) for \(G\sim N(0,I_r)\). Let \(\mu\) and \(K\) be the mean and
covariance matrix of the probability law with density \(w(g)\varphi_r(g)/p\).
Then \(K\preceq I_r\), and
\(\E[(GG^\top-I_r)w(G)]=p(K-I_r+\mu\mu^\top)\).
\end{fact}

\begin{proof}
The function \((z,g)\mapsto e^{-\norm{g-z}^2/2}w(g)\) is log-concave, since
\(\norm{g-z}^2\) is jointly convex in \((z,g)\), and it is bounded and
continuous. By Fact~\ref{tb:prekopa}(b), \(\ell(z)=\log\int e^{-\norm{g-z}^2/2}
w(g)\,dg\) is concave. Since \(\norm{g-z}^2=\norm g^2-2z\cdot g+\norm z^2\),
\[
 \ell(z)=\text{const}-\tfrac12\norm z^2+\log M(z),\qquad
 M(z)=\int e^{z\cdot g}w(g)\varphi_r(g)\,dg=p\,\E_\nu e^{z\cdot g},
\]
where \(\nu\) is the tilted law. The function \(M\) is finite and smooth,
because \(w\) is bounded, and the Hessian of \(\log M\) at \(0\) is the
covariance \(K\) of \(\nu\). A concave function of class \(C^2\) has a
negative semidefinite Hessian, so \(K-I_r\preceq0\). The identity is the
decomposition of the second moment \(\E_\nu gg^\top=K+\mu\mu^\top\).
\end{proof}

\usedin{Subsection~\PN{sec:outline-late}{2.4} (``Pr\'ekopa's theorem \dots\
has covariance \(K\preceq I_r\)'') and \PR{Lemma}{appx-an:lem:matrix}{C.4},
Step~1, where \(w=w_\alpha\) and \(K=\Sigma_\alpha\).}
\emph{Source.} For indicators of convex cells, the cell covariance bound
\TBc{GZ}, Lemma~11, is the same estimate \(K\preceq I_r\), obtained from the
Brascamp--Lieb inequality~\TBc{BL76}; the weights here are not indicators,
which is why the proof goes through Pr\'ekopa's theorem.

\begin{fact}[Log-concavity of the logistic weights]\label{tb:lc-weights}
With the notation of Fact~\ref{tb:logistic}, for every \(k\ge0\) and every
tuple \((\alpha_1,\ldots,\alpha_k)\), the function
\(F\pi_{\alpha_1}\cdots\pi_{\alpha_k}=\exp\{\sum_{a\le k}t_{\alpha_a}\}
(1+T)^{-k-1}\), where \(t_\alpha=(\nu_\alpha\cdot x-a_\alpha)/h\), is
log-concave on \(\R^n\), and so is its composition
\(g\mapsto F\pi_{\alpha_1}\cdots\pi_{\alpha_k}(v+\sigma g)\) with any affine map.
\end{fact}

\begin{proof}
The logarithm is a sum of the affine functions \(t_{\alpha_a}\) and
\(-(k+1)\log(1+\sum_\alpha e^{t_\alpha})\). The function
\(\log(1+\sum_\alpha e^{t_\alpha})\) is the logarithm of a sum of
exponentials of the affine functions \(0,t_1,\ldots,t_m\) and is convex by
H\"older's inequality: \(\sum_\alpha e^{(1-s)t_\alpha+st'_\alpha}\le
(\sum_\alpha e^{t_\alpha})^{1-s}(\sum_\alpha e^{t'_\alpha})^s\). Composition
with an affine map preserves concavity.
\end{proof}

\usedin{Subsection~\PN{sec:outline-late}{2.4} and Step~1 of the proof of
\PR{Lemma}{appx-an:lem:matrix}{C.4}.}

\subsection{Polyhedra: faces, normal cones, and nearest points}\label{tb:sec:poly}

A \emph{polyhedron} is a set
\[
 Q=\bigcap_{j=1}^m\{x\in\R^n:u_j\cdot x\le a_j\},
\]
with unit vectors \(u_j\), called \emph{outward normals}, and numbers
\(a_j\), called \emph{offsets}. It is closed and convex. It is
\emph{full-dimensional} if it has an interior point, and \emph{proper} if it
is neither empty nor \(\R^n\). The description is \emph{irredundant} if
deleting any one inequality changes the set; deleting redundant or repeated
inequalities one at a time turns every description into an irredundant one
of the same set. For \(j\le m\), put \(H_j=\{x:u_j\cdot x=a_j\}\) and
\(F_j=Q\cap H_j\). A \emph{face} of \(Q\) is \(Q\) itself or a set
\(Q\cap\{u\cdot x=c\}\) with \(u\ne0\) and \(Q\subseteq\{u\cdot x\le c\}\).
A \emph{facet} is a face of dimension \(n-1\), and a \emph{ridge} is a face
of dimension \(n-2\). In the plane the facets are the sides of a polygon
and the ridges its corners, in \(\R^3\) the facets are the faces of a solid
and the ridges its edges, and on the line the facets are the endpoints of an
interval, which has no ridges. We use these standard properties of faces:
the relative interiors of the faces of \(Q\) are pairwise disjoint and cover
\(Q\); a face of a face of \(Q\) is a face of \(Q\); and a face \(G'\)
of a face \(G\) with \(G'\ne G\) lies in the relative boundary
\(G\setminus\relint G\) of \(G\) (\TBc{Roc70}, Section~18; \TBc{Zie95}, Chapter~2).

\begin{fact}[Facets and ridges of a polyhedron]\label{tb:facets}
Let \(Q\) be a proper, full-dimensional polyhedron in \(\R^n\) with an
irredundant description.
\begin{enumerate}
\item[(a)] The normals \(u_1,\ldots,u_m\) are pairwise distinct, each
\(F_j\) is a facet, \(Q\) has no other facets, and
\(\partial Q=F_1\cup\cdots\cup F_m\).
\item[(b)] For \(j\ne k\), \(F_j\cap F_k\subseteq H_j\cap H_k\), and
\(H_j\cap H_k\) is empty if \(u_k=-u_j\), and an affine subspace of
dimension \(n-2\) otherwise. In particular \(F_j\cap F_k\) lies in a
finite union of affine subspaces of dimension at most \(n-2\).
\item[(c)] Each \(F_j\) is a full-dimensional polyhedron in the hyperplane
\(H_j\). Write the points of \(H_j\) as \(a_ju_j+y\) with \(y\perp u_j\), and,
for \(k\ne j\) with \(u_k\ne\pm u_j\), put
\(\alpha_{jk}=\arccos(u_j\cdot u_k)\in(0,\pi)\),
\[
 \nu_{k\mid j}=\frac{u_k-(u_j\cdot u_k)u_j}{\sin\alpha_{jk}},\qquad
 d_{k\mid j}=\frac{a_k-(u_j\cdot u_k)a_j}{\sin\alpha_{jk}} .
\]
Then \(F_j\) is described in the coordinate \(y\) by the inequalities
\(\nu_{k\mid j}\cdot y\le d_{k\mid j}\), and its facets are exactly the
ridges \(E_{jk}=F_j\cap F_k\) of \(Q\) of dimension \(n-2\); the outward
unit normal of \(F_j\) at \(E_{jk}\) is \(\nu_{k\mid j}\).
\item[(d)] Every ridge of \(Q\) lies in exactly two facets, so the ridges
correspond to the unordered pairs \(\{j,k\}\) with
\(\dim(F_j\cap F_k)=n-2\); for such a pair \(u_k\ne\pm u_j\), and the
only inequalities active at a point of \(\relint E_{jk}\) are \(j\) and \(k\).
\end{enumerate}
\end{fact}

\begin{proof}[Proof (sketch)]
(a) Fix \(j\). Irredundancy gives a point \(x_j\) with \(u_j\cdot x_j>a_j\)
and \(u_k\cdot x_j\le a_k\) for \(k\ne j\). Let \(x_0\) be an interior point of
\(Q\), and let \(x_*\) be the point where the segment from \(x_0\) to \(x_j\)
crosses \(H_j\). For \(k\ne j\), \(u_k\cdot x_*<a_k\), because \(x_*\) is a
convex combination of \(x_0\), where the inequality is strict, with positive
weight, and of \(x_j\), where it holds weakly. Hence a neighborhood of
\(x_*\) in \(H_j\) lies in \(F_j\), and \(\dim F_j=n-1\). Two inequalities
with the same normal cannot both be irredundant, since the one with the
larger offset is implied by the other. A point of \(Q\) at
which no inequality is active is interior, and at a point where
inequality \(j\) is active the points \(x+\varepsilon u_j\) lie outside \(Q\),
so \(\partial Q=\bigcup_jF_j\). If \(G=Q\cap\{u\cdot x=c\}\) is a facet and
\(x\in\relint G\), some inequality \(j\) is active at \(x\); on \(G\) the
affine function \(u_j\cdot x-a_j\le0\) attains its maximum at a relative
interior point and is therefore zero, so \(G\subseteq H_j\), and since \(G\) has
dimension \(n-1\), \(G=Q\cap H_j=F_j\).

(b) If \(u_k=-u_j\) and \(H_j\cap H_k\ne\emptyset\), then \(Q\) would lie in
the hyperplane \(H_j\), contradicting full dimension. If \(u_k\ne\pm u_j\),
the two equations are independent.

(c) The set \(F_j\) is the polyhedron in \(H_j\) cut out by the remaining
inequalities, which take the stated form in the coordinate \(y\) because
\(u_k\cdot(a_ju_j+y)=a_ju_j\cdot u_k+(u_k-(u_j\cdot u_k)u_j)\cdot y\); an
inequality with \(u_k=-u_j\) holds automatically on \(H_j\) since
\(F_j\ne\emptyset\). By (a) applied inside \(H_j\), the facets of \(F_j\) are
the faces of \(F_j\) of dimension \(n-2\), hence faces of \(Q\) of that
dimension, that is, ridges of \(Q\) in \(F_j\); each equals
\(F_j\cap H_k=F_j\cap F_k\) for the inequality \(k\) that defines it.

(d) Let \(E\) be a ridge and \(y\in\relint E\), and let \(J\) be the set of
inequalities active at \(y\). On \(E\) each affine function
\(u_i\cdot x-a_i\), \(i\in J\), is at most zero and vanishes at the relative
interior point \(y\), so it vanishes on \(E\); thus \(E\subseteq F_i\) for
\(i\in J\) and \(\dim\operatorname{span}\{u_i:i\in J\}\le2\). A single active
inequality is impossible, since then \(y\) would lie in the relative interior
of \(F_j\) and the only face of \(F_j\) containing it is \(F_j\) itself, of
dimension \(n-1\). Opposite normals \(u_i=-u_j\) are excluded by (b). Hence
the active normals span a plane \(\Pi\). Near \(y\), \(Q\) coincides with
\(y\) plus the product of \(\Pi^\perp\) and a planar cone \(K\subseteq\Pi\)
with nonempty interior, defined by the active inequalities and contained in
a wedge of opening \(\pi-\alpha_{jk}<\pi\); this cone has exactly two
boundary rays. Each active inequality \(i\) has \(\dim F_i=n-1\), so
\(F_i\) meets every neighborhood of \(y\) in a set of dimension \(n-1\), and
the line \(\{u_i\cdot v=0\}\) must therefore carry a boundary ray of \(K\).
With distinct normals this leaves exactly two active inequalities.
\end{proof}

\usedin{\PR{Lemma}{appx-an:lem:surface}{C.5}, in the paragraph on the terms
from the geometry of polyhedra and in Step~1 (``each inequality defines
exactly one facet, \dots\ two facets that share a ridge have normals that
are neither equal nor opposite'') and Step~2 (``every ridge lies in exactly
two facets'').}
\emph{Source.} \TBc{Zie95}, Chapter~2; \TBc{Roc70}, Sections~18 and~19;
the facet decomposition of \(\partial Q\) with the pairing of ridges is
the content of \TBc{GZ}, Lemma~27.

\begin{fact}[Normal cones]\label{tb:normalcone}
Let \(Q\) be a polyhedron as in Fact~\ref{tb:facets}. For \(x\in Q\) let
\(J(x)=\{j:u_j\cdot x=a_j\}\), and let
\[
 N_Q(x)=\{v\in\R^n:v\cdot(q-x)\le0\text{ for all }q\in Q\}
\]
be the \emph{normal cone} of \(Q\) at \(x\). Then
\(N_Q(x)=\{\sum_{j\in J(x)}c_ju_j:c_j\ge0\}\). In particular
\(N_Q(x)=\{0\}\) at interior points, \(N_Q(x)=\{cu_j:c\ge0\}\) for
\(x\in\relint F_j\), and \(N_Q(x)=\{c_ju_j+c_ku_k:c_j,c_k\ge0\}\), a sector of
angle \(\alpha_{jk}\) in the plane spanned by \(u_j\) and \(u_k\), for
\(x\in\relint E_{jk}\).
\end{fact}

\begin{proof}
Put \(T=\{d:u_j\cdot d\le0\text{ for }j\in J(x)\}\). A vector \(v\) lies in
\(N_Q(x)\) if and only if \(v\cdot d\le0\) for all \(d\in T\): if \(d\in T\)
then \(x+\varepsilon d\in Q\) for small \(\varepsilon>0\), because the
inactive inequalities are strict at \(x\), and conversely \(q-x\in T\) for
\(q\in Q\). By Farkas's lemma, which says that \(v\cdot d\le0\) for all
\(d\) with \(u_j\cdot d\le0\) (\(j\in J\)) if and only if \(v\) is a
nonnegative combination of the \(u_j\), \(j\in J\), this is the stated
cone. For \(x\in\relint F_j\), \(J(x)=\{j\}\), since an index \(k\ne j\) in
\(J(x)\) would put \(x\) in \(F_j\cap F_k\), a face of \(F_j\) of dimension
at most \(n-2\), and so in the relative boundary of \(F_j\). For
\(x\in\relint E_{jk}\), \(J(x)=\{j,k\}\) by Fact~\ref{tb:facets}(d).
\end{proof}

\usedin{\PR{Lemma}{appx-an:lem:surface}{C.5}, Step~2 (``the normal cone at a
relative-interior point of the ridge shared by facets \(j\) and \(k\) is the
cone spanned by \(u_j\) and \(u_k\)'') and the definition of the outward
tubes.}
\emph{Source.} \TBc{Roc70}, Sections~22 and~23, for Farkas's lemma and normal
cones; \TBc{RW98}, Chapter~6.

\begin{fact}[Nearest point]\label{tb:projection}
Let \(Q\subseteq\R^n\) be nonempty, closed, and convex. For each
\(y\in\R^n\) there is a unique point \(\prj_Q(y)\in Q\) with
\(\norm{y-\prj_Q(y)}=\dist(y,Q)\), and \(x=\prj_Q(y)\) if and only if
\(x\in Q\) and \(y-x\in N_Q(x)\). Consequently every \(y\in\R^n\) can be
written in exactly one way as \(y=x+v\) with \(x\in Q\) and \(v\in N_Q(x)\),
and then \(\norm v=\dist(y,Q)\). A direction \(v\) lies in \(N_Q(x)\) if and
only if \(x\) is the nearest point of \(x+tv\) in \(Q\) for every \(t>0\).
\end{fact}

\begin{proof}
For \(x,q\in Q\), \(\norm{y-q}^2=\norm{y-x}^2+\norm{x-q}^2-2(y-x)\cdot(q-x)\).
If \(y-x\in N_Q(x)\), the last term is nonnegative, so
\(\norm{y-q}^2\ge\norm{y-x}^2+\norm{x-q}^2>\norm{y-x}^2\) for \(q\ne x\); thus
\(x\) is the unique nearest point. Conversely, if \((y-x)\cdot(q-x)>0\) for
some \(q\in Q\), then moving from \(x\) toward \(q\) decreases the distance
to \(y\) for small steps, by the displayed identity applied to
\(x+\varepsilon(q-x)\in Q\), so \(x\) is not nearest. A nearest point exists
because \(Q\) is closed. The last statement follows because \(N_Q(x)\) is a
cone.
\end{proof}

\usedin{\PR{Lemma}{appx-an:lem:surface}{C.5}, Step~2, to show that the tubes
are disjoint (``their points have distinct metric projections'').}
\emph{Source.} \TBc{RW98}, Chapter~6 (Moreau's decomposition); \TBc{GZ},
proof of Lemma~28.

\subsection{Surface measure, Gaussian surface integrals, and the
divergence theorem}\label{tb:sec:surface}

\begin{fact}[Hausdorff measure on flat pieces]\label{tb:hausdorff}
For \(0\le k\le n\), the \(k\)-dimensional Hausdorff measure \(\Haus^k\) on
\(\R^n\) is the Borel measure
\[
 \Haus^k(S)=\lim_{\epsilon\downarrow0}\inf\Bigl\{\sum_i\omega_k
 \Bigl(\tfrac12\operatorname{diam}C_i\Bigr)^k:S\subseteq\bigcup_iC_i,\
 \operatorname{diam}C_i\le\epsilon\Bigr\},
\]
where \(\omega_k\) is the volume of the unit ball of \(\R^k\). (a) On an
affine subspace of dimension \(k\), \(\Haus^k\) is \(k\)-dimensional Lebesgue
measure; in particular \(\Haus^0\) is counting measure and \(\Haus^n\) is
Lebesgue measure on \(\R^n\). (b) \(\Haus^k\) is invariant under isometries.
(c) A set contained in a finite union of affine subspaces of dimension less
than \(k\) has \(\Haus^k\)-measure zero.
\end{fact}

Part (c) follows from (a), since each such subspace lies in a
\(k\)-dimensional affine subspace in which it is a Lebesgue null set. For a
proper, full-dimensional polyhedron \(Q\) with an irredundant description,
and for a function \(f\ge0\) or \(f\) integrable on \(\partial Q\), the
\emph{surface integral} is
\[
 \int_{\partial Q}f\,d\Haus^{n-1}=\sum_{j=1}^m\int_{F_j}f\,d\Haus^{n-1},
\]
where on each facet \(\Haus^{n-1}\) is Lebesgue measure on the hyperplane
\(H_j\) by part~(a). By Fact~\ref{tb:facets}(b) and part~(c), distinct
facets overlap in sets of \(\Haus^{n-1}\)-measure zero, so the sum is the
integral of \(f\) over the set \(\partial Q\) with respect to
\(\Haus^{n-1}\). It depends on the set \(Q\) and not on its description,
which justifies deleting redundant inequalities before computing it.
The Gaussian surface integral \(\surf(Q)=\int_{\partial Q}\Hh\,\varphi_r\,
d\Haus^{r-1}\) of Subsection~\PN{sec:outline-late}{2.4} is the case
\(n=r\) and \(f=\Hh\varphi_r\), with \(\Hh(g)=\norm g^2-r\).

\usedin{the definition of \(\surf\) in Subsection~\PN{sec:outline-late}{2.4}
and in \PR{Lemma}{appx-an:lem:surface}{C.5}.}
\emph{Source.} \TBc{EG15}, Chapter~2.

\begin{fact}[Gauss--Green formula (divergence theorem) for a polyhedron]\label{tb:divergence}
Let \(Q\subseteq\R^n\) (\(n\ge1\)) be a proper, full-dimensional polyhedron
with an irredundant description, facets \(F_1,\ldots,F_m\), and outward unit
normals \(u_1,\ldots,u_m\). Let \(X:\R^n\to\R^n\) be of class \(C^1\), with
\(\norm{X(x)}+|\operatorname{div}X(x)|\le C(1+\norm x)^k\varphi_n(x)\) for
some \(C\) and \(k\). Then
\[
 \int_Q\operatorname{div}X\,dx=\sum_{j=1}^m\int_{F_j}X\cdot u_j\,d\Haus^{n-1}.
\]
For \(n=1\) the facets are the endpoints of an interval, \(\Haus^0\) is
counting measure, and the identity is the fundamental theorem of calculus.
\end{fact}

\begin{proof}
It suffices to show, for a unit vector \(e\) and a function
\(f\in C^1(\R^n)\) with \(|f|+|\partial_ef|\le C(1+\norm x)^k\varphi_n(x)\),
that
\begin{equation}
 \int_Q\partial_ef\,dx=\sum_j(e\cdot u_j)\int_{F_j}f\,d\Haus^{n-1},
 \tag{\(*\)}
\end{equation}
since the theorem follows by taking \(e=e_i\) and \(f=X_i\) and summing over
\(i\). Write \(x=w+te\) with \(w\in e^\perp\). By Fubini's theorem,
\(\int_Q\partial_ef\,dx=\int_{e^\perp}\int_{I(w)}\partial_tf(w+te)\,dt\,dw\),
where \(I(w)=\{t:w+te\in Q\}\) is a closed interval, possibly empty, a point,
or unbounded, because \(Q\) is convex. Since \(f\) tends to zero along every
line, the inner integral equals \(f(w+t_+(w)e)-f(w+t_-(w)e)\), where
\(t_\pm(w)\) are the endpoints of \(I(w)\) and a term with an infinite
endpoint is read as zero.

Consider the upper endpoints \(x=w+t_+(w)e\). If \(x\in\relint F_j\), then
only the inequality \(j\) is active at \(x\) (Fact~\ref{tb:normalcone}), so
\(x+te\in Q\) for small \(t>0\) unless \(e\cdot u_j>0\), and \(x\) is an upper
endpoint if and only if \(e\cdot u_j>0\). The upper endpoints that lie in no
\(\relint F_j\) lie in \(\bigcup_{j\ne k}F_j\cap F_k\), which is a subset of a
finite union of affine subspaces of dimension at most \(n-2\)
(Fact~\ref{tb:facets}(b)); its projection along \(e\) onto \(e^\perp\) is a
Lebesgue null set of \(e^\perp\). For fixed \(j\) with \(e\cdot u_j>0\), the
map \(x\mapsto x-(x\cdot e)e\) is injective on the set of upper endpoints in
\(\relint F_j\), since a line meets \(Q\) in one interval, and it is the
restriction to \(H_j\) of a linear map whose Jacobian with respect to
\(\Haus^{n-1}\) on \(H_j\) is \(e\cdot u_j\): if \(b_1,\ldots,b_{n-1}\) is an
orthonormal basis of \(u_j^\perp\), the volume of the projection along \(e\)
of the cube they span equals the volume of the parallelepiped spanned by
\(b_1,\ldots,b_{n-1},e\), which is \(|e\cdot u_j|\). Therefore
\(\int f(w+t_+(w)e)\,dw=\sum_{j:\,e\cdot u_j>0}(e\cdot u_j)\int_{F_j}f\,d\Haus^{n-1}\),
the sets \(F_j\setminus\relint F_j\) being \(\Haus^{n-1}\)-null. The same
argument for the lower endpoints gives
\(-\sum_{j:\,e\cdot u_j<0}|e\cdot u_j|\int_{F_j}f\,d\Haus^{n-1}\), and
subtracting proves \((*)\); facets with \(e\cdot u_j=0\) contribute nothing
on either side.
\end{proof}

\usedin{\PR{Lemma}{appx-an:lem:surface}{C.5}, Step~1, applied in the hyperplane
\(H_j\) to the facet \(F_j\), regarded as a polyhedron of dimension
\(n=r-1\), whose facets are the ridges \(E_{jk}\) with outward normals
\(\nu_{k\mid j}\) (Fact~\ref{tb:facets}(c)), and to the vector field
\(X(y)=y\varphi_{r-1}(y)\), for which
\(\operatorname{div}X(y)=\{(r-1)-\norm y^2\}\varphi_{r-1}(y)\).
The decay of \(X\) is what replaces the truncation of unbounded facets.}
\emph{Source.} \TBc{GZ}, Lemma~24, proved there by truncation to a ball and
the Gauss--Green formula for Lipschitz domains~\TBc{EG15}, Chapters~4
and~5.

\subsection{Tubes around a polyhedron}\label{tb:sec:tubes}

Let \(Q\) be a proper, full-dimensional polyhedron in \(\R^n\) with an
irredundant description, facets \(F_j\), and ridges \(E_{jk}\) (\(n\ge2\)). For
a ridge, \(L_{jk}=\operatorname{span}(u_j,u_k)\), and \(z_{jk}\) is the unique
point of \(L_{jk}\) with \(u_j\cdot z_{jk}=a_j\) and \(u_k\cdot z_{jk}=a_k\). Every
point \(x\) of \(E_{jk}\) satisfies these two equations, so
\(E_{jk}\subseteq z_{jk}+L_{jk}^\perp\), and \(z_{jk}\) is the orthogonal
projection of every point of the ridge onto \(L_{jk}\). For \(\varepsilon>0\),
the \emph{outward tubes} of length \(\varepsilon\) are
\begin{align*}
 \mathcal T_j(\varepsilon)&=\{x+tu_j:x\in\relint F_j,\ 0<t<\varepsilon\},\\
 \mathcal T_{jk}(\varepsilon)&=\{x+v:x\in\relint E_{jk},\ v\in N_Q(x),\
 0<\norm v<\varepsilon\}.
\end{align*}
By Fact~\ref{tb:normalcone}, \(v\) ranges over the sector
\(S_{jk}(\varepsilon)=\{c_ju_j+c_ku_k:c_j,c_k\ge0\}\cap
\{\norm v<\varepsilon\}\) of angle \(\alpha_{jk}\) and area
\(\alpha_{jk}\varepsilon^2/2\) in \(L_{jk}\), without its apex. The tubes over the closed facets and ridges that
the proofs describe differ from these by sets of Lebesgue measure zero.
For \(Q(a)=\bigcap_j\{u_j\cdot x\le a_j\}\), write
\(Q(a)^t=Q(a_1+t,\ldots,a_m+t)\), the polyhedron with every offset raised by \(t\).

\begin{fact}[Tubes: disjointness and Gaussian mass]\label{tb:tubes}
(a) The sets \(\mathcal T_j(\varepsilon)\), \(j\le m\), and
\(\mathcal T_{jk}(\varepsilon)\), for the ridges \(E_{jk}\), are pairwise
disjoint, and they are contained in
\(\{y:0<\dist(y,Q)<\varepsilon\}\subseteq Q^\varepsilon\setminus Q\), where
\(Q^\varepsilon=Q(a)^\varepsilon\).

(b) With \(s_j=\int_{F_j}\varphi_n\,d\Haus^{n-1}\) and
\(\rho_{jk}=\int_{E_{jk}}\varphi_n\,d\Haus^{n-2}\),
\[
 \Phi_n(\mathcal T_j(\varepsilon))=\int_{F_j}\int_0^\varepsilon
 \varphi_n(x+tu_j)\,dt\,d\Haus^{n-1}(x),\qquad
 \Phi_n(\mathcal T_{jk}(\varepsilon))=\int_{E_{jk}}\int_{S_{jk}(\varepsilon)}
 \varphi_n(x+v)\,dv\,d\Haus^{n-2}(x).
\]

(c) Let \(R>0\) and \(\varepsilon=1/(1+R)\). If \(|a_j|\le R\), then
\(\Phi_n(\mathcal T_j(\varepsilon))\ge e^{-3/2}\,\varepsilon\,s_j\). If
\(\norm{z_{jk}}\le R\), then
\(\Phi_n(\mathcal T_{jk}(\varepsilon))\ge e^{-3/2}\,(\alpha_{jk}\varepsilon^2/2)\,
\rho_{jk}\).
\end{fact}

\begin{proof}
(a) By Fact~\ref{tb:projection}, a point \(y\notin Q\) is \(x+v\) for exactly
one pair with \(x=\prj_Q(y)\in Q\) and \(v\in N_Q(x)\setminus\{0\}\), and
\(\norm v=\dist(y,Q)\). A point of \(\mathcal T_j(\varepsilon)\) has its nearest
point in \(\relint F_j\), and a point of \(\mathcal T_{jk}(\varepsilon)\) has
its nearest point in \(\relint E_{jk}\), by Fact~\ref{tb:normalcone}. The
sets \(\relint F_j\) and \(\relint E_{jk}\) are pairwise disjoint, being
relative interiors of distinct faces, so the tubes are disjoint. If
\(\dist(y,Q)<\varepsilon\) and \(x=\prj_Q(y)\), then
\(u_j\cdot y\le u_j\cdot x+\norm{y-x}<a_j+\varepsilon\) for all \(j\), since the
\(u_j\) are unit vectors; thus \(y\in Q^\varepsilon\).

(b) The map \((x,t)\mapsto x+tu_j\) from \(H_j\times\R\) to \(\R^n\) is an
isometry, and \((x,v)\mapsto x+v\) is an isometry from
\((z_{jk}+L_{jk}^\perp)\times L_{jk}\) to \(\R^n\). Lebesgue measure on
\(\R^n\) is therefore the product of \(\Haus^{n-1}\) on \(H_j\) and \(dt\), and
of \(\Haus^{n-2}\) on \(z_{jk}+L_{jk}^\perp\) and \(dv\) on \(L_{jk}\)
(Fact~\ref{tb:hausdorff}(a)). Apply Fact~\ref{tb:fubini}.

(c) For \(x\in F_j\), \(x\cdot u_j=a_j\), and by Fact~\ref{tb:gauss-coord}(b)
\(\varphi_n(x+tu_j)=\varphi_n(x)\exp\{-ta_j-t^2/2\}\ge e^{-R\varepsilon-
\varepsilon^2/2}\varphi_n(x)\ge e^{-3/2}\varphi_n(x)\), since
\(R\varepsilon<1\) and \(\varepsilon<1\). Integrating over
\(t\in(0,\varepsilon)\) gives the first bound. For \(x\in E_{jk}\) and
\(v\in L_{jk}\), \(x\cdot v=z_{jk}\cdot v\), because \(x-z_{jk}\perp L_{jk}\),
so \(\varphi_n(x+v)=\varphi_n(x)\exp\{-z_{jk}\cdot v-\norm v^2/2\}
\ge e^{-3/2}\varphi_n(x)\) for \(\norm v<\varepsilon\). Integrating over the
sector, whose area is \(\alpha_{jk}\varepsilon^2/2\), gives the second bound.
\end{proof}

In the plane the tubes over the sides are strips and the tubes over the
corners are circular sectors, and together they tile the outer
\(\varepsilon\)-neighborhood of a polygon. When \(n\ge3\) the faces of
dimension at most \(n-3\) carry further strata of the neighborhood, which the
proofs do not need.

\usedin{\PR{Lemma}{appx-an:lem:surface}{C.5}, Step~2, for the bounds
\(q_j\ge c_*\varepsilon_Rs_j\) and \(q_{jk}\ge c_*\alpha_{jk}\varepsilon_R^2\rho_{jk}\),
and Step~3, for the disjointness that allows the Gaussian masses of the
tubes in a dyadic shell to be added, and for the inclusion
\(\mathcal O_\iota(V_0)\subseteq(Q')^{\varepsilon}\setminus Q'\) from which the
event \(V_0+\sigma G\in Q^{\sigma\varepsilon}\setminus Q\) is read: the
polyhedron \(Q'=(Q-V_0)/\sigma\) has offsets \((a_j-u_j\cdot V_0)/\sigma\),
and raising them by \(\varepsilon\) is raising those of \(Q\) by
\(\sigma\varepsilon\).}
\emph{Source.} \TBc{GZ}, Lemmas~28 and~29.

\subsection{The coarea formula and the first variation of a Gaussian
integral}\label{tb:sec:coarea}

\begin{fact}[Coarea formula]\label{tb:coarea}
Let \(\psi:\R^n\to\R\) be Lipschitz, so that \(\nabla\psi\) exists almost
everywhere. For every measurable \(f\ge0\),
\[
 \int_{\R^n}f\,\norm{\nabla\psi}\,dx=\int_{-\infty}^\infty
 \Bigl(\int_{\{\psi=s\}}f\,d\Haus^{n-1}\Bigr)ds ,
\]
and the same holds for integrable \(f\).
\end{fact}

\emph{Source.} \TBc{EG15}, Section~3.4. The proofs use the formula only for
a maximum of finitely many affine functions with gradients of equal length, and
for these Fact~\ref{tb:variation} gives a direct proof from Fubini's theorem.

\begin{fact}[First variation of a polyhedral sublevel integral]\label{tb:variation}
Let \(u_1,\ldots,u_m\) be unit vectors in \(\R^n\) (\(m\ge1\)), \(\kappa>0\),
\(b_1,\ldots,b_m\in\R\), and
\[
 \psi(g)=\max_{j\le m}\{\kappa\,u_j\cdot g-b_j\},\qquad Q_s=\{g:\psi(g)\le s\}.
\]
(a) For \(f\ge0\) measurable, or \(f\) integrable, and \(s\in\R\),
\[
 \int_{Q_s}f\,dg=\kappa^{-1}\int_{-\infty}^s\Bigl(\int_{\{\psi=s'\}}f\,d\Haus^{n-1}
 \Bigr)ds'.
\]
(b) If \(f\) is integrable, then \(\Lambda(s)=\int_{Q_s}f\,dg\) is absolutely
continuous, and for almost every \(s\)
\[
 \Lambda'(s)=\kappa^{-1}\int_{\{\psi=s\}}f\,d\Haus^{n-1},\qquad
 |\Lambda'(s)|\le\kappa^{-1}\sum_{j=1}^m\int_{\{\kappa u_j\cdot g-b_j=s\}}|f|\,d\Haus^{n-1}.
\]
(c) If \(s>\inf\psi\), then \(Q_s\) is a proper, full-dimensional polyhedron
with normals \(u_j\) and offsets \((s+b_j)/\kappa\), \(\{\psi=s\}=\partial Q_s\),
and \(\int_{\{\psi=s\}}f\,d\Haus^{n-1}\) is the surface integral of \(f\) over
\(\partial Q_s\) of Subsection~\ref{tb:sec:surface}. For \(s<\inf\psi\) the set
\(Q_s\) is empty, and \(s=\inf\psi\) is the only other level at which \(Q_s\)
may fail to be full-dimensional.

(d) If \(\Lambda\) is Lipschitz and \(S\) is a real random variable with a
density, then \(s\mapsto\E\Lambda(s+S)\) is differentiable with derivative
\(\E\Lambda'(s+S)\), where \(\Lambda'\) is defined almost everywhere.
\end{fact}

\begin{proof}
(a) Let \(C_j\) be the set where \(j\) is the smallest index attaining the
maximum defining \(\psi\); the Borel sets \(C_1,\ldots,C_m\) partition
\(\R^n\). On \(C_j\), \(\psi(g)=\kappa t-b_j\) in the coordinates
\(g=tu_j+y\), \(y\perp u_j\). By Fact~\ref{tb:fubini}, in these coordinates,
\[
 \int_{C_j\cap Q_s}f\,dg=\int_{u_j^\perp}\int_{-\infty}^{(s+b_j)/\kappa}
 (f\ind_{C_j})(tu_j+y)\,dt\,dy .
\]
The substitution \(s'=\kappa t-b_j\) turns the inner integral into
\(\kappa^{-1}\int_{-\infty}^s(f\ind_{C_j})(tu_j+y)\,ds'\) with
\(t=(s'+b_j)/\kappa\), and a further application of Fubini's theorem
exchanges the integrals. For fixed \(s'\), the set of \(tu_j+y\) in \(C_j\)
with this \(t\) is the slice \(C_j\cap\{\psi=s'\}\), and on the hyperplane
containing it \(\Haus^{n-1}\) is Lebesgue measure \(dy\)
(Fact~\ref{tb:hausdorff}(a)). Hence
\(\int_{C_j\cap Q_s}f\,dg=\kappa^{-1}\int_{-\infty}^s\int_{C_j\cap\{\psi=s'\}}
f\,d\Haus^{n-1}\,ds'\). Summing over \(j\) gives (a), since the slices
\(C_j\cap\{\psi=s'\}\) partition \(\{\psi=s'\}\).

(b) Applying (a) to \(|f|\) and letting \(s\to\infty\) shows that
\(I(s')=\kappa^{-1}\int_{\{\psi=s'\}}f\,d\Haus^{n-1}\) is integrable on
\(\R\), and (a) says \(\Lambda(s)=\int_{-\infty}^sI(s')\,ds'\). Hence \(\Lambda\) is
absolutely continuous, and \(\Lambda'=I\) almost everywhere by Lebesgue's
differentiation theorem. The bound follows because
\(\{\psi=s\}\subseteq\bigcup_j\{\kappa u_j\cdot g-b_j=s\}\), a union of \(m\)
hyperplanes.

(c) The set \(Q_s\) is the polyhedron \(\bigcap_j\{u_j\cdot g\le(s+b_j)/\kappa\}\),
and \(\psi\) is convex and continuous. If \(\psi(g)=s>\inf\psi\), choose \(g'\)
with \(\psi(g')<s\); for \(\tau>0\), \(g=(g_\tau+\tau g')/(1+\tau)\) with
\(g_\tau=g+\tau(g-g')\), so convexity gives
\(\psi(g_\tau)\ge s+\tau(s-\psi(g'))>s\). Thus \(g\) is a limit of points
outside \(Q_s\) and \(g\in\partial Q_s\); the open set \(\{\psi<s\}\) is
nonempty and lies in \(Q_s\), which is therefore full-dimensional, and it
is a proper subset because \(\psi(g)\ge\kappa u_1\cdot g-b_1\to\infty\) along
\(u_1\). A point with \(\psi<s\) is interior, so \(\partial Q_s=\{\psi=s\}\). After
the redundant and repeated inequalities are deleted, \(\partial Q_s=\{\psi=s\}\)
is the union of the sets \(Q_s\cap\{\kappa u_j\cdot g-b_j=s\}\) of the
remaining inequalities, which are faces of \(Q_s\) or empty. A face of
dimension \(n-1\) is a facet (Fact~\ref{tb:facets}(a)) and any other face is
\(\Haus^{n-1}\)-null, so the integral over \(\{\psi=s\}\) is the sum over the
facets of the irredundant description.

(d) A Lipschitz function is absolutely continuous and differentiable outside a
Lebesgue null set \(\mathcal N\). The random variable \(s+S\) avoids
\(\mathcal N\) almost surely, because \(S\) has a density, so the difference
quotients of \(\Lambda(s+S)\) converge almost surely and are bounded by the
Lipschitz constant. Dominated convergence gives the claim.
\end{proof}

\usedin{the paragraph ``First variation and common shifts'' after
\PR{Lemma}{appx-an:lem:surface}{C.5}, with
\(\psi(g)=\max_j\{u_j\cdot(V_0+\sigma g)-a_j\}\), \(\kappa=\sigma\),
\(b_j=a_j-u_j\cdot V_0\), \(Q_s=(Q(a+s\mathbf 1)-V_0)/\sigma\), and
\(f=\Hh\varphi_r\), which gives the derivative
\(\sigma^{-1}\surf((Q(a+s\mathbf 1)-V_0)/\sigma)\) for almost every \(s\) (parts
(a)--(c)); the Lipschitz bound in (b) is the bound \(Cm(r+1)/\sigma\) of that
paragraph, and part (d) justifies averaging the derivative over the
absolutely continuous common shift \(-h\log T_t\) in Step~4 of the proof of
\PR{Lemma}{appx-an:lem:matrix}{C.4}.}

\subsection{Where each fact is used}\label{tb:sec:index}

The table lists, for each fact of this appendix, the places in the proofs
that use it.

\medskip
\begin{center}
\small
\renewcommand{\arraystretch}{1.15}
\begin{tabular}{@{}p{0.27\textwidth}p{0.68\textwidth}@{}}
Facts~\ref{tb:fubini}, \ref{tb:taylor} &
every conditioning on an independent component in Appendix~\PN{appx-sec:support}{C};
the Taylor expansions in \PR{Lemma}{appx-an:lem:identity}{C.1}; Step~3 of
\PR{Lemma}{appx-an:lem:surface}{C.5}\\
Facts~\ref{tb:faa}, \ref{tb:logistic} &
the derivative formula and the local envelope in Subsection~\PN{appx-sec:mixture}{B.2};
Step~1 of \PR{Lemma}{appx-an:lem:restoration}{C.2}\\
Facts~\ref{tb:gauss-coord}--\ref{tb:heat} &
Subsection~\PN{sec:outline-path}{2.2} and \PR{Lemma}{appx-an:lem:identity}{C.1} (heat equation);
Step~1 of \PR{Lemma}{appx-an:lem:restoration}{C.2} (integration by parts);
\PR{Lemma}{appx-lem:endpoint-gap}{B.3}; Steps~1--2 of \PR{Lemma}{appx-an:lem:surface}{C.5}\\
Facts~\ref{tb:nuclear}, \ref{tb:logsum} &
Subsection~\PN{sec:outline-late}{2.4}; Step~4 of \PR{Lemma}{appx-an:lem:restoration}{C.2};
Steps~1 and~3 of \PR{Lemma}{appx-an:lem:matrix}{C.4}\\
Facts~\ref{tb:prekopa}--\ref{tb:lc-weights} &
Subsection~\PN{sec:outline-late}{2.4}; Step~1 of \PR{Lemma}{appx-an:lem:matrix}{C.4}\\
Facts~\ref{tb:facets}--\ref{tb:projection} &
the geometric paragraph and Steps~1--2 of \PR{Lemma}{appx-an:lem:surface}{C.5}\\
Facts~\ref{tb:hausdorff}, \ref{tb:divergence} &
the definition of \(\surf\); Step~1 of \PR{Lemma}{appx-an:lem:surface}{C.5}\\
Fact~\ref{tb:tubes} &
Steps~2 and~3 of \PR{Lemma}{appx-an:lem:surface}{C.5}\\
Facts~\ref{tb:coarea}, \ref{tb:variation} &
the paragraph on first variation after \PR{Lemma}{appx-an:lem:surface}{C.5}; Step~4 of
\PR{Lemma}{appx-an:lem:matrix}{C.4}\\
\end{tabular}
\end{center}

\subsection{Sources}\label{tb:sec:sources}

\begin{description}\small\setlength{\itemsep}{0pt}
\item[{[BL76]}] H.~J. Brascamp and E.~H. Lieb (1976). On extensions of the
Brunn--Minkowski and Pr\'ekopa--Leindler theorems, including inequalities
for log concave functions, and with an application to the diffusion
equation. \emph{Journal of Functional Analysis} \textbf{22}, 366--389.
\item[{[Com74]}] L. Comtet (1974). \emph{Advanced Combinatorics}. Reidel,
Dordrecht.
\item[{[CT06]}] T.~M. Cover and J.~A. Thomas (2006). \emph{Elements of
Information Theory}, 2nd edition. Wiley.
\item[{[EG15]}] L.~C. Evans and R.~F. Gariepy (2015). \emph{Measure Theory
and Fine Properties of Functions}, revised edition. CRC Press.
\item[{[Fol99]}] G.~B. Folland (1999). \emph{Real Analysis: Modern
Techniques and Their Applications}, 2nd edition. Wiley.
\item[{[Gar02]}] R.~J. Gardner (2002). The Brunn--Minkowski inequality.
\emph{Bulletin of the American Mathematical Society} \textbf{39}, 355--405.
\item[{[GZ]}] Z. Gao and W. Zhang (2026). Cubic-root Gaussian approximation
under unrestricted covariance. arXiv:2608.30221v1.
\item[{[Lei72]}] L. Leindler (1972). On a certain converse of H\"older's
inequality II. \emph{Acta Scientiarum Mathematicarum (Szeged)} \textbf{33},
217--223.
\item[{[Pre71]}] A. Pr\'ekopa (1971). Logarithmic concave measures with
application to stochastic programming. \emph{Acta Scientiarum Mathematicarum
(Szeged)} \textbf{32}, 301--316.
\item[{[Pre73]}] A. Pr\'ekopa (1973). On logarithmic concave measures and
functions. \emph{Acta Scientiarum Mathematicarum (Szeged)} \textbf{34},
335--343.
\item[{[Roc70]}] R.~T. Rockafellar (1970). \emph{Convex Analysis}. Princeton
University Press.
\item[{[RW98]}] R.~T. Rockafellar and R.~J.-B. Wets (1998). \emph{Variational
Analysis}. Springer, Berlin.
\item[{[Ste81]}] C.~M. Stein (1981). Estimation of the mean of a multivariate
normal distribution. \emph{The Annals of Statistics} \textbf{9}, 1135--1151.
\item[{[Zie95]}] G.~M. Ziegler (1995). \emph{Lectures on Polytopes}.
Springer, New York.
\end{description}

%
% after Appendix E and before \end{document}. The section then becomes
% Appendix F, and every reference to the paper below becomes a \ref to the
% paper's own labels. In the stand-alone version the same references print
% the numbers that the labels have in the paper on main (commit 12b6efd).
% The file defines the environment `fact` if the host does not, and defines
% its macros with \providecommand, so it leaves the host's definitions as
% they are. The list of sources at the end can be merged into the paper's
% bibliography.

\makeatletter
\providecommand{\E}{\mathbb E}
\providecommand{\Pp}{\mathbb P}
\providecommand{\R}{\mathbb R}
\providecommand{\tr}{\operatorname{tr}}
\providecommand{\op}{\mathrm{op}}
\providecommand{\norm}[1]{\left\lVert#1\right\rVert}
\providecommand{\ind}{\mathbf 1}
\providecommand{\dd}{\,d}
\providecommand{\Hh}{H_2}
\providecommand{\surf}{\mathcal J}
\providecommand{\relint}{\operatorname{relint}}
\providecommand{\aff}{\operatorname{aff}}
\providecommand{\cone}{\operatorname{cone}}
\providecommand{\dist}{\operatorname{dist}}
\providecommand{\prj}{\operatorname{pr}}
\providecommand{\Haus}{\mathcal H}
\providecommand{\TBc}[1]{[#1]}
% \PN{label}{C.5} prints a number of the paper (a \ref in the paper);
% \PR{Lemma}{label}{C.5} prints "Lemma C.5"; \PE{label}{C.21} an equation.
\providecommand{\PN}[2]{\ifdefined\toolboxinpaper\ref{#1}\else#2\fi}
\providecommand{\PR}[3]{#1~\PN{#2}{#3}}
\providecommand{\PE}[2]{(\PN{#1}{#2})}
\providecommand{\usedin}[1]{\par\smallskip\noindent{\small\emph{Used in:} #1}\par\smallskip}
\@ifundefined{c@fact}{\newtheorem{fact}{Fact}[section]}{}
\makeatother

\section{Toolbox: facts from analysis and geometry}\label{app:toolbox}

This appendix collects the facts from analysis, convex geometry, and
measure theory on which the proofs in Appendices~\PN{appx-sec:smoothing}{B}
and~\PN{appx-sec:support}{C} rely and which a mathematical statistician may
not have at hand. Every surface in those proofs is a finite union of flat
pieces of hyperplanes, so no differential geometry beyond the flat case is
needed. On a hyperplane the area formula is Fubini's theorem, the
Gauss--Green formula (the divergence theorem) is the fundamental theorem of
calculus along parallel lines, and the coarea formula is Fubini's theorem
applied piece by piece; the proofs below carry out each of these
reductions. Each fact is stated in the form in which
the proofs use it, and is followed by the places that use it and by a proof
or a reference. A proof is given when it takes a few lines, and a standard
source, listed at the end of the appendix, is cited otherwise. The facts are
standard results and not claims of the paper, so they carry their own
numbering. The symbols in the lines headed ``Used in'' are those of the
paper's appendices.

Subsection~\ref{tb:sec:calc} treats integration and calculus,
Subsection~\ref{tb:sec:gauss} the Gaussian identities, and
Subsection~\ref{tb:sec:matrix} the matrix and entropy inequalities.
Subsection~\ref{tb:sec:logconcave} covers log-concavity and
Pr\'ekopa's theorem. Subsections~\ref{tb:sec:poly}--\ref{tb:sec:coarea}
treat polyhedra, from faces and normal cones to the divergence theorem,
tubes, and the coarea formula, in the order in which
\PR{Lemma}{appx-an:lem:surface}{C.5} needs them, and
Subsection~\ref{tb:sec:index} lists where each fact is used.

\subsection{Notation}\label{tb:sec:notation}

The space \(\R^n\) carries the inner product \(x\cdot y\) and the Euclidean
norm \(\norm x\); in the proofs \(n\) is the rank \(r\) of the covariance,
or \(r-1\) when a hyperplane of \(\R^r\) is identified with
\(\R^{r-1}\). We write \(dx\) for Lebesgue measure, \(\varphi_n\) and
\(\Phi_n\) for the standard Gaussian density and law on \(\R^n\), and \(G\)
for a standard Gaussian vector. For a linear subspace \(L\), \(\varphi_L\)
is the standard Gaussian density on \(L\) with respect to Lebesgue measure
on \(L\). For a set \(S\), \(\ind_S\) is its indicator function, \(\aff S\)
is the smallest affine subspace containing it, the dimension of \(S\) is
that of \(\aff S\), and \(\relint S\) is the interior of \(S\) relative to
\(\aff S\). The distance from a point \(y\) to a set \(S\) is
\(\dist(y,S)=\inf_{x\in S}\norm{y-x}\).

For \(f\in C^q(\R^n)\), the derivative of order \(q\) in the directions
\(z_1,\ldots,z_q\in\R^n\) is
\[
 D^qf(x)[z_1,\ldots,z_q]
 =\partial_{s_1}\cdots\partial_{s_q}f(x+s_1z_1+\cdots+s_qz_q)\Big|_{s=0}
 =\sum_{i_1,\ldots,i_q}\partial_{i_1\cdots i_q}f(x)\,
   z_{1,i_1}\cdots z_{q,i_q},
\]
which is symmetric and multilinear in the directions. We write
\(D^qf(x)[z^{\otimes q}]\) when all directions equal \(z\), and
\(D^1f=\nabla f\). The Loewner order \(M\preceq N\) means that
\(N-M\) is positive semidefinite, \(\norm\Xi_\op\) is the operator norm
of a matrix \(\Xi\), and \(\norm M_*\) is the \emph{nuclear norm}, which for
a symmetric matrix \(M\) with eigenvalues \(\mu_i\) equals \(\sum_i|\mu_i|\).

\subsection{Integration and calculus}\label{tb:sec:calc}

\begin{fact}[Fubini and Tonelli]\label{tb:fubini}
Let \((S_1,\mu_1)\) and \((S_2,\mu_2)\) be \(\sigma\)-finite measure
spaces and let \(f\) be a measurable function on \(S_1\times S_2\). If
\(f\ge0\), then \(s_1\mapsto\int f(s_1,s_2)\,\mu_2(ds_2)\) is measurable
and
\[
 \int\Bigl(\int f\,d\mu_2\Bigr)d\mu_1
 =\int\Bigl(\int f\,d\mu_1\Bigr)d\mu_2
 =\int f\,d(\mu_1\otimes\mu_2)\in[0,\infty].
\]
If \(\int|f|\,d(\mu_1\otimes\mu_2)<\infty\), the same equalities hold, and
the inner integrals are finite almost everywhere.
\end{fact}

Three special cases occur in the proofs. If \(X\) and \(Y\) are
independent random elements and either \(f\ge0\) or \(f(X,Y)\) is
integrable, then
\[
 \E f(X,Y)=\E g(X),\qquad g(x)=\E f(x,Y),
\]
because the joint law is the product of the marginal laws. This is the
formal content of every argument that is carried out for a fixed value of
an independent component such as \(V_0\) and then averaged over it. With
\(\mu_2\) the counting measure, the theorem says that
\(\E\sum_k\xi_k=\sum_k\E\xi_k\) for nonnegative terms, and for terms with
\(\sum_k\E|\xi_k|<\infty\). For an orthogonal decomposition
\(\R^n=L\oplus L^\perp\), Lebesgue measure on \(\R^n\) is the product of the
Lebesgue measures on \(L\) and on \(L^\perp\), so that
\(\int_{\R^n}f\,dx=\int_{L^\perp}\int_Lf(l+w)\,dl\,dw\).

The appeal to Tonelli's theorem in Step~3 of the proof of
\PR{Lemma}{appx-an:lem:surface}{C.5}, whose notation we use in this
paragraph, is of the first kind. With \(V_0\) and
\(G\) independent, \(\Phi_r(\mathcal O_\iota(V_0))\) is the nonnegative
integral \(\int\ind\{g\in\mathcal O_\iota(V_0)\}\varphi_r(g)\,dg\), so
\(\E\sum_\iota R_\iota^3\,\Phi_r(\mathcal O_\iota(V_0))\) equals
\(\int\sum_\iota R_\iota^3\,\Pp\{g\in\mathcal O_\iota(V_0)\}\varphi_r(g)\,dg\)
with no check of integrability, since every term is nonnegative, and the
probability can then be bounded for each fixed \(g\).

\usedin{conditioning on \(V_0\) or on the \(\sigma\)-field \(\mathcal H\) in
\PR{Lemmas}{appx-an:lem:restoration}{C.2}, \PN{appx-an:lem:matrix}{C.4}
and~\PN{appx-an:lem:surface}{C.5}; conditioning on the exponentials in
\PR{Lemmas}{appx-lem:envelope}{B.1} and~\PN{appx-lem:endpoint-free}{B.2}
and on the candidate pairs in \PR{Appendix}{appx-app:concentration}{D};
interchanges of sums and expectations at the end of the proof of
\PR{Lemma}{appx-an:lem:restoration}{C.2}; Step~3 of the proof of
\PR{Lemma}{appx-an:lem:surface}{C.5}; iterated integrals over a hyperplane
and its normal in Subsections~\ref{tb:sec:surface}--\ref{tb:sec:coarea}.}
\emph{Source.} \TBc{Fol99}, Section~2.5.

\begin{fact}[Differentiation under the integral sign; Taylor's formula]
\label{tb:taylor}
(a) Let \(I\) be an open interval and let \(f(\lambda,\omega)\) be
integrable in \(\omega\) and differentiable in \(\lambda\in I\), with
\(|\partial_\lambda f(\lambda,\omega)|\le g(\omega)\) for an integrable
function \(g\). Then \(\frac d{d\lambda}\E f(\lambda,\cdot)
=\E\,\partial_\lambda f(\lambda,\cdot)\) on \(I\).

(b) For \(f\in C^q(\R^n)\) and \(x,v\in\R^n\),
\[
 f(x+v)=\sum_{k=0}^{q-1}\frac1{k!}D^kf(x)[v^{\otimes k}]
 +\int_0^1\frac{(1-s)^{q-1}}{(q-1)!}D^qf(x+sv)[v^{\otimes q}]\,ds.
\]
For \(q=1\) this is \(f(x+v)-f(x)=\int_0^1Df(x+sv)[v]\,ds\).
\end{fact}

\begin{proof}
For (a), the mean value theorem bounds the difference quotients by \(g\),
and they converge pointwise, so dominated convergence applies. For (b),
the function \(s\mapsto f(x+sv)\) is of class \(C^q\) with \(k\)th derivative
\(D^kf(x+sv)[v^{\otimes k}]\); apply the one-dimensional Taylor formula with
integral remainder at \(s=0\) and evaluate it at \(s=1\).
\end{proof}

\usedin{\PR{Lemma}{appx-an:lem:identity}{C.1}, where the derivative of
\(P_A(\lambda)=\E F(U_\lambda)\) is taken under the expectation (the
summands are finitely many and bounded) and Steps~1 and~3 use (b), to degree
three with the fourth-order integral remainder, to degree two for the
derivative, and to first order for the difference of third derivatives;
\PR{Lemma}{appx-lem:endpoint-gap}{B.3}, where (b) with \(q=1\) writes
\(\E[W\cdot\nabla\Psi(W)]\) as an integral of the Hessian along the segment
from \(W_i\) to \(W\).}
\emph{Source.} \TBc{Fol99}, Section~2.3.

\begin{fact}[Fa\`a di Bruno's formula]\label{tb:faa}
Let \(\phi\) be a smooth function on an open interval of \(\R\) and let
\(g:\R^n\to\R\) be smooth with values in that interval. For \(q\ge1\), let
\(\Pi_q\) be the set of partitions \(\mathcal P\) of \(\{1,\ldots,q\}\) into
nonempty blocks, let \(|\mathcal P|\) be the number of blocks, and for a
block \(J\) let \(z_J\) be the list of directions \(z_a\) with
\(a\in J\). Then
\[
 D^q(\phi\circ g)(x)[z_1,\ldots,z_q]
 =\sum_{\mathcal P\in\Pi_q}\phi^{(|\mathcal P|)}(g(x))
  \prod_{J\in\mathcal P}D^{|J|}g(x)[z_J].
\]
\end{fact}

\begin{proof}
For \(q=1\) this is the chain rule. Differentiate the right side for \(q\)
in the direction \(z_{q+1}\). By the product rule, the derivative either
falls on \(\phi^{(|\mathcal P|)}(g(x))\), which gives
\(\phi^{(|\mathcal P|+1)}(g(x))\,Dg(x)[z_{q+1}]\) and adds the block
\(\{q+1\}\) to \(\mathcal P\), or falls on one factor
\(D^{|J|}g(x)[z_J]\), which adds \(q+1\) to the block \(J\). Every partition
of \(\{1,\ldots,q+1\}\) arises exactly once in this way, by removing
\(q+1\) from its block.
\end{proof}

\begin{fact}[Derivatives of the logistic approximation]\label{tb:logistic}
Let \(\nu_1,\ldots,\nu_m\in\R^n\), \(a_1,\ldots,a_m\in\R\), \(h>0\), and
\[
 w_\alpha(x)=e^{(\nu_\alpha\cdot x-a_\alpha)/h},\qquad
 T=\sum_\alpha w_\alpha,\qquad
 F=\frac1{1+T},\qquad
 \pi_\alpha=\frac{w_\alpha}{1+T}=Fw_\alpha .
\]
For directions \(z_1,\ldots,z_q\), with the blocks \(J_1,\ldots,J_{|\mathcal P|}\)
of \(\mathcal P\) listed in any fixed order,
\[
 D^qF(x)[z_1,\ldots,z_q]
 =h^{-q}\sum_{\mathcal P\in\Pi_q}(-1)^{|\mathcal P|}|\mathcal P|!
 \sum_{\alpha_1,\ldots,\alpha_{|\mathcal P|}}
 F(x)\prod_{b=1}^{|\mathcal P|}\pi_{\alpha_b}(x)
 \prod_{b=1}^{|\mathcal P|}\prod_{a\in J_b}(\nu_{\alpha_b}\cdot z_a),
\]
and \(\sum_{\mathcal P\in\Pi_q}|\mathcal P|!\le q!\,2^{q-1}\). Consequently
\(|D^qF(x)[z_1,\ldots,z_q]|\le q!\,2^{q-1}h^{-q}\,F(x)\prod_{a=1}^q
\max_\alpha|\nu_\alpha\cdot z_a|\).
\end{fact}

\begin{proof}
Apply Fact~\ref{tb:faa} with \(\phi(\tau)=(1+\tau)^{-1}\) and \(g=T\). Then
\(\phi^{(k)}(\tau)=(-1)^kk!\,(1+\tau)^{-k-1}\), and since each \(w_\alpha\) is
an exponential of an affine function,
\(D^lT(x)[z_1,\ldots,z_l]=h^{-l}\sum_\alpha w_\alpha(x)\prod_{a\le l}
(\nu_\alpha\cdot z_a)\). Substituting, and using
\((1+T)^{-|\mathcal P|-1}\prod_bw_{\alpha_b}=F\prod_b\pi_{\alpha_b}\) and
\(\sum_b|J_b|=q\), gives the formula. The number
\(\sum_{\mathcal P}|\mathcal P|!\) counts ordered partitions
\((J_1,\ldots,J_k)\) of \(\{1,\ldots,q\}\). Writing the elements of
\(J_1\), then of \(J_2\), and so on, in increasing order inside each block,
gives a permutation together with the set of positions at which a new block
starts, which is a subset of the \(q-1\) gaps between consecutive entries.
Every ordered partition arises from at least one of these
\(q!\,2^{q-1}\) pairs. For the bound, the sum over \(\alpha_b\) of
\(\pi_{\alpha_b}\) is \(1-F\le1\), and each product of projections is
at most \(\prod_a\max_\alpha|\nu_\alpha\cdot z_a|\).
\end{proof}

\usedin{Subsection~\PN{appx-sec:mixture}{B.2}, where the first display
and the bound are \PE{appx-eq:partition-derivative}{B.1} and
\PE{appx-eq:all-orders}{B.2} (the text there cites the softmax cumulant formula,
\TBc{GZ}, Lemma~13); the envelope~\PE{appx-eq:local-envelope}{B.3} is
obtained by summing the same expansion over coordinate tuples;
\PR{Lemma}{appx-an:lem:restoration}{C.2}, Step~1, with \(u_\alpha\) in place
of \(\nu_\alpha\).}
\emph{Source.} \TBc{Com74}; the proof above is complete.

\subsection{Gaussian identities}\label{tb:sec:gauss}

\begin{fact}[Gaussian density in orthogonal coordinates]\label{tb:gauss-coord}
(a) Let \(\R^n=L\oplus L^\perp\) be an orthogonal decomposition. Then
\(\varphi_n(l+w)=\varphi_L(l)\varphi_{L^\perp}(w)\) for \(l\in L\) and
\(w\in L^\perp\), so the projections of a standard Gaussian vector onto
\(L\) and \(L^\perp\) are independent standard Gaussian vectors. In
particular, for a unit vector \(u\), \(a\in\R\), and \(y\perp u\),
\(\varphi_n(au+y)=\varphi_1(a)\varphi_{n-1}(y)\).

(b) For \(x,v\in\R^n\),
\(\varphi_n(x+v)=\varphi_n(x)\exp\{-x\cdot v-\tfrac12\norm v^2\}\).
\end{fact}

\begin{proof}
Both follow from \(\norm{l+w}^2=\norm l^2+\norm w^2\) and from expanding
\(\norm{x+v}^2\).
\end{proof}

\usedin{\PR{Lemma}{appx-an:lem:mean}{C.3}, where the normal and tangential
components of \(G\) relative to a face are independent;
\PR{Lemma}{appx-an:lem:surface}{C.5}, Steps~1 and~2, for the factorization of
\(\varphi_r\) on a facet and on a ridge and for the density ratio on a tube.}

\begin{fact}[Gaussian integration by parts]\label{tb:stein}
Let \(G\sim N(0,I_n)\) and let \(f\in C^2(\R^n)\) have \(f\) and its partial
derivatives of orders at most two of polynomial growth. Then
\begin{enumerate}
\item[(a)] \(\E[Gf(G)]=\E\nabla f(G)\);
\item[(b)] \(\E[(GG^\top-I_n)f(G)]=\E D^2f(G)\); equivalently, for \(z\in\R^n\),
\(\E D^2f(G)[z,z]=\E\bigl[\tr\{(GG^\top-I_n)zz^\top\}f(G)\bigr]\), where
\(\tr\{(GG^\top-I_n)zz^\top\}=(z\cdot G)^2-\norm z^2\);
\item[(c)] if \(v\in\R^n\), \(\sigma>0\), \(U=v+\sigma G\), and
\(F\in C^{s+2}(\R^n)\) has derivatives of polynomial growth, then for
\(z,z_1,\ldots,z_s\in\R^n\),
\[
 \E D^{s+2}F(U)[z,z,z_1,\ldots,z_s]
 =\sigma^{-2}\E\bigl[\tr\{(GG^\top-I_n)zz^\top\}\,
 D^sF(U)[z_1,\ldots,z_s]\bigr].
\]
\end{enumerate}
\end{fact}

\begin{proof}
For (a) in one coordinate, \(x\varphi_1(x)=-\varphi_1'(x)\), so integration by
parts gives
\[\int xf(x)\varphi_1(x)\,dx=\int f'(x)\varphi_1(x)\,dx,\]
and the boundary terms vanish because of the growth condition. The other coordinates
are handled by Fact~\ref{tb:fubini}. For (b), apply (a) to \(\partial_if\) and
to \(G_kf\): \(\E[G_iG_kf(G)]=\E\partial_i(G_kf)(G)=\delta_{ik}\E f(G)
+\E[G_k\partial_if(G)]=\delta_{ik}\E f(G)+\E\partial_{ik}f(G)\). For (c), the
function \(f(g)=D^sF(v+\sigma g)[z_1,\ldots,z_s]\) has
\(D^2f(g)[z,z]=\sigma^2D^{s+2}F(v+\sigma g)[z,z,z_1,\ldots,z_s]\); apply (b).
\end{proof}

\usedin{Subsection~\PN{sec:outline-late}{2.4} (the display after
``integration by parts against the covariance of the summands''),
\PR{Lemma}{appx-an:lem:restoration}{C.2}, Step~1, where (c) is applied
conditionally on \(\mathcal H\) with \(v=V_0+v_i\) and \(z=Z_i\);
\PR{Lemma}{appx-lem:endpoint-gap}{B.3}, for \(\E[Z\cdot\nabla\Psi(Z)]
=\E\tr(\Sigma D^2\Psi(Z))\), which is (a) with \(Z=\Gamma G\) for a factor
\(\Gamma\) of the covariance.}
\emph{Source.} Stein~\TBc{Ste81}; the proof above is complete.

\begin{fact}[Heat equation for the Gaussian component]\label{tb:heat}
Let \(Z\sim N(0,\Sigma)\) in \(\R^p\), where \(\Sigma\) may be singular, and
let \(f\in C^2(\R^p)\) be bounded with bounded first and second derivatives.
Then, for \(x\in\R^p\) and \(\lambda>0\),
\[
 \frac d{d\lambda}\E f(x+\sqrt\lambda Z)
 =\tfrac12\E\tr\{\Sigma D^2f(x+\sqrt\lambda Z)\}.
\]
\end{fact}

\begin{proof}
Write \(Z=\Gamma G\) with \(\Gamma\Gamma^\top=\Sigma\) and \(G\) standard
Gaussian. By Fact~\ref{tb:taylor}(a), which applies because the difference
quotients are bounded by \(\norm\nabla f_\infty\norm{\Gamma G}/(2\sqrt\lambda)\)
locally uniformly in \(\lambda>0\), the derivative equals
\((2\sqrt\lambda)^{-1}\E[G\cdot\Gamma^\top\nabla f(x+\sqrt\lambda\Gamma G)]\).
Fact~\ref{tb:stein}(a) applied to the function \(g\mapsto
\Gamma^\top\nabla f(x+\sqrt\lambda\Gamma g)\) turns this into
\((2\sqrt\lambda)^{-1}\sqrt\lambda\,\E\tr\{\Gamma^\top D^2f\,\Gamma\}
=\frac12\E\tr\{\Sigma D^2f\}\).
\end{proof}

\usedin{Subsection~\PN{sec:outline-path}{2.2} and the proof of
\PR{Lemma}{appx-an:lem:identity}{C.1}, Step~2, where the Gaussian part of the
path contributes \(\frac12\sum_iM_i:\E D^2F(U_\lambda)\).}

\subsection{Matrix and entropy inequalities}\label{tb:sec:matrix}

\begin{fact}[Duality between the operator and nuclear norms]\label{tb:nuclear}
(a) For every square matrix \(\Xi\) and every symmetric matrix \(M\),
\(|\tr(\Xi M)|\le\norm\Xi_\op\norm M_*\). Moreover
\(\norm M_*=\max\{\tr(\Xi M):\norm\Xi_\op\le1\}\), so that \(\norm\cdot_*\) is
a norm on symmetric matrices.

(b) If \(M=P-N\) with \(P\succeq0\) and \(N\succeq0\), then
\(\norm M_*\le\tr P+\tr N=2\tr P-\tr M\).
\end{fact}

\begin{proof}
For (a), let \(M=\sum_i\mu_ie_ie_i^\top\) be a spectral decomposition. Then
\(\tr(\Xi M)=\sum_i\mu_i\,e_i^\top\Xi e_i\) and \(|e_i^\top\Xi e_i|\le
\norm\Xi_\op\), which gives the inequality. Equality is attained at
\(\Xi=\sum_i\operatorname{sgn}(\mu_i)e_ie_i^\top\). A supremum of linear
functionals satisfies the triangle inequality. For (b), the norm of a
positive semidefinite matrix is its trace, so the triangle inequality gives
\(\norm M_*\le\tr P+\tr N\), and \(\tr N=\tr P-\tr M\).
\end{proof}

\usedin{Subsection~\PN{sec:outline-late}{2.4}; \PR{Lemma}{appx-an:lem:restoration}{C.2},
Step~4; \PR{Lemma}{appx-an:lem:matrix}{C.4}, Step~1, where
(b) is applied with \(P=p_\alpha\mu_\alpha\mu_\alpha^\top\) and
\(N=p_\alpha(I_r-\Sigma_\alpha)\), and gives
\PE{appx-an:matrix-trace-reduction}{C.17}.}

\begin{fact}[Log-sum inequality and a concave function]\label{tb:logsum}
(a) For positive numbers \(y_1,\ldots,y_l\) and \(z_1,\ldots,z_l\) with sums
\(y\) and \(z\),
\(\sum_jy_j\log(z_j/y_j)\le y\log(z/y)\).

(b) The function \((x,y)\mapsto x\log(y/x)\) is concave on \((0,\infty)^2\).
Hence, for positive integrable random variables \(X,Y\),
\(\E[X\log(Y/X)]\le\E X\,\log(\E Y/\E X)\).
\end{fact}

\begin{proof}
For (a), Jensen's inequality for the concave function \(\log\), with the
weights \(y_j/y\), gives \(\sum_j(y_j/y)\log(z_j/y_j)\le\log\sum_j(y_j/y)(z_j/y_j)
=\log(z/y)\). For (b), the function \(f(x,y)=x\log y-x\log x\) has Hessian
\(\begin{pmatrix}-1/x&1/y\\1/y&-x/y^2\end{pmatrix}\), whose determinant is
zero and whose trace is negative, so the Hessian is negative semidefinite;
the second statement is Jensen's inequality for the concave function
\(f\) and the random vector \((X,Y)\).
\end{proof}

\usedin{\PR{Lemma}{appx-an:lem:matrix}{C.4}, Step~3: (a) adds the reference
masses of the tuples that share a first face, and (b) averages over
\(V_0\) and over the faces.}
\emph{Source.} \TBc{CT06}, Theorem~2.7.1, for (a).

\subsection{Log-concavity}\label{tb:sec:logconcave}

A function \(f:\R^n\to[0,\infty)\) is \emph{log-concave} if
\(f((1-t)x+ty)\ge f(x)^{1-t}f(y)^t\) for all \(x,y\) and \(t\in[0,1]\);
when \(f>0\) this says that \(\log f\) is concave.

\begin{fact}[Pr\'ekopa--Leindler inequality and Pr\'ekopa's theorem]
\label{tb:prekopa}
(a) Let \(0<t<1\) and let \(f,g,h:\R^n\to[0,\infty)\) be bounded,
continuous, and integrable, with
\(h((1-t)x+ty)\ge f(x)^{1-t}g(y)^t\) for all \(x,y\). Then
\(\int h\ge(\int f)^{1-t}(\int g)^t\).

(b) Let \(f:\R^a\times\R^b\to[0,\infty)\) be bounded, continuous, and
log-concave, with \(f(z,\cdot)\) integrable for every \(z\). Then
\(z\mapsto\int f(z,y)\,dy\) is log-concave on \(\R^a\).
\end{fact}

\begin{proof}
(a) For \(n=1\), assume first that \(\sup f=\sup g=1\); the general case
follows by dividing \(f,g,h\) by \(\sup f\), \(\sup g\), and
\((\sup f)^{1-t}(\sup g)^t\). For \(0<s<1\) the level sets
\(\{f\ge s\}\) and \(\{g\ge s\}\) are nonempty, and the hypothesis gives
\(\{h\ge s\}\supseteq(1-t)\{f\ge s\}+t\{g\ge s\}\). In one dimension
\(|A+B|\ge|A|+|B|\) for nonempty compact sets \(A\) and \(B\), because
\(A+B\) contains the translates \(A+\min B\) and \(\max A+B\), which meet in
one point; the inequality passes to measurable sets by approximation from
inside by compact sets. Therefore
\(|\{h\ge s\}|\ge(1-t)|\{f\ge s\}|+t|\{g\ge s\}|\), and integration over
\(s\in(0,1)\) gives \(\int h\ge(1-t)\int f+t\int g\ge(\int f)^{1-t}(\int g)^t\),
by the inequality between the arithmetic and geometric means. For \(n\ge2\),
apply the case of dimension \(n-1\) to the sections of \(f,g,h\) at the last
coordinates \(s_0\), \(s_1\), and \((1-t)s_0+ts_1\), and then the case of
dimension one to the integrals of the sections, with Fubini's theorem
(Fact~\ref{tb:fubini}) connecting the steps; the routine measurability
checks are omitted here and can be found in~\TBc{Gar02}.

(b) Fix \(z_0,z_1\) and put \(z_t=(1-t)z_0+tz_1\). Log-concavity of \(f\)
gives \(f(z_t,(1-t)y_0+ty_1)\ge f(z_0,y_0)^{1-t}f(z_1,y_1)^t\), so (a) applies
to \(f(z_0,\cdot)\), \(f(z_1,\cdot)\), and \(f(z_t,\cdot)\).
\end{proof}

\emph{Source.} Pr\'ekopa~\TBc{Pre71,Pre73}, Leindler~\TBc{Lei72}; a
survey with the proof of the general case is~\TBc{Gar02}, Section~7. The
paper cites Theorem~6 of~\TBc{Pre73}.

\begin{fact}[Gaussian laws tilted by a log-concave weight]\label{tb:tilt}
Let \(w:\R^r\to[0,\infty)\) be bounded, continuous, and log-concave, with
\(p=\E w(G)>0\) for \(G\sim N(0,I_r)\). Let \(\mu\) and \(K\) be the mean and
covariance matrix of the probability law with density \(w(g)\varphi_r(g)/p\).
Then \(K\preceq I_r\), and
\(\E[(GG^\top-I_r)w(G)]=p(K-I_r+\mu\mu^\top)\).
\end{fact}

\begin{proof}
The function \((z,g)\mapsto e^{-\norm{g-z}^2/2}w(g)\) is log-concave, since
\(\norm{g-z}^2\) is jointly convex in \((z,g)\), and it is bounded and
continuous. By Fact~\ref{tb:prekopa}(b), \(\ell(z)=\log\int e^{-\norm{g-z}^2/2}
w(g)\,dg\) is concave. Since \(\norm{g-z}^2=\norm g^2-2z\cdot g+\norm z^2\),
\[
 \ell(z)=\text{const}-\tfrac12\norm z^2+\log M(z),\qquad
 M(z)=\int e^{z\cdot g}w(g)\varphi_r(g)\,dg=p\,\E_\nu e^{z\cdot g},
\]
where \(\nu\) is the tilted law. The function \(M\) is finite and smooth,
because \(w\) is bounded, and the Hessian of \(\log M\) at \(0\) is the
covariance \(K\) of \(\nu\). A concave function of class \(C^2\) has a
negative semidefinite Hessian, so \(K-I_r\preceq0\). The identity is the
decomposition of the second moment \(\E_\nu gg^\top=K+\mu\mu^\top\).
\end{proof}

\usedin{Subsection~\PN{sec:outline-late}{2.4} (``Pr\'ekopa's theorem \dots\
has covariance \(K\preceq I_r\)'') and \PR{Lemma}{appx-an:lem:matrix}{C.4},
Step~1, where \(w=w_\alpha\) and \(K=\Sigma_\alpha\).}
\emph{Source.} For indicators of convex cells, the cell covariance bound
\TBc{GZ}, Lemma~11, is the same estimate \(K\preceq I_r\), obtained from the
Brascamp--Lieb inequality~\TBc{BL76}; the weights here are not indicators,
which is why the proof goes through Pr\'ekopa's theorem.

\begin{fact}[Log-concavity of the logistic weights]\label{tb:lc-weights}
With the notation of Fact~\ref{tb:logistic}, for every \(k\ge0\) and every
tuple \((\alpha_1,\ldots,\alpha_k)\), the function
\(F\pi_{\alpha_1}\cdots\pi_{\alpha_k}=\exp\{\sum_{a\le k}t_{\alpha_a}\}
(1+T)^{-k-1}\), where \(t_\alpha=(\nu_\alpha\cdot x-a_\alpha)/h\), is
log-concave on \(\R^n\), and so is its composition
\(g\mapsto F\pi_{\alpha_1}\cdots\pi_{\alpha_k}(v+\sigma g)\) with any affine map.
\end{fact}

\begin{proof}
The logarithm is a sum of the affine functions \(t_{\alpha_a}\) and
\(-(k+1)\log(1+\sum_\alpha e^{t_\alpha})\). The function
\(\log(1+\sum_\alpha e^{t_\alpha})\) is the logarithm of a sum of
exponentials of the affine functions \(0,t_1,\ldots,t_m\) and is convex by
H\"older's inequality: \(\sum_\alpha e^{(1-s)t_\alpha+st'_\alpha}\le
(\sum_\alpha e^{t_\alpha})^{1-s}(\sum_\alpha e^{t'_\alpha})^s\). Composition
with an affine map preserves concavity.
\end{proof}

\usedin{Subsection~\PN{sec:outline-late}{2.4} and Step~1 of the proof of
\PR{Lemma}{appx-an:lem:matrix}{C.4}.}

\subsection{Polyhedra: faces, normal cones, and nearest points}\label{tb:sec:poly}

A \emph{polyhedron} is a set
\[
 Q=\bigcap_{j=1}^m\{x\in\R^n:u_j\cdot x\le a_j\},
\]
with unit vectors \(u_j\), called \emph{outward normals}, and numbers
\(a_j\), called \emph{offsets}. It is closed and convex. It is
\emph{full-dimensional} if it has an interior point, and \emph{proper} if it
is neither empty nor \(\R^n\). The description is \emph{irredundant} if
deleting any one inequality changes the set; deleting redundant or repeated
inequalities one at a time turns every description into an irredundant one
of the same set. For \(j\le m\), put \(H_j=\{x:u_j\cdot x=a_j\}\) and
\(F_j=Q\cap H_j\). A \emph{face} of \(Q\) is \(Q\) itself or a set
\(Q\cap\{u\cdot x=c\}\) with \(u\ne0\) and \(Q\subseteq\{u\cdot x\le c\}\).
A \emph{facet} is a face of dimension \(n-1\), and a \emph{ridge} is a face
of dimension \(n-2\). In the plane the facets are the sides of a polygon
and the ridges its corners, in \(\R^3\) the facets are the faces of a solid
and the ridges its edges, and on the line the facets are the endpoints of an
interval, which has no ridges. We use these standard properties of faces:
the relative interiors of the faces of \(Q\) are pairwise disjoint and cover
\(Q\); a face of a face of \(Q\) is a face of \(Q\); and a face \(G'\)
of a face \(G\) with \(G'\ne G\) lies in the relative boundary
\(G\setminus\relint G\) of \(G\) (\TBc{Roc70}, Section~18; \TBc{Zie95}, Chapter~2).

\begin{fact}[Facets and ridges of a polyhedron]\label{tb:facets}
Let \(Q\) be a proper, full-dimensional polyhedron in \(\R^n\) with an
irredundant description.
\begin{enumerate}
\item[(a)] The normals \(u_1,\ldots,u_m\) are pairwise distinct, each
\(F_j\) is a facet, \(Q\) has no other facets, and
\(\partial Q=F_1\cup\cdots\cup F_m\).
\item[(b)] For \(j\ne k\), \(F_j\cap F_k\subseteq H_j\cap H_k\), and
\(H_j\cap H_k\) is empty if \(u_k=-u_j\), and an affine subspace of
dimension \(n-2\) otherwise. In particular \(F_j\cap F_k\) lies in a
finite union of affine subspaces of dimension at most \(n-2\).
\item[(c)] Each \(F_j\) is a full-dimensional polyhedron in the hyperplane
\(H_j\). Write the points of \(H_j\) as \(a_ju_j+y\) with \(y\perp u_j\), and,
for \(k\ne j\) with \(u_k\ne\pm u_j\), put
\(\alpha_{jk}=\arccos(u_j\cdot u_k)\in(0,\pi)\),
\[
 \nu_{k\mid j}=\frac{u_k-(u_j\cdot u_k)u_j}{\sin\alpha_{jk}},\qquad
 d_{k\mid j}=\frac{a_k-(u_j\cdot u_k)a_j}{\sin\alpha_{jk}} .
\]
Then \(F_j\) is described in the coordinate \(y\) by the inequalities
\(\nu_{k\mid j}\cdot y\le d_{k\mid j}\), and its facets are exactly the
ridges \(E_{jk}=F_j\cap F_k\) of \(Q\) of dimension \(n-2\); the outward
unit normal of \(F_j\) at \(E_{jk}\) is \(\nu_{k\mid j}\).
\item[(d)] Every ridge of \(Q\) lies in exactly two facets, so the ridges
correspond to the unordered pairs \(\{j,k\}\) with
\(\dim(F_j\cap F_k)=n-2\); for such a pair \(u_k\ne\pm u_j\), and the
only inequalities active at a point of \(\relint E_{jk}\) are \(j\) and \(k\).
\end{enumerate}
\end{fact}

\begin{proof}[Proof (sketch)]
(a) Fix \(j\). Irredundancy gives a point \(x_j\) with \(u_j\cdot x_j>a_j\)
and \(u_k\cdot x_j\le a_k\) for \(k\ne j\). Let \(x_0\) be an interior point of
\(Q\), and let \(x_*\) be the point where the segment from \(x_0\) to \(x_j\)
crosses \(H_j\). For \(k\ne j\), \(u_k\cdot x_*<a_k\), because \(x_*\) is a
convex combination of \(x_0\), where the inequality is strict, with positive
weight, and of \(x_j\), where it holds weakly. Hence a neighborhood of
\(x_*\) in \(H_j\) lies in \(F_j\), and \(\dim F_j=n-1\). Two inequalities
with the same normal cannot both be irredundant, since the one with the
larger offset is implied by the other. A point of \(Q\) at
which no inequality is active is interior, and at a point where
inequality \(j\) is active the points \(x+\varepsilon u_j\) lie outside \(Q\),
so \(\partial Q=\bigcup_jF_j\). If \(G=Q\cap\{u\cdot x=c\}\) is a facet and
\(x\in\relint G\), some inequality \(j\) is active at \(x\); on \(G\) the
affine function \(u_j\cdot x-a_j\le0\) attains its maximum at a relative
interior point and is therefore zero, so \(G\subseteq H_j\), and since \(G\) has
dimension \(n-1\), \(G=Q\cap H_j=F_j\).

(b) If \(u_k=-u_j\) and \(H_j\cap H_k\ne\emptyset\), then \(Q\) would lie in
the hyperplane \(H_j\), contradicting full dimension. If \(u_k\ne\pm u_j\),
the two equations are independent.

(c) The set \(F_j\) is the polyhedron in \(H_j\) cut out by the remaining
inequalities, which take the stated form in the coordinate \(y\) because
\(u_k\cdot(a_ju_j+y)=a_ju_j\cdot u_k+(u_k-(u_j\cdot u_k)u_j)\cdot y\); an
inequality with \(u_k=-u_j\) holds automatically on \(H_j\) since
\(F_j\ne\emptyset\). By (a) applied inside \(H_j\), the facets of \(F_j\) are
the faces of \(F_j\) of dimension \(n-2\), hence faces of \(Q\) of that
dimension, that is, ridges of \(Q\) in \(F_j\); each equals
\(F_j\cap H_k=F_j\cap F_k\) for the inequality \(k\) that defines it.

(d) Let \(E\) be a ridge and \(y\in\relint E\), and let \(J\) be the set of
inequalities active at \(y\). On \(E\) each affine function
\(u_i\cdot x-a_i\), \(i\in J\), is at most zero and vanishes at the relative
interior point \(y\), so it vanishes on \(E\); thus \(E\subseteq F_i\) for
\(i\in J\) and \(\dim\operatorname{span}\{u_i:i\in J\}\le2\). A single active
inequality is impossible, since then \(y\) would lie in the relative interior
of \(F_j\) and the only face of \(F_j\) containing it is \(F_j\) itself, of
dimension \(n-1\). Opposite normals \(u_i=-u_j\) are excluded by (b). Hence
the active normals span a plane \(\Pi\). Near \(y\), \(Q\) coincides with
\(y\) plus the product of \(\Pi^\perp\) and a planar cone \(K\subseteq\Pi\)
with nonempty interior, defined by the active inequalities and contained in
a wedge of opening \(\pi-\alpha_{jk}<\pi\); this cone has exactly two
boundary rays. Each active inequality \(i\) has \(\dim F_i=n-1\), so
\(F_i\) meets every neighborhood of \(y\) in a set of dimension \(n-1\), and
the line \(\{u_i\cdot v=0\}\) must therefore carry a boundary ray of \(K\).
With distinct normals this leaves exactly two active inequalities.
\end{proof}

\usedin{\PR{Lemma}{appx-an:lem:surface}{C.5}, in the paragraph on the terms
from the geometry of polyhedra and in Step~1 (``each inequality defines
exactly one facet, \dots\ two facets that share a ridge have normals that
are neither equal nor opposite'') and Step~2 (``every ridge lies in exactly
two facets'').}
\emph{Source.} \TBc{Zie95}, Chapter~2; \TBc{Roc70}, Sections~18 and~19;
the facet decomposition of \(\partial Q\) with the pairing of ridges is
the content of \TBc{GZ}, Lemma~27.

\begin{fact}[Normal cones]\label{tb:normalcone}
Let \(Q\) be a polyhedron as in Fact~\ref{tb:facets}. For \(x\in Q\) let
\(J(x)=\{j:u_j\cdot x=a_j\}\), and let
\[
 N_Q(x)=\{v\in\R^n:v\cdot(q-x)\le0\text{ for all }q\in Q\}
\]
be the \emph{normal cone} of \(Q\) at \(x\). Then
\(N_Q(x)=\{\sum_{j\in J(x)}c_ju_j:c_j\ge0\}\). In particular
\(N_Q(x)=\{0\}\) at interior points, \(N_Q(x)=\{cu_j:c\ge0\}\) for
\(x\in\relint F_j\), and \(N_Q(x)=\{c_ju_j+c_ku_k:c_j,c_k\ge0\}\), a sector of
angle \(\alpha_{jk}\) in the plane spanned by \(u_j\) and \(u_k\), for
\(x\in\relint E_{jk}\).
\end{fact}

\begin{proof}
Put \(T=\{d:u_j\cdot d\le0\text{ for }j\in J(x)\}\). A vector \(v\) lies in
\(N_Q(x)\) if and only if \(v\cdot d\le0\) for all \(d\in T\): if \(d\in T\)
then \(x+\varepsilon d\in Q\) for small \(\varepsilon>0\), because the
inactive inequalities are strict at \(x\), and conversely \(q-x\in T\) for
\(q\in Q\). By Farkas's lemma, which says that \(v\cdot d\le0\) for all
\(d\) with \(u_j\cdot d\le0\) (\(j\in J\)) if and only if \(v\) is a
nonnegative combination of the \(u_j\), \(j\in J\), this is the stated
cone. For \(x\in\relint F_j\), \(J(x)=\{j\}\), since an index \(k\ne j\) in
\(J(x)\) would put \(x\) in \(F_j\cap F_k\), a face of \(F_j\) of dimension
at most \(n-2\), and so in the relative boundary of \(F_j\). For
\(x\in\relint E_{jk}\), \(J(x)=\{j,k\}\) by Fact~\ref{tb:facets}(d).
\end{proof}

\usedin{\PR{Lemma}{appx-an:lem:surface}{C.5}, Step~2 (``the normal cone at a
relative-interior point of the ridge shared by facets \(j\) and \(k\) is the
cone spanned by \(u_j\) and \(u_k\)'') and the definition of the outward
tubes.}
\emph{Source.} \TBc{Roc70}, Sections~22 and~23, for Farkas's lemma and normal
cones; \TBc{RW98}, Chapter~6.

\begin{fact}[Nearest point]\label{tb:projection}
Let \(Q\subseteq\R^n\) be nonempty, closed, and convex. For each
\(y\in\R^n\) there is a unique point \(\prj_Q(y)\in Q\) with
\(\norm{y-\prj_Q(y)}=\dist(y,Q)\), and \(x=\prj_Q(y)\) if and only if
\(x\in Q\) and \(y-x\in N_Q(x)\). Consequently every \(y\in\R^n\) can be
written in exactly one way as \(y=x+v\) with \(x\in Q\) and \(v\in N_Q(x)\),
and then \(\norm v=\dist(y,Q)\). A direction \(v\) lies in \(N_Q(x)\) if and
only if \(x\) is the nearest point of \(x+tv\) in \(Q\) for every \(t>0\).
\end{fact}

\begin{proof}
For \(x,q\in Q\), \(\norm{y-q}^2=\norm{y-x}^2+\norm{x-q}^2-2(y-x)\cdot(q-x)\).
If \(y-x\in N_Q(x)\), the last term is nonnegative, so
\(\norm{y-q}^2\ge\norm{y-x}^2+\norm{x-q}^2>\norm{y-x}^2\) for \(q\ne x\); thus
\(x\) is the unique nearest point. Conversely, if \((y-x)\cdot(q-x)>0\) for
some \(q\in Q\), then moving from \(x\) toward \(q\) decreases the distance
to \(y\) for small steps, by the displayed identity applied to
\(x+\varepsilon(q-x)\in Q\), so \(x\) is not nearest. A nearest point exists
because \(Q\) is closed. The last statement follows because \(N_Q(x)\) is a
cone.
\end{proof}

\usedin{\PR{Lemma}{appx-an:lem:surface}{C.5}, Step~2, to show that the tubes
are disjoint (``their points have distinct metric projections'').}
\emph{Source.} \TBc{RW98}, Chapter~6 (Moreau's decomposition); \TBc{GZ},
proof of Lemma~28.

\subsection{Surface measure, Gaussian surface integrals, and the
divergence theorem}\label{tb:sec:surface}

\begin{fact}[Hausdorff measure on flat pieces]\label{tb:hausdorff}
For \(0\le k\le n\), the \(k\)-dimensional Hausdorff measure \(\Haus^k\) on
\(\R^n\) is the Borel measure
\[
 \Haus^k(S)=\lim_{\epsilon\downarrow0}\inf\Bigl\{\sum_i\omega_k
 \Bigl(\tfrac12\operatorname{diam}C_i\Bigr)^k:S\subseteq\bigcup_iC_i,\
 \operatorname{diam}C_i\le\epsilon\Bigr\},
\]
where \(\omega_k\) is the volume of the unit ball of \(\R^k\). (a) On an
affine subspace of dimension \(k\), \(\Haus^k\) is \(k\)-dimensional Lebesgue
measure; in particular \(\Haus^0\) is counting measure and \(\Haus^n\) is
Lebesgue measure on \(\R^n\). (b) \(\Haus^k\) is invariant under isometries.
(c) A set contained in a finite union of affine subspaces of dimension less
than \(k\) has \(\Haus^k\)-measure zero.
\end{fact}

Part (c) follows from (a), since each such subspace lies in a
\(k\)-dimensional affine subspace in which it is a Lebesgue null set. For a
proper, full-dimensional polyhedron \(Q\) with an irredundant description,
and for a function \(f\ge0\) or \(f\) integrable on \(\partial Q\), the
\emph{surface integral} is
\[
 \int_{\partial Q}f\,d\Haus^{n-1}=\sum_{j=1}^m\int_{F_j}f\,d\Haus^{n-1},
\]
where on each facet \(\Haus^{n-1}\) is Lebesgue measure on the hyperplane
\(H_j\) by part~(a). By Fact~\ref{tb:facets}(b) and part~(c), distinct
facets overlap in sets of \(\Haus^{n-1}\)-measure zero, so the sum is the
integral of \(f\) over the set \(\partial Q\) with respect to
\(\Haus^{n-1}\). It depends on the set \(Q\) and not on its description,
which justifies deleting redundant inequalities before computing it.
The Gaussian surface integral \(\surf(Q)=\int_{\partial Q}\Hh\,\varphi_r\,
d\Haus^{r-1}\) of Subsection~\PN{sec:outline-late}{2.4} is the case
\(n=r\) and \(f=\Hh\varphi_r\), with \(\Hh(g)=\norm g^2-r\).

\usedin{the definition of \(\surf\) in Subsection~\PN{sec:outline-late}{2.4}
and in \PR{Lemma}{appx-an:lem:surface}{C.5}.}
\emph{Source.} \TBc{EG15}, Chapter~2.

\begin{fact}[Gauss--Green formula (divergence theorem) for a polyhedron]\label{tb:divergence}
Let \(Q\subseteq\R^n\) (\(n\ge1\)) be a proper, full-dimensional polyhedron
with an irredundant description, facets \(F_1,\ldots,F_m\), and outward unit
normals \(u_1,\ldots,u_m\). Let \(X:\R^n\to\R^n\) be of class \(C^1\), with
\(\norm{X(x)}+|\operatorname{div}X(x)|\le C(1+\norm x)^k\varphi_n(x)\) for
some \(C\) and \(k\). Then
\[
 \int_Q\operatorname{div}X\,dx=\sum_{j=1}^m\int_{F_j}X\cdot u_j\,d\Haus^{n-1}.
\]
For \(n=1\) the facets are the endpoints of an interval, \(\Haus^0\) is
counting measure, and the identity is the fundamental theorem of calculus.
\end{fact}

\begin{proof}
It suffices to show, for a unit vector \(e\) and a function
\(f\in C^1(\R^n)\) with \(|f|+|\partial_ef|\le C(1+\norm x)^k\varphi_n(x)\),
that
\begin{equation}
 \int_Q\partial_ef\,dx=\sum_j(e\cdot u_j)\int_{F_j}f\,d\Haus^{n-1},
 \tag{\(*\)}
\end{equation}
since the theorem follows by taking \(e=e_i\) and \(f=X_i\) and summing over
\(i\). Write \(x=w+te\) with \(w\in e^\perp\). By Fubini's theorem,
\(\int_Q\partial_ef\,dx=\int_{e^\perp}\int_{I(w)}\partial_tf(w+te)\,dt\,dw\),
where \(I(w)=\{t:w+te\in Q\}\) is a closed interval, possibly empty, a point,
or unbounded, because \(Q\) is convex. Since \(f\) tends to zero along every
line, the inner integral equals \(f(w+t_+(w)e)-f(w+t_-(w)e)\), where
\(t_\pm(w)\) are the endpoints of \(I(w)\) and a term with an infinite
endpoint is read as zero.

Consider the upper endpoints \(x=w+t_+(w)e\). If \(x\in\relint F_j\), then
only the inequality \(j\) is active at \(x\) (Fact~\ref{tb:normalcone}), so
\(x+te\in Q\) for small \(t>0\) unless \(e\cdot u_j>0\), and \(x\) is an upper
endpoint if and only if \(e\cdot u_j>0\). The upper endpoints that lie in no
\(\relint F_j\) lie in \(\bigcup_{j\ne k}F_j\cap F_k\), which is a subset of a
finite union of affine subspaces of dimension at most \(n-2\)
(Fact~\ref{tb:facets}(b)); its projection along \(e\) onto \(e^\perp\) is a
Lebesgue null set of \(e^\perp\). For fixed \(j\) with \(e\cdot u_j>0\), the
map \(x\mapsto x-(x\cdot e)e\) is injective on the set of upper endpoints in
\(\relint F_j\), since a line meets \(Q\) in one interval, and it is the
restriction to \(H_j\) of a linear map whose Jacobian with respect to
\(\Haus^{n-1}\) on \(H_j\) is \(e\cdot u_j\): if \(b_1,\ldots,b_{n-1}\) is an
orthonormal basis of \(u_j^\perp\), the volume of the projection along \(e\)
of the cube they span equals the volume of the parallelepiped spanned by
\(b_1,\ldots,b_{n-1},e\), which is \(|e\cdot u_j|\). Therefore
\(\int f(w+t_+(w)e)\,dw=\sum_{j:\,e\cdot u_j>0}(e\cdot u_j)\int_{F_j}f\,d\Haus^{n-1}\),
the sets \(F_j\setminus\relint F_j\) being \(\Haus^{n-1}\)-null. The same
argument for the lower endpoints gives
\(-\sum_{j:\,e\cdot u_j<0}|e\cdot u_j|\int_{F_j}f\,d\Haus^{n-1}\), and
subtracting proves \((*)\); facets with \(e\cdot u_j=0\) contribute nothing
on either side.
\end{proof}

\usedin{\PR{Lemma}{appx-an:lem:surface}{C.5}, Step~1, applied in the hyperplane
\(H_j\) to the facet \(F_j\), regarded as a polyhedron of dimension
\(n=r-1\), whose facets are the ridges \(E_{jk}\) with outward normals
\(\nu_{k\mid j}\) (Fact~\ref{tb:facets}(c)), and to the vector field
\(X(y)=y\varphi_{r-1}(y)\), for which
\(\operatorname{div}X(y)=\{(r-1)-\norm y^2\}\varphi_{r-1}(y)\).
The decay of \(X\) is what replaces the truncation of unbounded facets.}
\emph{Source.} \TBc{GZ}, Lemma~24, proved there by truncation to a ball and
the Gauss--Green formula for Lipschitz domains~\TBc{EG15}, Chapters~4
and~5.

\subsection{Tubes around a polyhedron}\label{tb:sec:tubes}

Let \(Q\) be a proper, full-dimensional polyhedron in \(\R^n\) with an
irredundant description, facets \(F_j\), and ridges \(E_{jk}\) (\(n\ge2\)). For
a ridge, \(L_{jk}=\operatorname{span}(u_j,u_k)\), and \(z_{jk}\) is the unique
point of \(L_{jk}\) with \(u_j\cdot z_{jk}=a_j\) and \(u_k\cdot z_{jk}=a_k\). Every
point \(x\) of \(E_{jk}\) satisfies these two equations, so
\(E_{jk}\subseteq z_{jk}+L_{jk}^\perp\), and \(z_{jk}\) is the orthogonal
projection of every point of the ridge onto \(L_{jk}\). For \(\varepsilon>0\),
the \emph{outward tubes} of length \(\varepsilon\) are
\begin{align*}
 \mathcal T_j(\varepsilon)&=\{x+tu_j:x\in\relint F_j,\ 0<t<\varepsilon\},\\
 \mathcal T_{jk}(\varepsilon)&=\{x+v:x\in\relint E_{jk},\ v\in N_Q(x),\
 0<\norm v<\varepsilon\}.
\end{align*}
By Fact~\ref{tb:normalcone}, \(v\) ranges over the sector
\(S_{jk}(\varepsilon)=\{c_ju_j+c_ku_k:c_j,c_k\ge0\}\cap
\{\norm v<\varepsilon\}\) of angle \(\alpha_{jk}\) and area
\(\alpha_{jk}\varepsilon^2/2\) in \(L_{jk}\), without its apex. The tubes over the closed facets and ridges that
the proofs describe differ from these by sets of Lebesgue measure zero.
For \(Q(a)=\bigcap_j\{u_j\cdot x\le a_j\}\), write
\(Q(a)^t=Q(a_1+t,\ldots,a_m+t)\), the polyhedron with every offset raised by \(t\).

\begin{fact}[Tubes: disjointness and Gaussian mass]\label{tb:tubes}
(a) The sets \(\mathcal T_j(\varepsilon)\), \(j\le m\), and
\(\mathcal T_{jk}(\varepsilon)\), for the ridges \(E_{jk}\), are pairwise
disjoint, and they are contained in
\(\{y:0<\dist(y,Q)<\varepsilon\}\subseteq Q^\varepsilon\setminus Q\), where
\(Q^\varepsilon=Q(a)^\varepsilon\).

(b) With \(s_j=\int_{F_j}\varphi_n\,d\Haus^{n-1}\) and
\(\rho_{jk}=\int_{E_{jk}}\varphi_n\,d\Haus^{n-2}\),
\[
 \Phi_n(\mathcal T_j(\varepsilon))=\int_{F_j}\int_0^\varepsilon
 \varphi_n(x+tu_j)\,dt\,d\Haus^{n-1}(x),\qquad
 \Phi_n(\mathcal T_{jk}(\varepsilon))=\int_{E_{jk}}\int_{S_{jk}(\varepsilon)}
 \varphi_n(x+v)\,dv\,d\Haus^{n-2}(x).
\]

(c) Let \(R>0\) and \(\varepsilon=1/(1+R)\). If \(|a_j|\le R\), then
\(\Phi_n(\mathcal T_j(\varepsilon))\ge e^{-3/2}\,\varepsilon\,s_j\). If
\(\norm{z_{jk}}\le R\), then
\(\Phi_n(\mathcal T_{jk}(\varepsilon))\ge e^{-3/2}\,(\alpha_{jk}\varepsilon^2/2)\,
\rho_{jk}\).
\end{fact}

\begin{proof}
(a) By Fact~\ref{tb:projection}, a point \(y\notin Q\) is \(x+v\) for exactly
one pair with \(x=\prj_Q(y)\in Q\) and \(v\in N_Q(x)\setminus\{0\}\), and
\(\norm v=\dist(y,Q)\). A point of \(\mathcal T_j(\varepsilon)\) has its nearest
point in \(\relint F_j\), and a point of \(\mathcal T_{jk}(\varepsilon)\) has
its nearest point in \(\relint E_{jk}\), by Fact~\ref{tb:normalcone}. The
sets \(\relint F_j\) and \(\relint E_{jk}\) are pairwise disjoint, being
relative interiors of distinct faces, so the tubes are disjoint. If
\(\dist(y,Q)<\varepsilon\) and \(x=\prj_Q(y)\), then
\(u_j\cdot y\le u_j\cdot x+\norm{y-x}<a_j+\varepsilon\) for all \(j\), since the
\(u_j\) are unit vectors; thus \(y\in Q^\varepsilon\).

(b) The map \((x,t)\mapsto x+tu_j\) from \(H_j\times\R\) to \(\R^n\) is an
isometry, and \((x,v)\mapsto x+v\) is an isometry from
\((z_{jk}+L_{jk}^\perp)\times L_{jk}\) to \(\R^n\). Lebesgue measure on
\(\R^n\) is therefore the product of \(\Haus^{n-1}\) on \(H_j\) and \(dt\), and
of \(\Haus^{n-2}\) on \(z_{jk}+L_{jk}^\perp\) and \(dv\) on \(L_{jk}\)
(Fact~\ref{tb:hausdorff}(a)). Apply Fact~\ref{tb:fubini}.

(c) For \(x\in F_j\), \(x\cdot u_j=a_j\), and by Fact~\ref{tb:gauss-coord}(b)
\(\varphi_n(x+tu_j)=\varphi_n(x)\exp\{-ta_j-t^2/2\}\ge e^{-R\varepsilon-
\varepsilon^2/2}\varphi_n(x)\ge e^{-3/2}\varphi_n(x)\), since
\(R\varepsilon<1\) and \(\varepsilon<1\). Integrating over
\(t\in(0,\varepsilon)\) gives the first bound. For \(x\in E_{jk}\) and
\(v\in L_{jk}\), \(x\cdot v=z_{jk}\cdot v\), because \(x-z_{jk}\perp L_{jk}\),
so \(\varphi_n(x+v)=\varphi_n(x)\exp\{-z_{jk}\cdot v-\norm v^2/2\}
\ge e^{-3/2}\varphi_n(x)\) for \(\norm v<\varepsilon\). Integrating over the
sector, whose area is \(\alpha_{jk}\varepsilon^2/2\), gives the second bound.
\end{proof}

In the plane the tubes over the sides are strips and the tubes over the
corners are circular sectors, and together they tile the outer
\(\varepsilon\)-neighborhood of a polygon. When \(n\ge3\) the faces of
dimension at most \(n-3\) carry further strata of the neighborhood, which the
proofs do not need.

\usedin{\PR{Lemma}{appx-an:lem:surface}{C.5}, Step~2, for the bounds
\(q_j\ge c_*\varepsilon_Rs_j\) and \(q_{jk}\ge c_*\alpha_{jk}\varepsilon_R^2\rho_{jk}\),
and Step~3, for the disjointness that allows the Gaussian masses of the
tubes in a dyadic shell to be added, and for the inclusion
\(\mathcal O_\iota(V_0)\subseteq(Q')^{\varepsilon}\setminus Q'\) from which the
event \(V_0+\sigma G\in Q^{\sigma\varepsilon}\setminus Q\) is read: the
polyhedron \(Q'=(Q-V_0)/\sigma\) has offsets \((a_j-u_j\cdot V_0)/\sigma\),
and raising them by \(\varepsilon\) is raising those of \(Q\) by
\(\sigma\varepsilon\).}
\emph{Source.} \TBc{GZ}, Lemmas~28 and~29.

\subsection{The coarea formula and the first variation of a Gaussian
integral}\label{tb:sec:coarea}

\begin{fact}[Coarea formula]\label{tb:coarea}
Let \(\psi:\R^n\to\R\) be Lipschitz, so that \(\nabla\psi\) exists almost
everywhere. For every measurable \(f\ge0\),
\[
 \int_{\R^n}f\,\norm{\nabla\psi}\,dx=\int_{-\infty}^\infty
 \Bigl(\int_{\{\psi=s\}}f\,d\Haus^{n-1}\Bigr)ds ,
\]
and the same holds for integrable \(f\).
\end{fact}

\emph{Source.} \TBc{EG15}, Section~3.4. The proofs use the formula only for
a maximum of finitely many affine functions with gradients of equal length, and
for these Fact~\ref{tb:variation} gives a direct proof from Fubini's theorem.

\begin{fact}[First variation of a polyhedral sublevel integral]\label{tb:variation}
Let \(u_1,\ldots,u_m\) be unit vectors in \(\R^n\) (\(m\ge1\)), \(\kappa>0\),
\(b_1,\ldots,b_m\in\R\), and
\[
 \psi(g)=\max_{j\le m}\{\kappa\,u_j\cdot g-b_j\},\qquad Q_s=\{g:\psi(g)\le s\}.
\]
(a) For \(f\ge0\) measurable, or \(f\) integrable, and \(s\in\R\),
\[
 \int_{Q_s}f\,dg=\kappa^{-1}\int_{-\infty}^s\Bigl(\int_{\{\psi=s'\}}f\,d\Haus^{n-1}
 \Bigr)ds'.
\]
(b) If \(f\) is integrable, then \(\Lambda(s)=\int_{Q_s}f\,dg\) is absolutely
continuous, and for almost every \(s\)
\[
 \Lambda'(s)=\kappa^{-1}\int_{\{\psi=s\}}f\,d\Haus^{n-1},\qquad
 |\Lambda'(s)|\le\kappa^{-1}\sum_{j=1}^m\int_{\{\kappa u_j\cdot g-b_j=s\}}|f|\,d\Haus^{n-1}.
\]
(c) If \(s>\inf\psi\), then \(Q_s\) is a proper, full-dimensional polyhedron
with normals \(u_j\) and offsets \((s+b_j)/\kappa\), \(\{\psi=s\}=\partial Q_s\),
and \(\int_{\{\psi=s\}}f\,d\Haus^{n-1}\) is the surface integral of \(f\) over
\(\partial Q_s\) of Subsection~\ref{tb:sec:surface}. For \(s<\inf\psi\) the set
\(Q_s\) is empty, and \(s=\inf\psi\) is the only other level at which \(Q_s\)
may fail to be full-dimensional.

(d) If \(\Lambda\) is Lipschitz and \(S\) is a real random variable with a
density, then \(s\mapsto\E\Lambda(s+S)\) is differentiable with derivative
\(\E\Lambda'(s+S)\), where \(\Lambda'\) is defined almost everywhere.
\end{fact}

\begin{proof}
(a) Let \(C_j\) be the set where \(j\) is the smallest index attaining the
maximum defining \(\psi\); the Borel sets \(C_1,\ldots,C_m\) partition
\(\R^n\). On \(C_j\), \(\psi(g)=\kappa t-b_j\) in the coordinates
\(g=tu_j+y\), \(y\perp u_j\). By Fact~\ref{tb:fubini}, in these coordinates,
\[
 \int_{C_j\cap Q_s}f\,dg=\int_{u_j^\perp}\int_{-\infty}^{(s+b_j)/\kappa}
 (f\ind_{C_j})(tu_j+y)\,dt\,dy .
\]
The substitution \(s'=\kappa t-b_j\) turns the inner integral into
\(\kappa^{-1}\int_{-\infty}^s(f\ind_{C_j})(tu_j+y)\,ds'\) with
\(t=(s'+b_j)/\kappa\), and a further application of Fubini's theorem
exchanges the integrals. For fixed \(s'\), the set of \(tu_j+y\) in \(C_j\)
with this \(t\) is the slice \(C_j\cap\{\psi=s'\}\), and on the hyperplane
containing it \(\Haus^{n-1}\) is Lebesgue measure \(dy\)
(Fact~\ref{tb:hausdorff}(a)). Hence
\(\int_{C_j\cap Q_s}f\,dg=\kappa^{-1}\int_{-\infty}^s\int_{C_j\cap\{\psi=s'\}}
f\,d\Haus^{n-1}\,ds'\). Summing over \(j\) gives (a), since the slices
\(C_j\cap\{\psi=s'\}\) partition \(\{\psi=s'\}\).

(b) Applying (a) to \(|f|\) and letting \(s\to\infty\) shows that
\(I(s')=\kappa^{-1}\int_{\{\psi=s'\}}f\,d\Haus^{n-1}\) is integrable on
\(\R\), and (a) says \(\Lambda(s)=\int_{-\infty}^sI(s')\,ds'\). Hence \(\Lambda\) is
absolutely continuous, and \(\Lambda'=I\) almost everywhere by Lebesgue's
differentiation theorem. The bound follows because
\(\{\psi=s\}\subseteq\bigcup_j\{\kappa u_j\cdot g-b_j=s\}\), a union of \(m\)
hyperplanes.

(c) The set \(Q_s\) is the polyhedron \(\bigcap_j\{u_j\cdot g\le(s+b_j)/\kappa\}\),
and \(\psi\) is convex and continuous. If \(\psi(g)=s>\inf\psi\), choose \(g'\)
with \(\psi(g')<s\); for \(\tau>0\), \(g=(g_\tau+\tau g')/(1+\tau)\) with
\(g_\tau=g+\tau(g-g')\), so convexity gives
\(\psi(g_\tau)\ge s+\tau(s-\psi(g'))>s\). Thus \(g\) is a limit of points
outside \(Q_s\) and \(g\in\partial Q_s\); the open set \(\{\psi<s\}\) is
nonempty and lies in \(Q_s\), which is therefore full-dimensional, and it
is a proper subset because \(\psi(g)\ge\kappa u_1\cdot g-b_1\to\infty\) along
\(u_1\). A point with \(\psi<s\) is interior, so \(\partial Q_s=\{\psi=s\}\). After
the redundant and repeated inequalities are deleted, \(\partial Q_s=\{\psi=s\}\)
is the union of the sets \(Q_s\cap\{\kappa u_j\cdot g-b_j=s\}\) of the
remaining inequalities, which are faces of \(Q_s\) or empty. A face of
dimension \(n-1\) is a facet (Fact~\ref{tb:facets}(a)) and any other face is
\(\Haus^{n-1}\)-null, so the integral over \(\{\psi=s\}\) is the sum over the
facets of the irredundant description.

(d) A Lipschitz function is absolutely continuous and differentiable outside a
Lebesgue null set \(\mathcal N\). The random variable \(s+S\) avoids
\(\mathcal N\) almost surely, because \(S\) has a density, so the difference
quotients of \(\Lambda(s+S)\) converge almost surely and are bounded by the
Lipschitz constant. Dominated convergence gives the claim.
\end{proof}

\usedin{the paragraph ``First variation and common shifts'' after
\PR{Lemma}{appx-an:lem:surface}{C.5}, with
\(\psi(g)=\max_j\{u_j\cdot(V_0+\sigma g)-a_j\}\), \(\kappa=\sigma\),
\(b_j=a_j-u_j\cdot V_0\), \(Q_s=(Q(a+s\mathbf 1)-V_0)/\sigma\), and
\(f=\Hh\varphi_r\), which gives the derivative
\(\sigma^{-1}\surf((Q(a+s\mathbf 1)-V_0)/\sigma)\) for almost every \(s\) (parts
(a)--(c)); the Lipschitz bound in (b) is the bound \(Cm(r+1)/\sigma\) of that
paragraph, and part (d) justifies averaging the derivative over the
absolutely continuous common shift \(-h\log T_t\) in Step~4 of the proof of
\PR{Lemma}{appx-an:lem:matrix}{C.4}.}

\subsection{Where each fact is used}\label{tb:sec:index}

The table lists, for each fact of this appendix, the places in the proofs
that use it.

\medskip
\begin{center}
\small
\renewcommand{\arraystretch}{1.15}
\begin{tabular}{@{}p{0.27\textwidth}p{0.68\textwidth}@{}}
Facts~\ref{tb:fubini}, \ref{tb:taylor} &
every conditioning on an independent component in Appendix~\PN{appx-sec:support}{C};
the Taylor expansions in \PR{Lemma}{appx-an:lem:identity}{C.1}; Step~3 of
\PR{Lemma}{appx-an:lem:surface}{C.5}\\
Facts~\ref{tb:faa}, \ref{tb:logistic} &
the derivative formula and the local envelope in Subsection~\PN{appx-sec:mixture}{B.2};
Step~1 of \PR{Lemma}{appx-an:lem:restoration}{C.2}\\
Facts~\ref{tb:gauss-coord}--\ref{tb:heat} &
Subsection~\PN{sec:outline-path}{2.2} and \PR{Lemma}{appx-an:lem:identity}{C.1} (heat equation);
Step~1 of \PR{Lemma}{appx-an:lem:restoration}{C.2} (integration by parts);
\PR{Lemma}{appx-lem:endpoint-gap}{B.3}; Steps~1--2 of \PR{Lemma}{appx-an:lem:surface}{C.5}\\
Facts~\ref{tb:nuclear}, \ref{tb:logsum} &
Subsection~\PN{sec:outline-late}{2.4}; Step~4 of \PR{Lemma}{appx-an:lem:restoration}{C.2};
Steps~1 and~3 of \PR{Lemma}{appx-an:lem:matrix}{C.4}\\
Facts~\ref{tb:prekopa}--\ref{tb:lc-weights} &
Subsection~\PN{sec:outline-late}{2.4}; Step~1 of \PR{Lemma}{appx-an:lem:matrix}{C.4}\\
Facts~\ref{tb:facets}--\ref{tb:projection} &
the geometric paragraph and Steps~1--2 of \PR{Lemma}{appx-an:lem:surface}{C.5}\\
Facts~\ref{tb:hausdorff}, \ref{tb:divergence} &
the definition of \(\surf\); Step~1 of \PR{Lemma}{appx-an:lem:surface}{C.5}\\
Fact~\ref{tb:tubes} &
Steps~2 and~3 of \PR{Lemma}{appx-an:lem:surface}{C.5}\\
Facts~\ref{tb:coarea}, \ref{tb:variation} &
the paragraph on first variation after \PR{Lemma}{appx-an:lem:surface}{C.5}; Step~4 of
\PR{Lemma}{appx-an:lem:matrix}{C.4}\\
\end{tabular}
\end{center}

\subsection{Sources}\label{tb:sec:sources}

\begin{description}\small\setlength{\itemsep}{0pt}
\item[{[BL76]}] H.~J. Brascamp and E.~H. Lieb (1976). On extensions of the
Brunn--Minkowski and Pr\'ekopa--Leindler theorems, including inequalities
for log concave functions, and with an application to the diffusion
equation. \emph{Journal of Functional Analysis} \textbf{22}, 366--389.
\item[{[Com74]}] L. Comtet (1974). \emph{Advanced Combinatorics}. Reidel,
Dordrecht.
\item[{[CT06]}] T.~M. Cover and J.~A. Thomas (2006). \emph{Elements of
Information Theory}, 2nd edition. Wiley.
\item[{[EG15]}] L.~C. Evans and R.~F. Gariepy (2015). \emph{Measure Theory
and Fine Properties of Functions}, revised edition. CRC Press.
\item[{[Fol99]}] G.~B. Folland (1999). \emph{Real Analysis: Modern
Techniques and Their Applications}, 2nd edition. Wiley.
\item[{[Gar02]}] R.~J. Gardner (2002). The Brunn--Minkowski inequality.
\emph{Bulletin of the American Mathematical Society} \textbf{39}, 355--405.
\item[{[GZ]}] Z. Gao and W. Zhang (2026). Cubic-root Gaussian approximation
under unrestricted covariance. arXiv:2608.30221v1.
\item[{[Lei72]}] L. Leindler (1972). On a certain converse of H\"older's
inequality II. \emph{Acta Scientiarum Mathematicarum (Szeged)} \textbf{33},
217--223.
\item[{[Pre71]}] A. Pr\'ekopa (1971). Logarithmic concave measures with
application to stochastic programming. \emph{Acta Scientiarum Mathematicarum
(Szeged)} \textbf{32}, 301--316.
\item[{[Pre73]}] A. Pr\'ekopa (1973). On logarithmic concave measures and
functions. \emph{Acta Scientiarum Mathematicarum (Szeged)} \textbf{34},
335--343.
\item[{[Roc70]}] R.~T. Rockafellar (1970). \emph{Convex Analysis}. Princeton
University Press.
\item[{[RW98]}] R.~T. Rockafellar and R.~J.-B. Wets (1998). \emph{Variational
Analysis}. Springer, Berlin.
\item[{[Ste81]}] C.~M. Stein (1981). Estimation of the mean of a multivariate
normal distribution. \emph{The Annals of Statistics} \textbf{9}, 1135--1151.
\item[{[Zie95]}] G.~M. Ziegler (1995). \emph{Lectures on Polytopes}.
Springer, New York.
\end{description}

%
% after Appendix E and before \end{document}. The section then becomes
% Appendix F, and every reference to the paper below becomes a \ref to the
% paper's own labels. In the stand-alone version the same references print
% the numbers that the labels have in the paper on main (commit 12b6efd).
% The file defines the environment `fact` if the host does not, and defines
% its macros with \providecommand, so it leaves the host's definitions as
% they are. The list of sources at the end can be merged into the paper's
% bibliography.

\makeatletter
\providecommand{\E}{\mathbb E}
\providecommand{\Pp}{\mathbb P}
\providecommand{\R}{\mathbb R}
\providecommand{\tr}{\operatorname{tr}}
\providecommand{\op}{\mathrm{op}}
\providecommand{\norm}[1]{\left\lVert#1\right\rVert}
\providecommand{\ind}{\mathbf 1}
\providecommand{\dd}{\,d}
\providecommand{\Hh}{H_2}
\providecommand{\surf}{\mathcal J}
\providecommand{\relint}{\operatorname{relint}}
\providecommand{\aff}{\operatorname{aff}}
\providecommand{\cone}{\operatorname{cone}}
\providecommand{\dist}{\operatorname{dist}}
\providecommand{\prj}{\operatorname{pr}}
\providecommand{\Haus}{\mathcal H}
\providecommand{\TBc}[1]{[#1]}
% \PN{label}{C.5} prints a number of the paper (a \ref in the paper);
% \PR{Lemma}{label}{C.5} prints "Lemma C.5"; \PE{label}{C.21} an equation.
\providecommand{\PN}[2]{\ifdefined\toolboxinpaper\ref{#1}\else#2\fi}
\providecommand{\PR}[3]{#1~\PN{#2}{#3}}
\providecommand{\PE}[2]{(\PN{#1}{#2})}
\providecommand{\usedin}[1]{\par\smallskip\noindent{\small\emph{Used in:} #1}\par\smallskip}
\@ifundefined{c@fact}{\newtheorem{fact}{Fact}[section]}{}
\makeatother

\section{Toolbox: facts from analysis and geometry}\label{app:toolbox}

This appendix collects the facts from analysis, convex geometry, and
measure theory on which the proofs in Appendices~\PN{appx-sec:smoothing}{B}
and~\PN{appx-sec:support}{C} rely and which a mathematical statistician may
not have at hand. Every surface in those proofs is a finite union of flat
pieces of hyperplanes, so no differential geometry beyond the flat case is
needed. On a hyperplane the area formula is Fubini's theorem, the
Gauss--Green formula (the divergence theorem) is the fundamental theorem of
calculus along parallel lines, and the coarea formula is Fubini's theorem
applied piece by piece; the proofs below carry out each of these
reductions. Each fact is stated in the form in which
the proofs use it, and is followed by the places that use it and by a proof
or a reference. A proof is given when it takes a few lines, and a standard
source, listed at the end of the appendix, is cited otherwise. The facts are
standard results and not claims of the paper, so they carry their own
numbering. The symbols in the lines headed ``Used in'' are those of the
paper's appendices.

Subsection~\ref{tb:sec:calc} treats integration and calculus,
Subsection~\ref{tb:sec:gauss} the Gaussian identities, and
Subsection~\ref{tb:sec:matrix} the matrix and entropy inequalities.
Subsection~\ref{tb:sec:logconcave} covers log-concavity and
Pr\'ekopa's theorem. Subsections~\ref{tb:sec:poly}--\ref{tb:sec:coarea}
treat polyhedra, from faces and normal cones to the divergence theorem,
tubes, and the coarea formula, in the order in which
\PR{Lemma}{appx-an:lem:surface}{C.5} needs them, and
Subsection~\ref{tb:sec:index} lists where each fact is used.

\subsection{Notation}\label{tb:sec:notation}

The space \(\R^n\) carries the inner product \(x\cdot y\) and the Euclidean
norm \(\norm x\); in the proofs \(n\) is the rank \(r\) of the covariance,
or \(r-1\) when a hyperplane of \(\R^r\) is identified with
\(\R^{r-1}\). We write \(dx\) for Lebesgue measure, \(\varphi_n\) and
\(\Phi_n\) for the standard Gaussian density and law on \(\R^n\), and \(G\)
for a standard Gaussian vector. For a linear subspace \(L\), \(\varphi_L\)
is the standard Gaussian density on \(L\) with respect to Lebesgue measure
on \(L\). For a set \(S\), \(\ind_S\) is its indicator function, \(\aff S\)
is the smallest affine subspace containing it, the dimension of \(S\) is
that of \(\aff S\), and \(\relint S\) is the interior of \(S\) relative to
\(\aff S\). The distance from a point \(y\) to a set \(S\) is
\(\dist(y,S)=\inf_{x\in S}\norm{y-x}\).

For \(f\in C^q(\R^n)\), the derivative of order \(q\) in the directions
\(z_1,\ldots,z_q\in\R^n\) is
\[
 D^qf(x)[z_1,\ldots,z_q]
 =\partial_{s_1}\cdots\partial_{s_q}f(x+s_1z_1+\cdots+s_qz_q)\Big|_{s=0}
 =\sum_{i_1,\ldots,i_q}\partial_{i_1\cdots i_q}f(x)\,
   z_{1,i_1}\cdots z_{q,i_q},
\]
which is symmetric and multilinear in the directions. We write
\(D^qf(x)[z^{\otimes q}]\) when all directions equal \(z\), and
\(D^1f=\nabla f\). The Loewner order \(M\preceq N\) means that
\(N-M\) is positive semidefinite, \(\norm\Xi_\op\) is the operator norm
of a matrix \(\Xi\), and \(\norm M_*\) is the \emph{nuclear norm}, which for
a symmetric matrix \(M\) with eigenvalues \(\mu_i\) equals \(\sum_i|\mu_i|\).

\subsection{Integration and calculus}\label{tb:sec:calc}

\begin{fact}[Fubini and Tonelli]\label{tb:fubini}
Let \((S_1,\mu_1)\) and \((S_2,\mu_2)\) be \(\sigma\)-finite measure
spaces and let \(f\) be a measurable function on \(S_1\times S_2\). If
\(f\ge0\), then \(s_1\mapsto\int f(s_1,s_2)\,\mu_2(ds_2)\) is measurable
and
\[
 \int\Bigl(\int f\,d\mu_2\Bigr)d\mu_1
 =\int\Bigl(\int f\,d\mu_1\Bigr)d\mu_2
 =\int f\,d(\mu_1\otimes\mu_2)\in[0,\infty].
\]
If \(\int|f|\,d(\mu_1\otimes\mu_2)<\infty\), the same equalities hold, and
the inner integrals are finite almost everywhere.
\end{fact}

Three special cases occur in the proofs. If \(X\) and \(Y\) are
independent random elements and either \(f\ge0\) or \(f(X,Y)\) is
integrable, then
\[
 \E f(X,Y)=\E g(X),\qquad g(x)=\E f(x,Y),
\]
because the joint law is the product of the marginal laws. This is the
formal content of every argument that is carried out for a fixed value of
an independent component such as \(V_0\) and then averaged over it. With
\(\mu_2\) the counting measure, the theorem says that
\(\E\sum_k\xi_k=\sum_k\E\xi_k\) for nonnegative terms, and for terms with
\(\sum_k\E|\xi_k|<\infty\). For an orthogonal decomposition
\(\R^n=L\oplus L^\perp\), Lebesgue measure on \(\R^n\) is the product of the
Lebesgue measures on \(L\) and on \(L^\perp\), so that
\(\int_{\R^n}f\,dx=\int_{L^\perp}\int_Lf(l+w)\,dl\,dw\).

The appeal to Tonelli's theorem in Step~3 of the proof of
\PR{Lemma}{appx-an:lem:surface}{C.5}, whose notation we use in this
paragraph, is of the first kind. With \(V_0\) and
\(G\) independent, \(\Phi_r(\mathcal O_\iota(V_0))\) is the nonnegative
integral \(\int\ind\{g\in\mathcal O_\iota(V_0)\}\varphi_r(g)\,dg\), so
\(\E\sum_\iota R_\iota^3\,\Phi_r(\mathcal O_\iota(V_0))\) equals
\(\int\sum_\iota R_\iota^3\,\Pp\{g\in\mathcal O_\iota(V_0)\}\varphi_r(g)\,dg\)
with no check of integrability, since every term is nonnegative, and the
probability can then be bounded for each fixed \(g\).

\usedin{conditioning on \(V_0\) or on the \(\sigma\)-field \(\mathcal H\) in
\PR{Lemmas}{appx-an:lem:restoration}{C.2}, \PN{appx-an:lem:matrix}{C.4}
and~\PN{appx-an:lem:surface}{C.5}; conditioning on the exponentials in
\PR{Lemmas}{appx-lem:envelope}{B.1} and~\PN{appx-lem:endpoint-free}{B.2}
and on the candidate pairs in \PR{Appendix}{appx-app:concentration}{D};
interchanges of sums and expectations at the end of the proof of
\PR{Lemma}{appx-an:lem:restoration}{C.2}; Step~3 of the proof of
\PR{Lemma}{appx-an:lem:surface}{C.5}; iterated integrals over a hyperplane
and its normal in Subsections~\ref{tb:sec:surface}--\ref{tb:sec:coarea}.}
\emph{Source.} \TBc{Fol99}, Section~2.5.

\begin{fact}[Differentiation under the integral sign; Taylor's formula]
\label{tb:taylor}
(a) Let \(I\) be an open interval and let \(f(\lambda,\omega)\) be
integrable in \(\omega\) and differentiable in \(\lambda\in I\), with
\(|\partial_\lambda f(\lambda,\omega)|\le g(\omega)\) for an integrable
function \(g\). Then \(\frac d{d\lambda}\E f(\lambda,\cdot)
=\E\,\partial_\lambda f(\lambda,\cdot)\) on \(I\).

(b) For \(f\in C^q(\R^n)\) and \(x,v\in\R^n\),
\[
 f(x+v)=\sum_{k=0}^{q-1}\frac1{k!}D^kf(x)[v^{\otimes k}]
 +\int_0^1\frac{(1-s)^{q-1}}{(q-1)!}D^qf(x+sv)[v^{\otimes q}]\,ds.
\]
For \(q=1\) this is \(f(x+v)-f(x)=\int_0^1Df(x+sv)[v]\,ds\).
\end{fact}

\begin{proof}
For (a), the mean value theorem bounds the difference quotients by \(g\),
and they converge pointwise, so dominated convergence applies. For (b),
the function \(s\mapsto f(x+sv)\) is of class \(C^q\) with \(k\)th derivative
\(D^kf(x+sv)[v^{\otimes k}]\); apply the one-dimensional Taylor formula with
integral remainder at \(s=0\) and evaluate it at \(s=1\).
\end{proof}

\usedin{\PR{Lemma}{appx-an:lem:identity}{C.1}, where the derivative of
\(P_A(\lambda)=\E F(U_\lambda)\) is taken under the expectation (the
summands are finitely many and bounded) and Steps~1 and~3 use (b), to degree
three with the fourth-order integral remainder, to degree two for the
derivative, and to first order for the difference of third derivatives;
\PR{Lemma}{appx-lem:endpoint-gap}{B.3}, where (b) with \(q=1\) writes
\(\E[W\cdot\nabla\Psi(W)]\) as an integral of the Hessian along the segment
from \(W_i\) to \(W\).}
\emph{Source.} \TBc{Fol99}, Section~2.3.

\begin{fact}[Fa\`a di Bruno's formula]\label{tb:faa}
Let \(\phi\) be a smooth function on an open interval of \(\R\) and let
\(g:\R^n\to\R\) be smooth with values in that interval. For \(q\ge1\), let
\(\Pi_q\) be the set of partitions \(\mathcal P\) of \(\{1,\ldots,q\}\) into
nonempty blocks, let \(|\mathcal P|\) be the number of blocks, and for a
block \(J\) let \(z_J\) be the list of directions \(z_a\) with
\(a\in J\). Then
\[
 D^q(\phi\circ g)(x)[z_1,\ldots,z_q]
 =\sum_{\mathcal P\in\Pi_q}\phi^{(|\mathcal P|)}(g(x))
  \prod_{J\in\mathcal P}D^{|J|}g(x)[z_J].
\]
\end{fact}

\begin{proof}
For \(q=1\) this is the chain rule. Differentiate the right side for \(q\)
in the direction \(z_{q+1}\). By the product rule, the derivative either
falls on \(\phi^{(|\mathcal P|)}(g(x))\), which gives
\(\phi^{(|\mathcal P|+1)}(g(x))\,Dg(x)[z_{q+1}]\) and adds the block
\(\{q+1\}\) to \(\mathcal P\), or falls on one factor
\(D^{|J|}g(x)[z_J]\), which adds \(q+1\) to the block \(J\). Every partition
of \(\{1,\ldots,q+1\}\) arises exactly once in this way, by removing
\(q+1\) from its block.
\end{proof}

\begin{fact}[Derivatives of the logistic approximation]\label{tb:logistic}
Let \(\nu_1,\ldots,\nu_m\in\R^n\), \(a_1,\ldots,a_m\in\R\), \(h>0\), and
\[
 w_\alpha(x)=e^{(\nu_\alpha\cdot x-a_\alpha)/h},\qquad
 T=\sum_\alpha w_\alpha,\qquad
 F=\frac1{1+T},\qquad
 \pi_\alpha=\frac{w_\alpha}{1+T}=Fw_\alpha .
\]
For directions \(z_1,\ldots,z_q\), with the blocks \(J_1,\ldots,J_{|\mathcal P|}\)
of \(\mathcal P\) listed in any fixed order,
\[
 D^qF(x)[z_1,\ldots,z_q]
 =h^{-q}\sum_{\mathcal P\in\Pi_q}(-1)^{|\mathcal P|}|\mathcal P|!
 \sum_{\alpha_1,\ldots,\alpha_{|\mathcal P|}}
 F(x)\prod_{b=1}^{|\mathcal P|}\pi_{\alpha_b}(x)
 \prod_{b=1}^{|\mathcal P|}\prod_{a\in J_b}(\nu_{\alpha_b}\cdot z_a),
\]
and \(\sum_{\mathcal P\in\Pi_q}|\mathcal P|!\le q!\,2^{q-1}\). Consequently
\(|D^qF(x)[z_1,\ldots,z_q]|\le q!\,2^{q-1}h^{-q}\,F(x)\prod_{a=1}^q
\max_\alpha|\nu_\alpha\cdot z_a|\).
\end{fact}

\begin{proof}
Apply Fact~\ref{tb:faa} with \(\phi(\tau)=(1+\tau)^{-1}\) and \(g=T\). Then
\(\phi^{(k)}(\tau)=(-1)^kk!\,(1+\tau)^{-k-1}\), and since each \(w_\alpha\) is
an exponential of an affine function,
\(D^lT(x)[z_1,\ldots,z_l]=h^{-l}\sum_\alpha w_\alpha(x)\prod_{a\le l}
(\nu_\alpha\cdot z_a)\). Substituting, and using
\((1+T)^{-|\mathcal P|-1}\prod_bw_{\alpha_b}=F\prod_b\pi_{\alpha_b}\) and
\(\sum_b|J_b|=q\), gives the formula. The number
\(\sum_{\mathcal P}|\mathcal P|!\) counts ordered partitions
\((J_1,\ldots,J_k)\) of \(\{1,\ldots,q\}\). Writing the elements of
\(J_1\), then of \(J_2\), and so on, in increasing order inside each block,
gives a permutation together with the set of positions at which a new block
starts, which is a subset of the \(q-1\) gaps between consecutive entries.
Every ordered partition arises from at least one of these
\(q!\,2^{q-1}\) pairs. For the bound, the sum over \(\alpha_b\) of
\(\pi_{\alpha_b}\) is \(1-F\le1\), and each product of projections is
at most \(\prod_a\max_\alpha|\nu_\alpha\cdot z_a|\).
\end{proof}

\usedin{Subsection~\PN{appx-sec:mixture}{B.2}, where the first display
and the bound are \PE{appx-eq:partition-derivative}{B.1} and
\PE{appx-eq:all-orders}{B.2} (the text there cites the softmax cumulant formula,
\TBc{GZ}, Lemma~13); the envelope~\PE{appx-eq:local-envelope}{B.3} is
obtained by summing the same expansion over coordinate tuples;
\PR{Lemma}{appx-an:lem:restoration}{C.2}, Step~1, with \(u_\alpha\) in place
of \(\nu_\alpha\).}
\emph{Source.} \TBc{Com74}; the proof above is complete.

\subsection{Gaussian identities}\label{tb:sec:gauss}

\begin{fact}[Gaussian density in orthogonal coordinates]\label{tb:gauss-coord}
(a) Let \(\R^n=L\oplus L^\perp\) be an orthogonal decomposition. Then
\(\varphi_n(l+w)=\varphi_L(l)\varphi_{L^\perp}(w)\) for \(l\in L\) and
\(w\in L^\perp\), so the projections of a standard Gaussian vector onto
\(L\) and \(L^\perp\) are independent standard Gaussian vectors. In
particular, for a unit vector \(u\), \(a\in\R\), and \(y\perp u\),
\(\varphi_n(au+y)=\varphi_1(a)\varphi_{n-1}(y)\).

(b) For \(x,v\in\R^n\),
\(\varphi_n(x+v)=\varphi_n(x)\exp\{-x\cdot v-\tfrac12\norm v^2\}\).
\end{fact}

\begin{proof}
Both follow from \(\norm{l+w}^2=\norm l^2+\norm w^2\) and from expanding
\(\norm{x+v}^2\).
\end{proof}

\usedin{\PR{Lemma}{appx-an:lem:mean}{C.3}, where the normal and tangential
components of \(G\) relative to a face are independent;
\PR{Lemma}{appx-an:lem:surface}{C.5}, Steps~1 and~2, for the factorization of
\(\varphi_r\) on a facet and on a ridge and for the density ratio on a tube.}

\begin{fact}[Gaussian integration by parts]\label{tb:stein}
Let \(G\sim N(0,I_n)\) and let \(f\in C^2(\R^n)\) have \(f\) and its partial
derivatives of orders at most two of polynomial growth. Then
\begin{enumerate}
\item[(a)] \(\E[Gf(G)]=\E\nabla f(G)\);
\item[(b)] \(\E[(GG^\top-I_n)f(G)]=\E D^2f(G)\); equivalently, for \(z\in\R^n\),
\(\E D^2f(G)[z,z]=\E\bigl[\tr\{(GG^\top-I_n)zz^\top\}f(G)\bigr]\), where
\(\tr\{(GG^\top-I_n)zz^\top\}=(z\cdot G)^2-\norm z^2\);
\item[(c)] if \(v\in\R^n\), \(\sigma>0\), \(U=v+\sigma G\), and
\(F\in C^{s+2}(\R^n)\) has derivatives of polynomial growth, then for
\(z,z_1,\ldots,z_s\in\R^n\),
\[
 \E D^{s+2}F(U)[z,z,z_1,\ldots,z_s]
 =\sigma^{-2}\E\bigl[\tr\{(GG^\top-I_n)zz^\top\}\,
 D^sF(U)[z_1,\ldots,z_s]\bigr].
\]
\end{enumerate}
\end{fact}

\begin{proof}
For (a) in one coordinate, \(x\varphi_1(x)=-\varphi_1'(x)\), so integration by
parts gives
\[\int xf(x)\varphi_1(x)\,dx=\int f'(x)\varphi_1(x)\,dx,\]
and the boundary terms vanish because of the growth condition. The other coordinates
are handled by Fact~\ref{tb:fubini}. For (b), apply (a) to \(\partial_if\) and
to \(G_kf\): \(\E[G_iG_kf(G)]=\E\partial_i(G_kf)(G)=\delta_{ik}\E f(G)
+\E[G_k\partial_if(G)]=\delta_{ik}\E f(G)+\E\partial_{ik}f(G)\). For (c), the
function \(f(g)=D^sF(v+\sigma g)[z_1,\ldots,z_s]\) has
\(D^2f(g)[z,z]=\sigma^2D^{s+2}F(v+\sigma g)[z,z,z_1,\ldots,z_s]\); apply (b).
\end{proof}

\usedin{Subsection~\PN{sec:outline-late}{2.4} (the display after
``integration by parts against the covariance of the summands''),
\PR{Lemma}{appx-an:lem:restoration}{C.2}, Step~1, where (c) is applied
conditionally on \(\mathcal H\) with \(v=V_0+v_i\) and \(z=Z_i\);
\PR{Lemma}{appx-lem:endpoint-gap}{B.3}, for \(\E[Z\cdot\nabla\Psi(Z)]
=\E\tr(\Sigma D^2\Psi(Z))\), which is (a) with \(Z=\Gamma G\) for a factor
\(\Gamma\) of the covariance.}
\emph{Source.} Stein~\TBc{Ste81}; the proof above is complete.

\begin{fact}[Heat equation for the Gaussian component]\label{tb:heat}
Let \(Z\sim N(0,\Sigma)\) in \(\R^p\), where \(\Sigma\) may be singular, and
let \(f\in C^2(\R^p)\) be bounded with bounded first and second derivatives.
Then, for \(x\in\R^p\) and \(\lambda>0\),
\[
 \frac d{d\lambda}\E f(x+\sqrt\lambda Z)
 =\tfrac12\E\tr\{\Sigma D^2f(x+\sqrt\lambda Z)\}.
\]
\end{fact}

\begin{proof}
Write \(Z=\Gamma G\) with \(\Gamma\Gamma^\top=\Sigma\) and \(G\) standard
Gaussian. By Fact~\ref{tb:taylor}(a), which applies because the difference
quotients are bounded by \(\norm\nabla f_\infty\norm{\Gamma G}/(2\sqrt\lambda)\)
locally uniformly in \(\lambda>0\), the derivative equals
\((2\sqrt\lambda)^{-1}\E[G\cdot\Gamma^\top\nabla f(x+\sqrt\lambda\Gamma G)]\).
Fact~\ref{tb:stein}(a) applied to the function \(g\mapsto
\Gamma^\top\nabla f(x+\sqrt\lambda\Gamma g)\) turns this into
\((2\sqrt\lambda)^{-1}\sqrt\lambda\,\E\tr\{\Gamma^\top D^2f\,\Gamma\}
=\frac12\E\tr\{\Sigma D^2f\}\).
\end{proof}

\usedin{Subsection~\PN{sec:outline-path}{2.2} and the proof of
\PR{Lemma}{appx-an:lem:identity}{C.1}, Step~2, where the Gaussian part of the
path contributes \(\frac12\sum_iM_i:\E D^2F(U_\lambda)\).}

\subsection{Matrix and entropy inequalities}\label{tb:sec:matrix}

\begin{fact}[Duality between the operator and nuclear norms]\label{tb:nuclear}
(a) For every square matrix \(\Xi\) and every symmetric matrix \(M\),
\(|\tr(\Xi M)|\le\norm\Xi_\op\norm M_*\). Moreover
\(\norm M_*=\max\{\tr(\Xi M):\norm\Xi_\op\le1\}\), so that \(\norm\cdot_*\) is
a norm on symmetric matrices.

(b) If \(M=P-N\) with \(P\succeq0\) and \(N\succeq0\), then
\(\norm M_*\le\tr P+\tr N=2\tr P-\tr M\).
\end{fact}

\begin{proof}
For (a), let \(M=\sum_i\mu_ie_ie_i^\top\) be a spectral decomposition. Then
\(\tr(\Xi M)=\sum_i\mu_i\,e_i^\top\Xi e_i\) and \(|e_i^\top\Xi e_i|\le
\norm\Xi_\op\), which gives the inequality. Equality is attained at
\(\Xi=\sum_i\operatorname{sgn}(\mu_i)e_ie_i^\top\). A supremum of linear
functionals satisfies the triangle inequality. For (b), the norm of a
positive semidefinite matrix is its trace, so the triangle inequality gives
\(\norm M_*\le\tr P+\tr N\), and \(\tr N=\tr P-\tr M\).
\end{proof}

\usedin{Subsection~\PN{sec:outline-late}{2.4}; \PR{Lemma}{appx-an:lem:restoration}{C.2},
Step~4; \PR{Lemma}{appx-an:lem:matrix}{C.4}, Step~1, where
(b) is applied with \(P=p_\alpha\mu_\alpha\mu_\alpha^\top\) and
\(N=p_\alpha(I_r-\Sigma_\alpha)\), and gives
\PE{appx-an:matrix-trace-reduction}{C.17}.}

\begin{fact}[Log-sum inequality and a concave function]\label{tb:logsum}
(a) For positive numbers \(y_1,\ldots,y_l\) and \(z_1,\ldots,z_l\) with sums
\(y\) and \(z\),
\(\sum_jy_j\log(z_j/y_j)\le y\log(z/y)\).

(b) The function \((x,y)\mapsto x\log(y/x)\) is concave on \((0,\infty)^2\).
Hence, for positive integrable random variables \(X,Y\),
\(\E[X\log(Y/X)]\le\E X\,\log(\E Y/\E X)\).
\end{fact}

\begin{proof}
For (a), Jensen's inequality for the concave function \(\log\), with the
weights \(y_j/y\), gives \(\sum_j(y_j/y)\log(z_j/y_j)\le\log\sum_j(y_j/y)(z_j/y_j)
=\log(z/y)\). For (b), the function \(f(x,y)=x\log y-x\log x\) has Hessian
\(\begin{pmatrix}-1/x&1/y\\1/y&-x/y^2\end{pmatrix}\), whose determinant is
zero and whose trace is negative, so the Hessian is negative semidefinite;
the second statement is Jensen's inequality for the concave function
\(f\) and the random vector \((X,Y)\).
\end{proof}

\usedin{\PR{Lemma}{appx-an:lem:matrix}{C.4}, Step~3: (a) adds the reference
masses of the tuples that share a first face, and (b) averages over
\(V_0\) and over the faces.}
\emph{Source.} \TBc{CT06}, Theorem~2.7.1, for (a).

\subsection{Log-concavity}\label{tb:sec:logconcave}

A function \(f:\R^n\to[0,\infty)\) is \emph{log-concave} if
\(f((1-t)x+ty)\ge f(x)^{1-t}f(y)^t\) for all \(x,y\) and \(t\in[0,1]\);
when \(f>0\) this says that \(\log f\) is concave.

\begin{fact}[Pr\'ekopa--Leindler inequality and Pr\'ekopa's theorem]
\label{tb:prekopa}
(a) Let \(0<t<1\) and let \(f,g,h:\R^n\to[0,\infty)\) be bounded,
continuous, and integrable, with
\(h((1-t)x+ty)\ge f(x)^{1-t}g(y)^t\) for all \(x,y\). Then
\(\int h\ge(\int f)^{1-t}(\int g)^t\).

(b) Let \(f:\R^a\times\R^b\to[0,\infty)\) be bounded, continuous, and
log-concave, with \(f(z,\cdot)\) integrable for every \(z\). Then
\(z\mapsto\int f(z,y)\,dy\) is log-concave on \(\R^a\).
\end{fact}

\begin{proof}
(a) For \(n=1\), assume first that \(\sup f=\sup g=1\); the general case
follows by dividing \(f,g,h\) by \(\sup f\), \(\sup g\), and
\((\sup f)^{1-t}(\sup g)^t\). For \(0<s<1\) the level sets
\(\{f\ge s\}\) and \(\{g\ge s\}\) are nonempty, and the hypothesis gives
\(\{h\ge s\}\supseteq(1-t)\{f\ge s\}+t\{g\ge s\}\). In one dimension
\(|A+B|\ge|A|+|B|\) for nonempty compact sets \(A\) and \(B\), because
\(A+B\) contains the translates \(A+\min B\) and \(\max A+B\), which meet in
one point; the inequality passes to measurable sets by approximation from
inside by compact sets. Therefore
\(|\{h\ge s\}|\ge(1-t)|\{f\ge s\}|+t|\{g\ge s\}|\), and integration over
\(s\in(0,1)\) gives \(\int h\ge(1-t)\int f+t\int g\ge(\int f)^{1-t}(\int g)^t\),
by the inequality between the arithmetic and geometric means. For \(n\ge2\),
apply the case of dimension \(n-1\) to the sections of \(f,g,h\) at the last
coordinates \(s_0\), \(s_1\), and \((1-t)s_0+ts_1\), and then the case of
dimension one to the integrals of the sections, with Fubini's theorem
(Fact~\ref{tb:fubini}) connecting the steps; the routine measurability
checks are omitted here and can be found in~\TBc{Gar02}.

(b) Fix \(z_0,z_1\) and put \(z_t=(1-t)z_0+tz_1\). Log-concavity of \(f\)
gives \(f(z_t,(1-t)y_0+ty_1)\ge f(z_0,y_0)^{1-t}f(z_1,y_1)^t\), so (a) applies
to \(f(z_0,\cdot)\), \(f(z_1,\cdot)\), and \(f(z_t,\cdot)\).
\end{proof}

\emph{Source.} Pr\'ekopa~\TBc{Pre71,Pre73}, Leindler~\TBc{Lei72}; a
survey with the proof of the general case is~\TBc{Gar02}, Section~7. The
paper cites Theorem~6 of~\TBc{Pre73}.

\begin{fact}[Gaussian laws tilted by a log-concave weight]\label{tb:tilt}
Let \(w:\R^r\to[0,\infty)\) be bounded, continuous, and log-concave, with
\(p=\E w(G)>0\) for \(G\sim N(0,I_r)\). Let \(\mu\) and \(K\) be the mean and
covariance matrix of the probability law with density \(w(g)\varphi_r(g)/p\).
Then \(K\preceq I_r\), and
\(\E[(GG^\top-I_r)w(G)]=p(K-I_r+\mu\mu^\top)\).
\end{fact}

\begin{proof}
The function \((z,g)\mapsto e^{-\norm{g-z}^2/2}w(g)\) is log-concave, since
\(\norm{g-z}^2\) is jointly convex in \((z,g)\), and it is bounded and
continuous. By Fact~\ref{tb:prekopa}(b), \(\ell(z)=\log\int e^{-\norm{g-z}^2/2}
w(g)\,dg\) is concave. Since \(\norm{g-z}^2=\norm g^2-2z\cdot g+\norm z^2\),
\[
 \ell(z)=\text{const}-\tfrac12\norm z^2+\log M(z),\qquad
 M(z)=\int e^{z\cdot g}w(g)\varphi_r(g)\,dg=p\,\E_\nu e^{z\cdot g},
\]
where \(\nu\) is the tilted law. The function \(M\) is finite and smooth,
because \(w\) is bounded, and the Hessian of \(\log M\) at \(0\) is the
covariance \(K\) of \(\nu\). A concave function of class \(C^2\) has a
negative semidefinite Hessian, so \(K-I_r\preceq0\). The identity is the
decomposition of the second moment \(\E_\nu gg^\top=K+\mu\mu^\top\).
\end{proof}

\usedin{Subsection~\PN{sec:outline-late}{2.4} (``Pr\'ekopa's theorem \dots\
has covariance \(K\preceq I_r\)'') and \PR{Lemma}{appx-an:lem:matrix}{C.4},
Step~1, where \(w=w_\alpha\) and \(K=\Sigma_\alpha\).}
\emph{Source.} For indicators of convex cells, the cell covariance bound
\TBc{GZ}, Lemma~11, is the same estimate \(K\preceq I_r\), obtained from the
Brascamp--Lieb inequality~\TBc{BL76}; the weights here are not indicators,
which is why the proof goes through Pr\'ekopa's theorem.

\begin{fact}[Log-concavity of the logistic weights]\label{tb:lc-weights}
With the notation of Fact~\ref{tb:logistic}, for every \(k\ge0\) and every
tuple \((\alpha_1,\ldots,\alpha_k)\), the function
\(F\pi_{\alpha_1}\cdots\pi_{\alpha_k}=\exp\{\sum_{a\le k}t_{\alpha_a}\}
(1+T)^{-k-1}\), where \(t_\alpha=(\nu_\alpha\cdot x-a_\alpha)/h\), is
log-concave on \(\R^n\), and so is its composition
\(g\mapsto F\pi_{\alpha_1}\cdots\pi_{\alpha_k}(v+\sigma g)\) with any affine map.
\end{fact}

\begin{proof}
The logarithm is a sum of the affine functions \(t_{\alpha_a}\) and
\(-(k+1)\log(1+\sum_\alpha e^{t_\alpha})\). The function
\(\log(1+\sum_\alpha e^{t_\alpha})\) is the logarithm of a sum of
exponentials of the affine functions \(0,t_1,\ldots,t_m\) and is convex by
H\"older's inequality: \(\sum_\alpha e^{(1-s)t_\alpha+st'_\alpha}\le
(\sum_\alpha e^{t_\alpha})^{1-s}(\sum_\alpha e^{t'_\alpha})^s\). Composition
with an affine map preserves concavity.
\end{proof}

\usedin{Subsection~\PN{sec:outline-late}{2.4} and Step~1 of the proof of
\PR{Lemma}{appx-an:lem:matrix}{C.4}.}

\subsection{Polyhedra: faces, normal cones, and nearest points}\label{tb:sec:poly}

A \emph{polyhedron} is a set
\[
 Q=\bigcap_{j=1}^m\{x\in\R^n:u_j\cdot x\le a_j\},
\]
with unit vectors \(u_j\), called \emph{outward normals}, and numbers
\(a_j\), called \emph{offsets}. It is closed and convex. It is
\emph{full-dimensional} if it has an interior point, and \emph{proper} if it
is neither empty nor \(\R^n\). The description is \emph{irredundant} if
deleting any one inequality changes the set; deleting redundant or repeated
inequalities one at a time turns every description into an irredundant one
of the same set. For \(j\le m\), put \(H_j=\{x:u_j\cdot x=a_j\}\) and
\(F_j=Q\cap H_j\). A \emph{face} of \(Q\) is \(Q\) itself or a set
\(Q\cap\{u\cdot x=c\}\) with \(u\ne0\) and \(Q\subseteq\{u\cdot x\le c\}\).
A \emph{facet} is a face of dimension \(n-1\), and a \emph{ridge} is a face
of dimension \(n-2\). In the plane the facets are the sides of a polygon
and the ridges its corners, in \(\R^3\) the facets are the faces of a solid
and the ridges its edges, and on the line the facets are the endpoints of an
interval, which has no ridges. We use these standard properties of faces:
the relative interiors of the faces of \(Q\) are pairwise disjoint and cover
\(Q\); a face of a face of \(Q\) is a face of \(Q\); and a face \(G'\)
of a face \(G\) with \(G'\ne G\) lies in the relative boundary
\(G\setminus\relint G\) of \(G\) (\TBc{Roc70}, Section~18; \TBc{Zie95}, Chapter~2).

\begin{fact}[Facets and ridges of a polyhedron]\label{tb:facets}
Let \(Q\) be a proper, full-dimensional polyhedron in \(\R^n\) with an
irredundant description.
\begin{enumerate}
\item[(a)] The normals \(u_1,\ldots,u_m\) are pairwise distinct, each
\(F_j\) is a facet, \(Q\) has no other facets, and
\(\partial Q=F_1\cup\cdots\cup F_m\).
\item[(b)] For \(j\ne k\), \(F_j\cap F_k\subseteq H_j\cap H_k\), and
\(H_j\cap H_k\) is empty if \(u_k=-u_j\), and an affine subspace of
dimension \(n-2\) otherwise. In particular \(F_j\cap F_k\) lies in a
finite union of affine subspaces of dimension at most \(n-2\).
\item[(c)] Each \(F_j\) is a full-dimensional polyhedron in the hyperplane
\(H_j\). Write the points of \(H_j\) as \(a_ju_j+y\) with \(y\perp u_j\), and,
for \(k\ne j\) with \(u_k\ne\pm u_j\), put
\(\alpha_{jk}=\arccos(u_j\cdot u_k)\in(0,\pi)\),
\[
 \nu_{k\mid j}=\frac{u_k-(u_j\cdot u_k)u_j}{\sin\alpha_{jk}},\qquad
 d_{k\mid j}=\frac{a_k-(u_j\cdot u_k)a_j}{\sin\alpha_{jk}} .
\]
Then \(F_j\) is described in the coordinate \(y\) by the inequalities
\(\nu_{k\mid j}\cdot y\le d_{k\mid j}\), and its facets are exactly the
ridges \(E_{jk}=F_j\cap F_k\) of \(Q\) of dimension \(n-2\); the outward
unit normal of \(F_j\) at \(E_{jk}\) is \(\nu_{k\mid j}\).
\item[(d)] Every ridge of \(Q\) lies in exactly two facets, so the ridges
correspond to the unordered pairs \(\{j,k\}\) with
\(\dim(F_j\cap F_k)=n-2\); for such a pair \(u_k\ne\pm u_j\), and the
only inequalities active at a point of \(\relint E_{jk}\) are \(j\) and \(k\).
\end{enumerate}
\end{fact}

\begin{proof}[Proof (sketch)]
(a) Fix \(j\). Irredundancy gives a point \(x_j\) with \(u_j\cdot x_j>a_j\)
and \(u_k\cdot x_j\le a_k\) for \(k\ne j\). Let \(x_0\) be an interior point of
\(Q\), and let \(x_*\) be the point where the segment from \(x_0\) to \(x_j\)
crosses \(H_j\). For \(k\ne j\), \(u_k\cdot x_*<a_k\), because \(x_*\) is a
convex combination of \(x_0\), where the inequality is strict, with positive
weight, and of \(x_j\), where it holds weakly. Hence a neighborhood of
\(x_*\) in \(H_j\) lies in \(F_j\), and \(\dim F_j=n-1\). Two inequalities
with the same normal cannot both be irredundant, since the one with the
larger offset is implied by the other. A point of \(Q\) at
which no inequality is active is interior, and at a point where
inequality \(j\) is active the points \(x+\varepsilon u_j\) lie outside \(Q\),
so \(\partial Q=\bigcup_jF_j\). If \(G=Q\cap\{u\cdot x=c\}\) is a facet and
\(x\in\relint G\), some inequality \(j\) is active at \(x\); on \(G\) the
affine function \(u_j\cdot x-a_j\le0\) attains its maximum at a relative
interior point and is therefore zero, so \(G\subseteq H_j\), and since \(G\) has
dimension \(n-1\), \(G=Q\cap H_j=F_j\).

(b) If \(u_k=-u_j\) and \(H_j\cap H_k\ne\emptyset\), then \(Q\) would lie in
the hyperplane \(H_j\), contradicting full dimension. If \(u_k\ne\pm u_j\),
the two equations are independent.

(c) The set \(F_j\) is the polyhedron in \(H_j\) cut out by the remaining
inequalities, which take the stated form in the coordinate \(y\) because
\(u_k\cdot(a_ju_j+y)=a_ju_j\cdot u_k+(u_k-(u_j\cdot u_k)u_j)\cdot y\); an
inequality with \(u_k=-u_j\) holds automatically on \(H_j\) since
\(F_j\ne\emptyset\). By (a) applied inside \(H_j\), the facets of \(F_j\) are
the faces of \(F_j\) of dimension \(n-2\), hence faces of \(Q\) of that
dimension, that is, ridges of \(Q\) in \(F_j\); each equals
\(F_j\cap H_k=F_j\cap F_k\) for the inequality \(k\) that defines it.

(d) Let \(E\) be a ridge and \(y\in\relint E\), and let \(J\) be the set of
inequalities active at \(y\). On \(E\) each affine function
\(u_i\cdot x-a_i\), \(i\in J\), is at most zero and vanishes at the relative
interior point \(y\), so it vanishes on \(E\); thus \(E\subseteq F_i\) for
\(i\in J\) and \(\dim\operatorname{span}\{u_i:i\in J\}\le2\). A single active
inequality is impossible, since then \(y\) would lie in the relative interior
of \(F_j\) and the only face of \(F_j\) containing it is \(F_j\) itself, of
dimension \(n-1\). Opposite normals \(u_i=-u_j\) are excluded by (b). Hence
the active normals span a plane \(\Pi\). Near \(y\), \(Q\) coincides with
\(y\) plus the product of \(\Pi^\perp\) and a planar cone \(K\subseteq\Pi\)
with nonempty interior, defined by the active inequalities and contained in
a wedge of opening \(\pi-\alpha_{jk}<\pi\); this cone has exactly two
boundary rays. Each active inequality \(i\) has \(\dim F_i=n-1\), so
\(F_i\) meets every neighborhood of \(y\) in a set of dimension \(n-1\), and
the line \(\{u_i\cdot v=0\}\) must therefore carry a boundary ray of \(K\).
With distinct normals this leaves exactly two active inequalities.
\end{proof}

\usedin{\PR{Lemma}{appx-an:lem:surface}{C.5}, in the paragraph on the terms
from the geometry of polyhedra and in Step~1 (``each inequality defines
exactly one facet, \dots\ two facets that share a ridge have normals that
are neither equal nor opposite'') and Step~2 (``every ridge lies in exactly
two facets'').}
\emph{Source.} \TBc{Zie95}, Chapter~2; \TBc{Roc70}, Sections~18 and~19;
the facet decomposition of \(\partial Q\) with the pairing of ridges is
the content of \TBc{GZ}, Lemma~27.

\begin{fact}[Normal cones]\label{tb:normalcone}
Let \(Q\) be a polyhedron as in Fact~\ref{tb:facets}. For \(x\in Q\) let
\(J(x)=\{j:u_j\cdot x=a_j\}\), and let
\[
 N_Q(x)=\{v\in\R^n:v\cdot(q-x)\le0\text{ for all }q\in Q\}
\]
be the \emph{normal cone} of \(Q\) at \(x\). Then
\(N_Q(x)=\{\sum_{j\in J(x)}c_ju_j:c_j\ge0\}\). In particular
\(N_Q(x)=\{0\}\) at interior points, \(N_Q(x)=\{cu_j:c\ge0\}\) for
\(x\in\relint F_j\), and \(N_Q(x)=\{c_ju_j+c_ku_k:c_j,c_k\ge0\}\), a sector of
angle \(\alpha_{jk}\) in the plane spanned by \(u_j\) and \(u_k\), for
\(x\in\relint E_{jk}\).
\end{fact}

\begin{proof}
Put \(T=\{d:u_j\cdot d\le0\text{ for }j\in J(x)\}\). A vector \(v\) lies in
\(N_Q(x)\) if and only if \(v\cdot d\le0\) for all \(d\in T\): if \(d\in T\)
then \(x+\varepsilon d\in Q\) for small \(\varepsilon>0\), because the
inactive inequalities are strict at \(x\), and conversely \(q-x\in T\) for
\(q\in Q\). By Farkas's lemma, which says that \(v\cdot d\le0\) for all
\(d\) with \(u_j\cdot d\le0\) (\(j\in J\)) if and only if \(v\) is a
nonnegative combination of the \(u_j\), \(j\in J\), this is the stated
cone. For \(x\in\relint F_j\), \(J(x)=\{j\}\), since an index \(k\ne j\) in
\(J(x)\) would put \(x\) in \(F_j\cap F_k\), a face of \(F_j\) of dimension
at most \(n-2\), and so in the relative boundary of \(F_j\). For
\(x\in\relint E_{jk}\), \(J(x)=\{j,k\}\) by Fact~\ref{tb:facets}(d).
\end{proof}

\usedin{\PR{Lemma}{appx-an:lem:surface}{C.5}, Step~2 (``the normal cone at a
relative-interior point of the ridge shared by facets \(j\) and \(k\) is the
cone spanned by \(u_j\) and \(u_k\)'') and the definition of the outward
tubes.}
\emph{Source.} \TBc{Roc70}, Sections~22 and~23, for Farkas's lemma and normal
cones; \TBc{RW98}, Chapter~6.

\begin{fact}[Nearest point]\label{tb:projection}
Let \(Q\subseteq\R^n\) be nonempty, closed, and convex. For each
\(y\in\R^n\) there is a unique point \(\prj_Q(y)\in Q\) with
\(\norm{y-\prj_Q(y)}=\dist(y,Q)\), and \(x=\prj_Q(y)\) if and only if
\(x\in Q\) and \(y-x\in N_Q(x)\). Consequently every \(y\in\R^n\) can be
written in exactly one way as \(y=x+v\) with \(x\in Q\) and \(v\in N_Q(x)\),
and then \(\norm v=\dist(y,Q)\). A direction \(v\) lies in \(N_Q(x)\) if and
only if \(x\) is the nearest point of \(x+tv\) in \(Q\) for every \(t>0\).
\end{fact}

\begin{proof}
For \(x,q\in Q\), \(\norm{y-q}^2=\norm{y-x}^2+\norm{x-q}^2-2(y-x)\cdot(q-x)\).
If \(y-x\in N_Q(x)\), the last term is nonnegative, so
\(\norm{y-q}^2\ge\norm{y-x}^2+\norm{x-q}^2>\norm{y-x}^2\) for \(q\ne x\); thus
\(x\) is the unique nearest point. Conversely, if \((y-x)\cdot(q-x)>0\) for
some \(q\in Q\), then moving from \(x\) toward \(q\) decreases the distance
to \(y\) for small steps, by the displayed identity applied to
\(x+\varepsilon(q-x)\in Q\), so \(x\) is not nearest. A nearest point exists
because \(Q\) is closed. The last statement follows because \(N_Q(x)\) is a
cone.
\end{proof}

\usedin{\PR{Lemma}{appx-an:lem:surface}{C.5}, Step~2, to show that the tubes
are disjoint (``their points have distinct metric projections'').}
\emph{Source.} \TBc{RW98}, Chapter~6 (Moreau's decomposition); \TBc{GZ},
proof of Lemma~28.

\subsection{Surface measure, Gaussian surface integrals, and the
divergence theorem}\label{tb:sec:surface}

\begin{fact}[Hausdorff measure on flat pieces]\label{tb:hausdorff}
For \(0\le k\le n\), the \(k\)-dimensional Hausdorff measure \(\Haus^k\) on
\(\R^n\) is the Borel measure
\[
 \Haus^k(S)=\lim_{\epsilon\downarrow0}\inf\Bigl\{\sum_i\omega_k
 \Bigl(\tfrac12\operatorname{diam}C_i\Bigr)^k:S\subseteq\bigcup_iC_i,\
 \operatorname{diam}C_i\le\epsilon\Bigr\},
\]
where \(\omega_k\) is the volume of the unit ball of \(\R^k\). (a) On an
affine subspace of dimension \(k\), \(\Haus^k\) is \(k\)-dimensional Lebesgue
measure; in particular \(\Haus^0\) is counting measure and \(\Haus^n\) is
Lebesgue measure on \(\R^n\). (b) \(\Haus^k\) is invariant under isometries.
(c) A set contained in a finite union of affine subspaces of dimension less
than \(k\) has \(\Haus^k\)-measure zero.
\end{fact}

Part (c) follows from (a), since each such subspace lies in a
\(k\)-dimensional affine subspace in which it is a Lebesgue null set. For a
proper, full-dimensional polyhedron \(Q\) with an irredundant description,
and for a function \(f\ge0\) or \(f\) integrable on \(\partial Q\), the
\emph{surface integral} is
\[
 \int_{\partial Q}f\,d\Haus^{n-1}=\sum_{j=1}^m\int_{F_j}f\,d\Haus^{n-1},
\]
where on each facet \(\Haus^{n-1}\) is Lebesgue measure on the hyperplane
\(H_j\) by part~(a). By Fact~\ref{tb:facets}(b) and part~(c), distinct
facets overlap in sets of \(\Haus^{n-1}\)-measure zero, so the sum is the
integral of \(f\) over the set \(\partial Q\) with respect to
\(\Haus^{n-1}\). It depends on the set \(Q\) and not on its description,
which justifies deleting redundant inequalities before computing it.
The Gaussian surface integral \(\surf(Q)=\int_{\partial Q}\Hh\,\varphi_r\,
d\Haus^{r-1}\) of Subsection~\PN{sec:outline-late}{2.4} is the case
\(n=r\) and \(f=\Hh\varphi_r\), with \(\Hh(g)=\norm g^2-r\).

\usedin{the definition of \(\surf\) in Subsection~\PN{sec:outline-late}{2.4}
and in \PR{Lemma}{appx-an:lem:surface}{C.5}.}
\emph{Source.} \TBc{EG15}, Chapter~2.

\begin{fact}[Gauss--Green formula (divergence theorem) for a polyhedron]\label{tb:divergence}
Let \(Q\subseteq\R^n\) (\(n\ge1\)) be a proper, full-dimensional polyhedron
with an irredundant description, facets \(F_1,\ldots,F_m\), and outward unit
normals \(u_1,\ldots,u_m\). Let \(X:\R^n\to\R^n\) be of class \(C^1\), with
\(\norm{X(x)}+|\operatorname{div}X(x)|\le C(1+\norm x)^k\varphi_n(x)\) for
some \(C\) and \(k\). Then
\[
 \int_Q\operatorname{div}X\,dx=\sum_{j=1}^m\int_{F_j}X\cdot u_j\,d\Haus^{n-1}.
\]
For \(n=1\) the facets are the endpoints of an interval, \(\Haus^0\) is
counting measure, and the identity is the fundamental theorem of calculus.
\end{fact}

\begin{proof}
It suffices to show, for a unit vector \(e\) and a function
\(f\in C^1(\R^n)\) with \(|f|+|\partial_ef|\le C(1+\norm x)^k\varphi_n(x)\),
that
\begin{equation}
 \int_Q\partial_ef\,dx=\sum_j(e\cdot u_j)\int_{F_j}f\,d\Haus^{n-1},
 \tag{\(*\)}
\end{equation}
since the theorem follows by taking \(e=e_i\) and \(f=X_i\) and summing over
\(i\). Write \(x=w+te\) with \(w\in e^\perp\). By Fubini's theorem,
\(\int_Q\partial_ef\,dx=\int_{e^\perp}\int_{I(w)}\partial_tf(w+te)\,dt\,dw\),
where \(I(w)=\{t:w+te\in Q\}\) is a closed interval, possibly empty, a point,
or unbounded, because \(Q\) is convex. Since \(f\) tends to zero along every
line, the inner integral equals \(f(w+t_+(w)e)-f(w+t_-(w)e)\), where
\(t_\pm(w)\) are the endpoints of \(I(w)\) and a term with an infinite
endpoint is read as zero.

Consider the upper endpoints \(x=w+t_+(w)e\). If \(x\in\relint F_j\), then
only the inequality \(j\) is active at \(x\) (Fact~\ref{tb:normalcone}), so
\(x+te\in Q\) for small \(t>0\) unless \(e\cdot u_j>0\), and \(x\) is an upper
endpoint if and only if \(e\cdot u_j>0\). The upper endpoints that lie in no
\(\relint F_j\) lie in \(\bigcup_{j\ne k}F_j\cap F_k\), which is a subset of a
finite union of affine subspaces of dimension at most \(n-2\)
(Fact~\ref{tb:facets}(b)); its projection along \(e\) onto \(e^\perp\) is a
Lebesgue null set of \(e^\perp\). For fixed \(j\) with \(e\cdot u_j>0\), the
map \(x\mapsto x-(x\cdot e)e\) is injective on the set of upper endpoints in
\(\relint F_j\), since a line meets \(Q\) in one interval, and it is the
restriction to \(H_j\) of a linear map whose Jacobian with respect to
\(\Haus^{n-1}\) on \(H_j\) is \(e\cdot u_j\): if \(b_1,\ldots,b_{n-1}\) is an
orthonormal basis of \(u_j^\perp\), the volume of the projection along \(e\)
of the cube they span equals the volume of the parallelepiped spanned by
\(b_1,\ldots,b_{n-1},e\), which is \(|e\cdot u_j|\). Therefore
\(\int f(w+t_+(w)e)\,dw=\sum_{j:\,e\cdot u_j>0}(e\cdot u_j)\int_{F_j}f\,d\Haus^{n-1}\),
the sets \(F_j\setminus\relint F_j\) being \(\Haus^{n-1}\)-null. The same
argument for the lower endpoints gives
\(-\sum_{j:\,e\cdot u_j<0}|e\cdot u_j|\int_{F_j}f\,d\Haus^{n-1}\), and
subtracting proves \((*)\); facets with \(e\cdot u_j=0\) contribute nothing
on either side.
\end{proof}

\usedin{\PR{Lemma}{appx-an:lem:surface}{C.5}, Step~1, applied in the hyperplane
\(H_j\) to the facet \(F_j\), regarded as a polyhedron of dimension
\(n=r-1\), whose facets are the ridges \(E_{jk}\) with outward normals
\(\nu_{k\mid j}\) (Fact~\ref{tb:facets}(c)), and to the vector field
\(X(y)=y\varphi_{r-1}(y)\), for which
\(\operatorname{div}X(y)=\{(r-1)-\norm y^2\}\varphi_{r-1}(y)\).
The decay of \(X\) is what replaces the truncation of unbounded facets.}
\emph{Source.} \TBc{GZ}, Lemma~24, proved there by truncation to a ball and
the Gauss--Green formula for Lipschitz domains~\TBc{EG15}, Chapters~4
and~5.

\subsection{Tubes around a polyhedron}\label{tb:sec:tubes}

Let \(Q\) be a proper, full-dimensional polyhedron in \(\R^n\) with an
irredundant description, facets \(F_j\), and ridges \(E_{jk}\) (\(n\ge2\)). For
a ridge, \(L_{jk}=\operatorname{span}(u_j,u_k)\), and \(z_{jk}\) is the unique
point of \(L_{jk}\) with \(u_j\cdot z_{jk}=a_j\) and \(u_k\cdot z_{jk}=a_k\). Every
point \(x\) of \(E_{jk}\) satisfies these two equations, so
\(E_{jk}\subseteq z_{jk}+L_{jk}^\perp\), and \(z_{jk}\) is the orthogonal
projection of every point of the ridge onto \(L_{jk}\). For \(\varepsilon>0\),
the \emph{outward tubes} of length \(\varepsilon\) are
\begin{align*}
 \mathcal T_j(\varepsilon)&=\{x+tu_j:x\in\relint F_j,\ 0<t<\varepsilon\},\\
 \mathcal T_{jk}(\varepsilon)&=\{x+v:x\in\relint E_{jk},\ v\in N_Q(x),\
 0<\norm v<\varepsilon\}.
\end{align*}
By Fact~\ref{tb:normalcone}, \(v\) ranges over the sector
\(S_{jk}(\varepsilon)=\{c_ju_j+c_ku_k:c_j,c_k\ge0\}\cap
\{\norm v<\varepsilon\}\) of angle \(\alpha_{jk}\) and area
\(\alpha_{jk}\varepsilon^2/2\) in \(L_{jk}\), without its apex. The tubes over the closed facets and ridges that
the proofs describe differ from these by sets of Lebesgue measure zero.
For \(Q(a)=\bigcap_j\{u_j\cdot x\le a_j\}\), write
\(Q(a)^t=Q(a_1+t,\ldots,a_m+t)\), the polyhedron with every offset raised by \(t\).

\begin{fact}[Tubes: disjointness and Gaussian mass]\label{tb:tubes}
(a) The sets \(\mathcal T_j(\varepsilon)\), \(j\le m\), and
\(\mathcal T_{jk}(\varepsilon)\), for the ridges \(E_{jk}\), are pairwise
disjoint, and they are contained in
\(\{y:0<\dist(y,Q)<\varepsilon\}\subseteq Q^\varepsilon\setminus Q\), where
\(Q^\varepsilon=Q(a)^\varepsilon\).

(b) With \(s_j=\int_{F_j}\varphi_n\,d\Haus^{n-1}\) and
\(\rho_{jk}=\int_{E_{jk}}\varphi_n\,d\Haus^{n-2}\),
\[
 \Phi_n(\mathcal T_j(\varepsilon))=\int_{F_j}\int_0^\varepsilon
 \varphi_n(x+tu_j)\,dt\,d\Haus^{n-1}(x),\qquad
 \Phi_n(\mathcal T_{jk}(\varepsilon))=\int_{E_{jk}}\int_{S_{jk}(\varepsilon)}
 \varphi_n(x+v)\,dv\,d\Haus^{n-2}(x).
\]

(c) Let \(R>0\) and \(\varepsilon=1/(1+R)\). If \(|a_j|\le R\), then
\(\Phi_n(\mathcal T_j(\varepsilon))\ge e^{-3/2}\,\varepsilon\,s_j\). If
\(\norm{z_{jk}}\le R\), then
\(\Phi_n(\mathcal T_{jk}(\varepsilon))\ge e^{-3/2}\,(\alpha_{jk}\varepsilon^2/2)\,
\rho_{jk}\).
\end{fact}

\begin{proof}
(a) By Fact~\ref{tb:projection}, a point \(y\notin Q\) is \(x+v\) for exactly
one pair with \(x=\prj_Q(y)\in Q\) and \(v\in N_Q(x)\setminus\{0\}\), and
\(\norm v=\dist(y,Q)\). A point of \(\mathcal T_j(\varepsilon)\) has its nearest
point in \(\relint F_j\), and a point of \(\mathcal T_{jk}(\varepsilon)\) has
its nearest point in \(\relint E_{jk}\), by Fact~\ref{tb:normalcone}. The
sets \(\relint F_j\) and \(\relint E_{jk}\) are pairwise disjoint, being
relative interiors of distinct faces, so the tubes are disjoint. If
\(\dist(y,Q)<\varepsilon\) and \(x=\prj_Q(y)\), then
\(u_j\cdot y\le u_j\cdot x+\norm{y-x}<a_j+\varepsilon\) for all \(j\), since the
\(u_j\) are unit vectors; thus \(y\in Q^\varepsilon\).

(b) The map \((x,t)\mapsto x+tu_j\) from \(H_j\times\R\) to \(\R^n\) is an
isometry, and \((x,v)\mapsto x+v\) is an isometry from
\((z_{jk}+L_{jk}^\perp)\times L_{jk}\) to \(\R^n\). Lebesgue measure on
\(\R^n\) is therefore the product of \(\Haus^{n-1}\) on \(H_j\) and \(dt\), and
of \(\Haus^{n-2}\) on \(z_{jk}+L_{jk}^\perp\) and \(dv\) on \(L_{jk}\)
(Fact~\ref{tb:hausdorff}(a)). Apply Fact~\ref{tb:fubini}.

(c) For \(x\in F_j\), \(x\cdot u_j=a_j\), and by Fact~\ref{tb:gauss-coord}(b)
\(\varphi_n(x+tu_j)=\varphi_n(x)\exp\{-ta_j-t^2/2\}\ge e^{-R\varepsilon-
\varepsilon^2/2}\varphi_n(x)\ge e^{-3/2}\varphi_n(x)\), since
\(R\varepsilon<1\) and \(\varepsilon<1\). Integrating over
\(t\in(0,\varepsilon)\) gives the first bound. For \(x\in E_{jk}\) and
\(v\in L_{jk}\), \(x\cdot v=z_{jk}\cdot v\), because \(x-z_{jk}\perp L_{jk}\),
so \(\varphi_n(x+v)=\varphi_n(x)\exp\{-z_{jk}\cdot v-\norm v^2/2\}
\ge e^{-3/2}\varphi_n(x)\) for \(\norm v<\varepsilon\). Integrating over the
sector, whose area is \(\alpha_{jk}\varepsilon^2/2\), gives the second bound.
\end{proof}

In the plane the tubes over the sides are strips and the tubes over the
corners are circular sectors, and together they tile the outer
\(\varepsilon\)-neighborhood of a polygon. When \(n\ge3\) the faces of
dimension at most \(n-3\) carry further strata of the neighborhood, which the
proofs do not need.

\usedin{\PR{Lemma}{appx-an:lem:surface}{C.5}, Step~2, for the bounds
\(q_j\ge c_*\varepsilon_Rs_j\) and \(q_{jk}\ge c_*\alpha_{jk}\varepsilon_R^2\rho_{jk}\),
and Step~3, for the disjointness that allows the Gaussian masses of the
tubes in a dyadic shell to be added, and for the inclusion
\(\mathcal O_\iota(V_0)\subseteq(Q')^{\varepsilon}\setminus Q'\) from which the
event \(V_0+\sigma G\in Q^{\sigma\varepsilon}\setminus Q\) is read: the
polyhedron \(Q'=(Q-V_0)/\sigma\) has offsets \((a_j-u_j\cdot V_0)/\sigma\),
and raising them by \(\varepsilon\) is raising those of \(Q\) by
\(\sigma\varepsilon\).}
\emph{Source.} \TBc{GZ}, Lemmas~28 and~29.

\subsection{The coarea formula and the first variation of a Gaussian
integral}\label{tb:sec:coarea}

\begin{fact}[Coarea formula]\label{tb:coarea}
Let \(\psi:\R^n\to\R\) be Lipschitz, so that \(\nabla\psi\) exists almost
everywhere. For every measurable \(f\ge0\),
\[
 \int_{\R^n}f\,\norm{\nabla\psi}\,dx=\int_{-\infty}^\infty
 \Bigl(\int_{\{\psi=s\}}f\,d\Haus^{n-1}\Bigr)ds ,
\]
and the same holds for integrable \(f\).
\end{fact}

\emph{Source.} \TBc{EG15}, Section~3.4. The proofs use the formula only for
a maximum of finitely many affine functions with gradients of equal length, and
for these Fact~\ref{tb:variation} gives a direct proof from Fubini's theorem.

\begin{fact}[First variation of a polyhedral sublevel integral]\label{tb:variation}
Let \(u_1,\ldots,u_m\) be unit vectors in \(\R^n\) (\(m\ge1\)), \(\kappa>0\),
\(b_1,\ldots,b_m\in\R\), and
\[
 \psi(g)=\max_{j\le m}\{\kappa\,u_j\cdot g-b_j\},\qquad Q_s=\{g:\psi(g)\le s\}.
\]
(a) For \(f\ge0\) measurable, or \(f\) integrable, and \(s\in\R\),
\[
 \int_{Q_s}f\,dg=\kappa^{-1}\int_{-\infty}^s\Bigl(\int_{\{\psi=s'\}}f\,d\Haus^{n-1}
 \Bigr)ds'.
\]
(b) If \(f\) is integrable, then \(\Lambda(s)=\int_{Q_s}f\,dg\) is absolutely
continuous, and for almost every \(s\)
\[
 \Lambda'(s)=\kappa^{-1}\int_{\{\psi=s\}}f\,d\Haus^{n-1},\qquad
 |\Lambda'(s)|\le\kappa^{-1}\sum_{j=1}^m\int_{\{\kappa u_j\cdot g-b_j=s\}}|f|\,d\Haus^{n-1}.
\]
(c) If \(s>\inf\psi\), then \(Q_s\) is a proper, full-dimensional polyhedron
with normals \(u_j\) and offsets \((s+b_j)/\kappa\), \(\{\psi=s\}=\partial Q_s\),
and \(\int_{\{\psi=s\}}f\,d\Haus^{n-1}\) is the surface integral of \(f\) over
\(\partial Q_s\) of Subsection~\ref{tb:sec:surface}. For \(s<\inf\psi\) the set
\(Q_s\) is empty, and \(s=\inf\psi\) is the only other level at which \(Q_s\)
may fail to be full-dimensional.

(d) If \(\Lambda\) is Lipschitz and \(S\) is a real random variable with a
density, then \(s\mapsto\E\Lambda(s+S)\) is differentiable with derivative
\(\E\Lambda'(s+S)\), where \(\Lambda'\) is defined almost everywhere.
\end{fact}

\begin{proof}
(a) Let \(C_j\) be the set where \(j\) is the smallest index attaining the
maximum defining \(\psi\); the Borel sets \(C_1,\ldots,C_m\) partition
\(\R^n\). On \(C_j\), \(\psi(g)=\kappa t-b_j\) in the coordinates
\(g=tu_j+y\), \(y\perp u_j\). By Fact~\ref{tb:fubini}, in these coordinates,
\[
 \int_{C_j\cap Q_s}f\,dg=\int_{u_j^\perp}\int_{-\infty}^{(s+b_j)/\kappa}
 (f\ind_{C_j})(tu_j+y)\,dt\,dy .
\]
The substitution \(s'=\kappa t-b_j\) turns the inner integral into
\(\kappa^{-1}\int_{-\infty}^s(f\ind_{C_j})(tu_j+y)\,ds'\) with
\(t=(s'+b_j)/\kappa\), and a further application of Fubini's theorem
exchanges the integrals. For fixed \(s'\), the set of \(tu_j+y\) in \(C_j\)
with this \(t\) is the slice \(C_j\cap\{\psi=s'\}\), and on the hyperplane
containing it \(\Haus^{n-1}\) is Lebesgue measure \(dy\)
(Fact~\ref{tb:hausdorff}(a)). Hence
\(\int_{C_j\cap Q_s}f\,dg=\kappa^{-1}\int_{-\infty}^s\int_{C_j\cap\{\psi=s'\}}
f\,d\Haus^{n-1}\,ds'\). Summing over \(j\) gives (a), since the slices
\(C_j\cap\{\psi=s'\}\) partition \(\{\psi=s'\}\).

(b) Applying (a) to \(|f|\) and letting \(s\to\infty\) shows that
\(I(s')=\kappa^{-1}\int_{\{\psi=s'\}}f\,d\Haus^{n-1}\) is integrable on
\(\R\), and (a) says \(\Lambda(s)=\int_{-\infty}^sI(s')\,ds'\). Hence \(\Lambda\) is
absolutely continuous, and \(\Lambda'=I\) almost everywhere by Lebesgue's
differentiation theorem. The bound follows because
\(\{\psi=s\}\subseteq\bigcup_j\{\kappa u_j\cdot g-b_j=s\}\), a union of \(m\)
hyperplanes.

(c) The set \(Q_s\) is the polyhedron \(\bigcap_j\{u_j\cdot g\le(s+b_j)/\kappa\}\),
and \(\psi\) is convex and continuous. If \(\psi(g)=s>\inf\psi\), choose \(g'\)
with \(\psi(g')<s\); for \(\tau>0\), \(g=(g_\tau+\tau g')/(1+\tau)\) with
\(g_\tau=g+\tau(g-g')\), so convexity gives
\(\psi(g_\tau)\ge s+\tau(s-\psi(g'))>s\). Thus \(g\) is a limit of points
outside \(Q_s\) and \(g\in\partial Q_s\); the open set \(\{\psi<s\}\) is
nonempty and lies in \(Q_s\), which is therefore full-dimensional, and it
is a proper subset because \(\psi(g)\ge\kappa u_1\cdot g-b_1\to\infty\) along
\(u_1\). A point with \(\psi<s\) is interior, so \(\partial Q_s=\{\psi=s\}\). After
the redundant and repeated inequalities are deleted, \(\partial Q_s=\{\psi=s\}\)
is the union of the sets \(Q_s\cap\{\kappa u_j\cdot g-b_j=s\}\) of the
remaining inequalities, which are faces of \(Q_s\) or empty. A face of
dimension \(n-1\) is a facet (Fact~\ref{tb:facets}(a)) and any other face is
\(\Haus^{n-1}\)-null, so the integral over \(\{\psi=s\}\) is the sum over the
facets of the irredundant description.

(d) A Lipschitz function is absolutely continuous and differentiable outside a
Lebesgue null set \(\mathcal N\). The random variable \(s+S\) avoids
\(\mathcal N\) almost surely, because \(S\) has a density, so the difference
quotients of \(\Lambda(s+S)\) converge almost surely and are bounded by the
Lipschitz constant. Dominated convergence gives the claim.
\end{proof}

\usedin{the paragraph ``First variation and common shifts'' after
\PR{Lemma}{appx-an:lem:surface}{C.5}, with
\(\psi(g)=\max_j\{u_j\cdot(V_0+\sigma g)-a_j\}\), \(\kappa=\sigma\),
\(b_j=a_j-u_j\cdot V_0\), \(Q_s=(Q(a+s\mathbf 1)-V_0)/\sigma\), and
\(f=\Hh\varphi_r\), which gives the derivative
\(\sigma^{-1}\surf((Q(a+s\mathbf 1)-V_0)/\sigma)\) for almost every \(s\) (parts
(a)--(c)); the Lipschitz bound in (b) is the bound \(Cm(r+1)/\sigma\) of that
paragraph, and part (d) justifies averaging the derivative over the
absolutely continuous common shift \(-h\log T_t\) in Step~4 of the proof of
\PR{Lemma}{appx-an:lem:matrix}{C.4}.}

\subsection{Where each fact is used}\label{tb:sec:index}

The table lists, for each fact of this appendix, the places in the proofs
that use it.

\medskip
\begin{center}
\small
\renewcommand{\arraystretch}{1.15}
\begin{tabular}{@{}p{0.27\textwidth}p{0.68\textwidth}@{}}
Facts~\ref{tb:fubini}, \ref{tb:taylor} &
every conditioning on an independent component in Appendix~\PN{appx-sec:support}{C};
the Taylor expansions in \PR{Lemma}{appx-an:lem:identity}{C.1}; Step~3 of
\PR{Lemma}{appx-an:lem:surface}{C.5}\\
Facts~\ref{tb:faa}, \ref{tb:logistic} &
the derivative formula and the local envelope in Subsection~\PN{appx-sec:mixture}{B.2};
Step~1 of \PR{Lemma}{appx-an:lem:restoration}{C.2}\\
Facts~\ref{tb:gauss-coord}--\ref{tb:heat} &
Subsection~\PN{sec:outline-path}{2.2} and \PR{Lemma}{appx-an:lem:identity}{C.1} (heat equation);
Step~1 of \PR{Lemma}{appx-an:lem:restoration}{C.2} (integration by parts);
\PR{Lemma}{appx-lem:endpoint-gap}{B.3}; Steps~1--2 of \PR{Lemma}{appx-an:lem:surface}{C.5}\\
Facts~\ref{tb:nuclear}, \ref{tb:logsum} &
Subsection~\PN{sec:outline-late}{2.4}; Step~4 of \PR{Lemma}{appx-an:lem:restoration}{C.2};
Steps~1 and~3 of \PR{Lemma}{appx-an:lem:matrix}{C.4}\\
Facts~\ref{tb:prekopa}--\ref{tb:lc-weights} &
Subsection~\PN{sec:outline-late}{2.4}; Step~1 of \PR{Lemma}{appx-an:lem:matrix}{C.4}\\
Facts~\ref{tb:facets}--\ref{tb:projection} &
the geometric paragraph and Steps~1--2 of \PR{Lemma}{appx-an:lem:surface}{C.5}\\
Facts~\ref{tb:hausdorff}, \ref{tb:divergence} &
the definition of \(\surf\); Step~1 of \PR{Lemma}{appx-an:lem:surface}{C.5}\\
Fact~\ref{tb:tubes} &
Steps~2 and~3 of \PR{Lemma}{appx-an:lem:surface}{C.5}\\
Facts~\ref{tb:coarea}, \ref{tb:variation} &
the paragraph on first variation after \PR{Lemma}{appx-an:lem:surface}{C.5}; Step~4 of
\PR{Lemma}{appx-an:lem:matrix}{C.4}\\
\end{tabular}
\end{center}

\subsection{Sources}\label{tb:sec:sources}

\begin{description}\small\setlength{\itemsep}{0pt}
\item[{[BL76]}] H.~J. Brascamp and E.~H. Lieb (1976). On extensions of the
Brunn--Minkowski and Pr\'ekopa--Leindler theorems, including inequalities
for log concave functions, and with an application to the diffusion
equation. \emph{Journal of Functional Analysis} \textbf{22}, 366--389.
\item[{[Com74]}] L. Comtet (1974). \emph{Advanced Combinatorics}. Reidel,
Dordrecht.
\item[{[CT06]}] T.~M. Cover and J.~A. Thomas (2006). \emph{Elements of
Information Theory}, 2nd edition. Wiley.
\item[{[EG15]}] L.~C. Evans and R.~F. Gariepy (2015). \emph{Measure Theory
and Fine Properties of Functions}, revised edition. CRC Press.
\item[{[Fol99]}] G.~B. Folland (1999). \emph{Real Analysis: Modern
Techniques and Their Applications}, 2nd edition. Wiley.
\item[{[Gar02]}] R.~J. Gardner (2002). The Brunn--Minkowski inequality.
\emph{Bulletin of the American Mathematical Society} \textbf{39}, 355--405.
\item[{[GZ]}] Z. Gao and W. Zhang (2026). Cubic-root Gaussian approximation
under unrestricted covariance. arXiv:2608.30221v1.
\item[{[Lei72]}] L. Leindler (1972). On a certain converse of H\"older's
inequality II. \emph{Acta Scientiarum Mathematicarum (Szeged)} \textbf{33},
217--223.
\item[{[Pre71]}] A. Pr\'ekopa (1971). Logarithmic concave measures with
application to stochastic programming. \emph{Acta Scientiarum Mathematicarum
(Szeged)} \textbf{32}, 301--316.
\item[{[Pre73]}] A. Pr\'ekopa (1973). On logarithmic concave measures and
functions. \emph{Acta Scientiarum Mathematicarum (Szeged)} \textbf{34},
335--343.
\item[{[Roc70]}] R.~T. Rockafellar (1970). \emph{Convex Analysis}. Princeton
University Press.
\item[{[RW98]}] R.~T. Rockafellar and R.~J.-B. Wets (1998). \emph{Variational
Analysis}. Springer, Berlin.
\item[{[Ste81]}] C.~M. Stein (1981). Estimation of the mean of a multivariate
normal distribution. \emph{The Annals of Statistics} \textbf{9}, 1135--1151.
\item[{[Zie95]}] G.~M. Ziegler (1995). \emph{Lectures on Polytopes}.
Springer, New York.
\end{description}

%
% after Appendix E and before \end{document}. The section then becomes
% Appendix F, and every reference to the paper below becomes a \ref to the
% paper's own labels. In the stand-alone version the same references print
% the numbers that the labels have in the paper on main (commit 12b6efd).
% The file defines the environment `fact` if the host does not, and defines
% its macros with \providecommand, so it leaves the host's definitions as
% they are. The list of sources at the end can be merged into the paper's
% bibliography.

\makeatletter
\providecommand{\E}{\mathbb E}
\providecommand{\Pp}{\mathbb P}
\providecommand{\R}{\mathbb R}
\providecommand{\tr}{\operatorname{tr}}
\providecommand{\op}{\mathrm{op}}
\providecommand{\norm}[1]{\left\lVert#1\right\rVert}
\providecommand{\ind}{\mathbf 1}
\providecommand{\dd}{\,d}
\providecommand{\Hh}{H_2}
\providecommand{\surf}{\mathcal J}
\providecommand{\relint}{\operatorname{relint}}
\providecommand{\aff}{\operatorname{aff}}
\providecommand{\cone}{\operatorname{cone}}
\providecommand{\dist}{\operatorname{dist}}
\providecommand{\prj}{\operatorname{pr}}
\providecommand{\Haus}{\mathcal H}
\providecommand{\TBc}[1]{[#1]}
% \PN{label}{C.5} prints a number of the paper (a \ref in the paper);
% \PR{Lemma}{label}{C.5} prints "Lemma C.5"; \PE{label}{C.21} an equation.
\providecommand{\PN}[2]{\ifdefined\toolboxinpaper\ref{#1}\else#2\fi}
\providecommand{\PR}[3]{#1~\PN{#2}{#3}}
\providecommand{\PE}[2]{(\PN{#1}{#2})}
\providecommand{\usedin}[1]{\par\smallskip\noindent{\small\emph{Used in:} #1}\par\smallskip}
\@ifundefined{c@fact}{\newtheorem{fact}{Fact}[section]}{}
\makeatother

\section{Toolbox: facts from analysis and geometry}\label{app:toolbox}

This appendix collects the facts from analysis, convex geometry, and
measure theory on which the proofs in Appendices~\PN{appx-sec:smoothing}{B}
and~\PN{appx-sec:support}{C} rely and which a mathematical statistician may
not have at hand. Every surface in those proofs is a finite union of flat
pieces of hyperplanes, so no differential geometry beyond the flat case is
needed. On a hyperplane the area formula is Fubini's theorem, the
Gauss--Green formula (the divergence theorem) is the fundamental theorem of
calculus along parallel lines, and the coarea formula is Fubini's theorem
applied piece by piece; the proofs below carry out each of these
reductions. Each fact is stated in the form in which
the proofs use it, and is followed by the places that use it and by a proof
or a reference. A proof is given when it takes a few lines, and a standard
source, listed at the end of the appendix, is cited otherwise. The facts are
standard results and not claims of the paper, so they carry their own
numbering. The symbols in the lines headed ``Used in'' are those of the
paper's appendices.

Subsection~\ref{tb:sec:calc} treats integration and calculus,
Subsection~\ref{tb:sec:gauss} the Gaussian identities, and
Subsection~\ref{tb:sec:matrix} the matrix and entropy inequalities.
Subsection~\ref{tb:sec:logconcave} covers log-concavity and
Pr\'ekopa's theorem. Subsections~\ref{tb:sec:poly}--\ref{tb:sec:coarea}
treat polyhedra, from faces and normal cones to the divergence theorem,
tubes, and the coarea formula, in the order in which
\PR{Lemma}{appx-an:lem:surface}{C.5} needs them, and
Subsection~\ref{tb:sec:index} lists where each fact is used.

\subsection{Notation}\label{tb:sec:notation}

The space \(\R^n\) carries the inner product \(x\cdot y\) and the Euclidean
norm \(\norm x\); in the proofs \(n\) is the rank \(r\) of the covariance,
or \(r-1\) when a hyperplane of \(\R^r\) is identified with
\(\R^{r-1}\). We write \(dx\) for Lebesgue measure, \(\varphi_n\) and
\(\Phi_n\) for the standard Gaussian density and law on \(\R^n\), and \(G\)
for a standard Gaussian vector. For a linear subspace \(L\), \(\varphi_L\)
is the standard Gaussian density on \(L\) with respect to Lebesgue measure
on \(L\). For a set \(S\), \(\ind_S\) is its indicator function, \(\aff S\)
is the smallest affine subspace containing it, the dimension of \(S\) is
that of \(\aff S\), and \(\relint S\) is the interior of \(S\) relative to
\(\aff S\). The distance from a point \(y\) to a set \(S\) is
\(\dist(y,S)=\inf_{x\in S}\norm{y-x}\).

For \(f\in C^q(\R^n)\), the derivative of order \(q\) in the directions
\(z_1,\ldots,z_q\in\R^n\) is
\[
 D^qf(x)[z_1,\ldots,z_q]
 =\partial_{s_1}\cdots\partial_{s_q}f(x+s_1z_1+\cdots+s_qz_q)\Big|_{s=0}
 =\sum_{i_1,\ldots,i_q}\partial_{i_1\cdots i_q}f(x)\,
   z_{1,i_1}\cdots z_{q,i_q},
\]
which is symmetric and multilinear in the directions. We write
\(D^qf(x)[z^{\otimes q}]\) when all directions equal \(z\), and
\(D^1f=\nabla f\). The Loewner order \(M\preceq N\) means that
\(N-M\) is positive semidefinite, \(\norm\Xi_\op\) is the operator norm
of a matrix \(\Xi\), and \(\norm M_*\) is the \emph{nuclear norm}, which for
a symmetric matrix \(M\) with eigenvalues \(\mu_i\) equals \(\sum_i|\mu_i|\).

\subsection{Integration and calculus}\label{tb:sec:calc}

\begin{fact}[Fubini and Tonelli]\label{tb:fubini}
Let \((S_1,\mu_1)\) and \((S_2,\mu_2)\) be \(\sigma\)-finite measure
spaces and let \(f\) be a measurable function on \(S_1\times S_2\). If
\(f\ge0\), then \(s_1\mapsto\int f(s_1,s_2)\,\mu_2(ds_2)\) is measurable
and
\[
 \int\Bigl(\int f\,d\mu_2\Bigr)d\mu_1
 =\int\Bigl(\int f\,d\mu_1\Bigr)d\mu_2
 =\int f\,d(\mu_1\otimes\mu_2)\in[0,\infty].
\]
If \(\int|f|\,d(\mu_1\otimes\mu_2)<\infty\), the same equalities hold, and
the inner integrals are finite almost everywhere.
\end{fact}

Three special cases occur in the proofs. If \(X\) and \(Y\) are
independent random elements and either \(f\ge0\) or \(f(X,Y)\) is
integrable, then
\[
 \E f(X,Y)=\E g(X),\qquad g(x)=\E f(x,Y),
\]
because the joint law is the product of the marginal laws. This is the
formal content of every argument that is carried out for a fixed value of
an independent component such as \(V_0\) and then averaged over it. With
\(\mu_2\) the counting measure, the theorem says that
\(\E\sum_k\xi_k=\sum_k\E\xi_k\) for nonnegative terms, and for terms with
\(\sum_k\E|\xi_k|<\infty\). For an orthogonal decomposition
\(\R^n=L\oplus L^\perp\), Lebesgue measure on \(\R^n\) is the product of the
Lebesgue measures on \(L\) and on \(L^\perp\), so that
\(\int_{\R^n}f\,dx=\int_{L^\perp}\int_Lf(l+w)\,dl\,dw\).

The appeal to Tonelli's theorem in Step~3 of the proof of
\PR{Lemma}{appx-an:lem:surface}{C.5}, whose notation we use in this
paragraph, is of the first kind. With \(V_0\) and
\(G\) independent, \(\Phi_r(\mathcal O_\iota(V_0))\) is the nonnegative
integral \(\int\ind\{g\in\mathcal O_\iota(V_0)\}\varphi_r(g)\,dg\), so
\(\E\sum_\iota R_\iota^3\,\Phi_r(\mathcal O_\iota(V_0))\) equals
\(\int\sum_\iota R_\iota^3\,\Pp\{g\in\mathcal O_\iota(V_0)\}\varphi_r(g)\,dg\)
with no check of integrability, since every term is nonnegative, and the
probability can then be bounded for each fixed \(g\).

\usedin{conditioning on \(V_0\) or on the \(\sigma\)-field \(\mathcal H\) in
\PR{Lemmas}{appx-an:lem:restoration}{C.2}, \PN{appx-an:lem:matrix}{C.4}
and~\PN{appx-an:lem:surface}{C.5}; conditioning on the exponentials in
\PR{Lemmas}{appx-lem:envelope}{B.1} and~\PN{appx-lem:endpoint-free}{B.2}
and on the candidate pairs in \PR{Appendix}{appx-app:concentration}{D};
interchanges of sums and expectations at the end of the proof of
\PR{Lemma}{appx-an:lem:restoration}{C.2}; Step~3 of the proof of
\PR{Lemma}{appx-an:lem:surface}{C.5}; iterated integrals over a hyperplane
and its normal in Subsections~\ref{tb:sec:surface}--\ref{tb:sec:coarea}.}
\emph{Source.} \TBc{Fol99}, Section~2.5.

\begin{fact}[Differentiation under the integral sign; Taylor's formula]
\label{tb:taylor}
(a) Let \(I\) be an open interval and let \(f(\lambda,\omega)\) be
integrable in \(\omega\) and differentiable in \(\lambda\in I\), with
\(|\partial_\lambda f(\lambda,\omega)|\le g(\omega)\) for an integrable
function \(g\). Then \(\frac d{d\lambda}\E f(\lambda,\cdot)
=\E\,\partial_\lambda f(\lambda,\cdot)\) on \(I\).

(b) For \(f\in C^q(\R^n)\) and \(x,v\in\R^n\),
\[
 f(x+v)=\sum_{k=0}^{q-1}\frac1{k!}D^kf(x)[v^{\otimes k}]
 +\int_0^1\frac{(1-s)^{q-1}}{(q-1)!}D^qf(x+sv)[v^{\otimes q}]\,ds.
\]
For \(q=1\) this is \(f(x+v)-f(x)=\int_0^1Df(x+sv)[v]\,ds\).
\end{fact}

\begin{proof}
For (a), the mean value theorem bounds the difference quotients by \(g\),
and they converge pointwise, so dominated convergence applies. For (b),
the function \(s\mapsto f(x+sv)\) is of class \(C^q\) with \(k\)th derivative
\(D^kf(x+sv)[v^{\otimes k}]\); apply the one-dimensional Taylor formula with
integral remainder at \(s=0\) and evaluate it at \(s=1\).
\end{proof}

\usedin{\PR{Lemma}{appx-an:lem:identity}{C.1}, where the derivative of
\(P_A(\lambda)=\E F(U_\lambda)\) is taken under the expectation (the
summands are finitely many and bounded) and Steps~1 and~3 use (b), to degree
three with the fourth-order integral remainder, to degree two for the
derivative, and to first order for the difference of third derivatives;
\PR{Lemma}{appx-lem:endpoint-gap}{B.3}, where (b) with \(q=1\) writes
\(\E[W\cdot\nabla\Psi(W)]\) as an integral of the Hessian along the segment
from \(W_i\) to \(W\).}
\emph{Source.} \TBc{Fol99}, Section~2.3.

\begin{fact}[Fa\`a di Bruno's formula]\label{tb:faa}
Let \(\phi\) be a smooth function on an open interval of \(\R\) and let
\(g:\R^n\to\R\) be smooth with values in that interval. For \(q\ge1\), let
\(\Pi_q\) be the set of partitions \(\mathcal P\) of \(\{1,\ldots,q\}\) into
nonempty blocks, let \(|\mathcal P|\) be the number of blocks, and for a
block \(J\) let \(z_J\) be the list of directions \(z_a\) with
\(a\in J\). Then
\[
 D^q(\phi\circ g)(x)[z_1,\ldots,z_q]
 =\sum_{\mathcal P\in\Pi_q}\phi^{(|\mathcal P|)}(g(x))
  \prod_{J\in\mathcal P}D^{|J|}g(x)[z_J].
\]
\end{fact}

\begin{proof}
For \(q=1\) this is the chain rule. Differentiate the right side for \(q\)
in the direction \(z_{q+1}\). By the product rule, the derivative either
falls on \(\phi^{(|\mathcal P|)}(g(x))\), which gives
\(\phi^{(|\mathcal P|+1)}(g(x))\,Dg(x)[z_{q+1}]\) and adds the block
\(\{q+1\}\) to \(\mathcal P\), or falls on one factor
\(D^{|J|}g(x)[z_J]\), which adds \(q+1\) to the block \(J\). Every partition
of \(\{1,\ldots,q+1\}\) arises exactly once in this way, by removing
\(q+1\) from its block.
\end{proof}

\begin{fact}[Derivatives of the logistic approximation]\label{tb:logistic}
Let \(\nu_1,\ldots,\nu_m\in\R^n\), \(a_1,\ldots,a_m\in\R\), \(h>0\), and
\[
 w_\alpha(x)=e^{(\nu_\alpha\cdot x-a_\alpha)/h},\qquad
 T=\sum_\alpha w_\alpha,\qquad
 F=\frac1{1+T},\qquad
 \pi_\alpha=\frac{w_\alpha}{1+T}=Fw_\alpha .
\]
For directions \(z_1,\ldots,z_q\), with the blocks \(J_1,\ldots,J_{|\mathcal P|}\)
of \(\mathcal P\) listed in any fixed order,
\[
 D^qF(x)[z_1,\ldots,z_q]
 =h^{-q}\sum_{\mathcal P\in\Pi_q}(-1)^{|\mathcal P|}|\mathcal P|!
 \sum_{\alpha_1,\ldots,\alpha_{|\mathcal P|}}
 F(x)\prod_{b=1}^{|\mathcal P|}\pi_{\alpha_b}(x)
 \prod_{b=1}^{|\mathcal P|}\prod_{a\in J_b}(\nu_{\alpha_b}\cdot z_a),
\]
and \(\sum_{\mathcal P\in\Pi_q}|\mathcal P|!\le q!\,2^{q-1}\). Consequently
\(|D^qF(x)[z_1,\ldots,z_q]|\le q!\,2^{q-1}h^{-q}\,F(x)\prod_{a=1}^q
\max_\alpha|\nu_\alpha\cdot z_a|\).
\end{fact}

\begin{proof}
Apply Fact~\ref{tb:faa} with \(\phi(\tau)=(1+\tau)^{-1}\) and \(g=T\). Then
\(\phi^{(k)}(\tau)=(-1)^kk!\,(1+\tau)^{-k-1}\), and since each \(w_\alpha\) is
an exponential of an affine function,
\(D^lT(x)[z_1,\ldots,z_l]=h^{-l}\sum_\alpha w_\alpha(x)\prod_{a\le l}
(\nu_\alpha\cdot z_a)\). Substituting, and using
\((1+T)^{-|\mathcal P|-1}\prod_bw_{\alpha_b}=F\prod_b\pi_{\alpha_b}\) and
\(\sum_b|J_b|=q\), gives the formula. The number
\(\sum_{\mathcal P}|\mathcal P|!\) counts ordered partitions
\((J_1,\ldots,J_k)\) of \(\{1,\ldots,q\}\). Writing the elements of
\(J_1\), then of \(J_2\), and so on, in increasing order inside each block,
gives a permutation together with the set of positions at which a new block
starts, which is a subset of the \(q-1\) gaps between consecutive entries.
Every ordered partition arises from at least one of these
\(q!\,2^{q-1}\) pairs. For the bound, the sum over \(\alpha_b\) of
\(\pi_{\alpha_b}\) is \(1-F\le1\), and each product of projections is
at most \(\prod_a\max_\alpha|\nu_\alpha\cdot z_a|\).
\end{proof}

\usedin{Subsection~\PN{appx-sec:mixture}{B.2}, where the first display
and the bound are \PE{appx-eq:partition-derivative}{B.1} and
\PE{appx-eq:all-orders}{B.2} (the text there cites the softmax cumulant formula,
\TBc{GZ}, Lemma~13); the envelope~\PE{appx-eq:local-envelope}{B.3} is
obtained by summing the same expansion over coordinate tuples;
\PR{Lemma}{appx-an:lem:restoration}{C.2}, Step~1, with \(u_\alpha\) in place
of \(\nu_\alpha\).}
\emph{Source.} \TBc{Com74}; the proof above is complete.

\subsection{Gaussian identities}\label{tb:sec:gauss}

\begin{fact}[Gaussian density in orthogonal coordinates]\label{tb:gauss-coord}
(a) Let \(\R^n=L\oplus L^\perp\) be an orthogonal decomposition. Then
\(\varphi_n(l+w)=\varphi_L(l)\varphi_{L^\perp}(w)\) for \(l\in L\) and
\(w\in L^\perp\), so the projections of a standard Gaussian vector onto
\(L\) and \(L^\perp\) are independent standard Gaussian vectors. In
particular, for a unit vector \(u\), \(a\in\R\), and \(y\perp u\),
\(\varphi_n(au+y)=\varphi_1(a)\varphi_{n-1}(y)\).

(b) For \(x,v\in\R^n\),
\(\varphi_n(x+v)=\varphi_n(x)\exp\{-x\cdot v-\tfrac12\norm v^2\}\).
\end{fact}

\begin{proof}
Both follow from \(\norm{l+w}^2=\norm l^2+\norm w^2\) and from expanding
\(\norm{x+v}^2\).
\end{proof}

\usedin{\PR{Lemma}{appx-an:lem:mean}{C.3}, where the normal and tangential
components of \(G\) relative to a face are independent;
\PR{Lemma}{appx-an:lem:surface}{C.5}, Steps~1 and~2, for the factorization of
\(\varphi_r\) on a facet and on a ridge and for the density ratio on a tube.}

\begin{fact}[Gaussian integration by parts]\label{tb:stein}
Let \(G\sim N(0,I_n)\) and let \(f\in C^2(\R^n)\) have \(f\) and its partial
derivatives of orders at most two of polynomial growth. Then
\begin{enumerate}
\item[(a)] \(\E[Gf(G)]=\E\nabla f(G)\);
\item[(b)] \(\E[(GG^\top-I_n)f(G)]=\E D^2f(G)\); equivalently, for \(z\in\R^n\),
\(\E D^2f(G)[z,z]=\E\bigl[\tr\{(GG^\top-I_n)zz^\top\}f(G)\bigr]\), where
\(\tr\{(GG^\top-I_n)zz^\top\}=(z\cdot G)^2-\norm z^2\);
\item[(c)] if \(v\in\R^n\), \(\sigma>0\), \(U=v+\sigma G\), and
\(F\in C^{s+2}(\R^n)\) has derivatives of polynomial growth, then for
\(z,z_1,\ldots,z_s\in\R^n\),
\[
 \E D^{s+2}F(U)[z,z,z_1,\ldots,z_s]
 =\sigma^{-2}\E\bigl[\tr\{(GG^\top-I_n)zz^\top\}\,
 D^sF(U)[z_1,\ldots,z_s]\bigr].
\]
\end{enumerate}
\end{fact}

\begin{proof}
For (a) in one coordinate, \(x\varphi_1(x)=-\varphi_1'(x)\), so integration by
parts gives
\[\int xf(x)\varphi_1(x)\,dx=\int f'(x)\varphi_1(x)\,dx,\]
and the boundary terms vanish because of the growth condition. The other coordinates
are handled by Fact~\ref{tb:fubini}. For (b), apply (a) to \(\partial_if\) and
to \(G_kf\): \(\E[G_iG_kf(G)]=\E\partial_i(G_kf)(G)=\delta_{ik}\E f(G)
+\E[G_k\partial_if(G)]=\delta_{ik}\E f(G)+\E\partial_{ik}f(G)\). For (c), the
function \(f(g)=D^sF(v+\sigma g)[z_1,\ldots,z_s]\) has
\(D^2f(g)[z,z]=\sigma^2D^{s+2}F(v+\sigma g)[z,z,z_1,\ldots,z_s]\); apply (b).
\end{proof}

\usedin{Subsection~\PN{sec:outline-late}{2.4} (the display after
``integration by parts against the covariance of the summands''),
\PR{Lemma}{appx-an:lem:restoration}{C.2}, Step~1, where (c) is applied
conditionally on \(\mathcal H\) with \(v=V_0+v_i\) and \(z=Z_i\);
\PR{Lemma}{appx-lem:endpoint-gap}{B.3}, for \(\E[Z\cdot\nabla\Psi(Z)]
=\E\tr(\Sigma D^2\Psi(Z))\), which is (a) with \(Z=\Gamma G\) for a factor
\(\Gamma\) of the covariance.}
\emph{Source.} Stein~\TBc{Ste81}; the proof above is complete.

\begin{fact}[Heat equation for the Gaussian component]\label{tb:heat}
Let \(Z\sim N(0,\Sigma)\) in \(\R^p\), where \(\Sigma\) may be singular, and
let \(f\in C^2(\R^p)\) be bounded with bounded first and second derivatives.
Then, for \(x\in\R^p\) and \(\lambda>0\),
\[
 \frac d{d\lambda}\E f(x+\sqrt\lambda Z)
 =\tfrac12\E\tr\{\Sigma D^2f(x+\sqrt\lambda Z)\}.
\]
\end{fact}

\begin{proof}
Write \(Z=\Gamma G\) with \(\Gamma\Gamma^\top=\Sigma\) and \(G\) standard
Gaussian. By Fact~\ref{tb:taylor}(a), which applies because the difference
quotients are bounded by \(\norm\nabla f_\infty\norm{\Gamma G}/(2\sqrt\lambda)\)
locally uniformly in \(\lambda>0\), the derivative equals
\((2\sqrt\lambda)^{-1}\E[G\cdot\Gamma^\top\nabla f(x+\sqrt\lambda\Gamma G)]\).
Fact~\ref{tb:stein}(a) applied to the function \(g\mapsto
\Gamma^\top\nabla f(x+\sqrt\lambda\Gamma g)\) turns this into
\((2\sqrt\lambda)^{-1}\sqrt\lambda\,\E\tr\{\Gamma^\top D^2f\,\Gamma\}
=\frac12\E\tr\{\Sigma D^2f\}\).
\end{proof}

\usedin{Subsection~\PN{sec:outline-path}{2.2} and the proof of
\PR{Lemma}{appx-an:lem:identity}{C.1}, Step~2, where the Gaussian part of the
path contributes \(\frac12\sum_iM_i:\E D^2F(U_\lambda)\).}

\subsection{Matrix and entropy inequalities}\label{tb:sec:matrix}

\begin{fact}[Duality between the operator and nuclear norms]\label{tb:nuclear}
(a) For every square matrix \(\Xi\) and every symmetric matrix \(M\),
\(|\tr(\Xi M)|\le\norm\Xi_\op\norm M_*\). Moreover
\(\norm M_*=\max\{\tr(\Xi M):\norm\Xi_\op\le1\}\), so that \(\norm\cdot_*\) is
a norm on symmetric matrices.

(b) If \(M=P-N\) with \(P\succeq0\) and \(N\succeq0\), then
\(\norm M_*\le\tr P+\tr N=2\tr P-\tr M\).
\end{fact}

\begin{proof}
For (a), let \(M=\sum_i\mu_ie_ie_i^\top\) be a spectral decomposition. Then
\(\tr(\Xi M)=\sum_i\mu_i\,e_i^\top\Xi e_i\) and \(|e_i^\top\Xi e_i|\le
\norm\Xi_\op\), which gives the inequality. Equality is attained at
\(\Xi=\sum_i\operatorname{sgn}(\mu_i)e_ie_i^\top\). A supremum of linear
functionals satisfies the triangle inequality. For (b), the norm of a
positive semidefinite matrix is its trace, so the triangle inequality gives
\(\norm M_*\le\tr P+\tr N\), and \(\tr N=\tr P-\tr M\).
\end{proof}

\usedin{Subsection~\PN{sec:outline-late}{2.4}; \PR{Lemma}{appx-an:lem:restoration}{C.2},
Step~4; \PR{Lemma}{appx-an:lem:matrix}{C.4}, Step~1, where
(b) is applied with \(P=p_\alpha\mu_\alpha\mu_\alpha^\top\) and
\(N=p_\alpha(I_r-\Sigma_\alpha)\), and gives
\PE{appx-an:matrix-trace-reduction}{C.17}.}

\begin{fact}[Log-sum inequality and a concave function]\label{tb:logsum}
(a) For positive numbers \(y_1,\ldots,y_l\) and \(z_1,\ldots,z_l\) with sums
\(y\) and \(z\),
\(\sum_jy_j\log(z_j/y_j)\le y\log(z/y)\).

(b) The function \((x,y)\mapsto x\log(y/x)\) is concave on \((0,\infty)^2\).
Hence, for positive integrable random variables \(X,Y\),
\(\E[X\log(Y/X)]\le\E X\,\log(\E Y/\E X)\).
\end{fact}

\begin{proof}
For (a), Jensen's inequality for the concave function \(\log\), with the
weights \(y_j/y\), gives \(\sum_j(y_j/y)\log(z_j/y_j)\le\log\sum_j(y_j/y)(z_j/y_j)
=\log(z/y)\). For (b), the function \(f(x,y)=x\log y-x\log x\) has Hessian
\(\begin{pmatrix}-1/x&1/y\\1/y&-x/y^2\end{pmatrix}\), whose determinant is
zero and whose trace is negative, so the Hessian is negative semidefinite;
the second statement is Jensen's inequality for the concave function
\(f\) and the random vector \((X,Y)\).
\end{proof}

\usedin{\PR{Lemma}{appx-an:lem:matrix}{C.4}, Step~3: (a) adds the reference
masses of the tuples that share a first face, and (b) averages over
\(V_0\) and over the faces.}
\emph{Source.} \TBc{CT06}, Theorem~2.7.1, for (a).

\subsection{Log-concavity}\label{tb:sec:logconcave}

A function \(f:\R^n\to[0,\infty)\) is \emph{log-concave} if
\(f((1-t)x+ty)\ge f(x)^{1-t}f(y)^t\) for all \(x,y\) and \(t\in[0,1]\);
when \(f>0\) this says that \(\log f\) is concave.

\begin{fact}[Pr\'ekopa--Leindler inequality and Pr\'ekopa's theorem]
\label{tb:prekopa}
(a) Let \(0<t<1\) and let \(f,g,h:\R^n\to[0,\infty)\) be bounded,
continuous, and integrable, with
\(h((1-t)x+ty)\ge f(x)^{1-t}g(y)^t\) for all \(x,y\). Then
\(\int h\ge(\int f)^{1-t}(\int g)^t\).

(b) Let \(f:\R^a\times\R^b\to[0,\infty)\) be bounded, continuous, and
log-concave, with \(f(z,\cdot)\) integrable for every \(z\). Then
\(z\mapsto\int f(z,y)\,dy\) is log-concave on \(\R^a\).
\end{fact}

\begin{proof}
(a) For \(n=1\), assume first that \(\sup f=\sup g=1\); the general case
follows by dividing \(f,g,h\) by \(\sup f\), \(\sup g\), and
\((\sup f)^{1-t}(\sup g)^t\). For \(0<s<1\) the level sets
\(\{f\ge s\}\) and \(\{g\ge s\}\) are nonempty, and the hypothesis gives
\(\{h\ge s\}\supseteq(1-t)\{f\ge s\}+t\{g\ge s\}\). In one dimension
\(|A+B|\ge|A|+|B|\) for nonempty compact sets \(A\) and \(B\), because
\(A+B\) contains the translates \(A+\min B\) and \(\max A+B\), which meet in
one point; the inequality passes to measurable sets by approximation from
inside by compact sets. Therefore
\(|\{h\ge s\}|\ge(1-t)|\{f\ge s\}|+t|\{g\ge s\}|\), and integration over
\(s\in(0,1)\) gives \(\int h\ge(1-t)\int f+t\int g\ge(\int f)^{1-t}(\int g)^t\),
by the inequality between the arithmetic and geometric means. For \(n\ge2\),
apply the case of dimension \(n-1\) to the sections of \(f,g,h\) at the last
coordinates \(s_0\), \(s_1\), and \((1-t)s_0+ts_1\), and then the case of
dimension one to the integrals of the sections, with Fubini's theorem
(Fact~\ref{tb:fubini}) connecting the steps; the routine measurability
checks are omitted here and can be found in~\TBc{Gar02}.

(b) Fix \(z_0,z_1\) and put \(z_t=(1-t)z_0+tz_1\). Log-concavity of \(f\)
gives \(f(z_t,(1-t)y_0+ty_1)\ge f(z_0,y_0)^{1-t}f(z_1,y_1)^t\), so (a) applies
to \(f(z_0,\cdot)\), \(f(z_1,\cdot)\), and \(f(z_t,\cdot)\).
\end{proof}

\emph{Source.} Pr\'ekopa~\TBc{Pre71,Pre73}, Leindler~\TBc{Lei72}; a
survey with the proof of the general case is~\TBc{Gar02}, Section~7. The
paper cites Theorem~6 of~\TBc{Pre73}.

\begin{fact}[Gaussian laws tilted by a log-concave weight]\label{tb:tilt}
Let \(w:\R^r\to[0,\infty)\) be bounded, continuous, and log-concave, with
\(p=\E w(G)>0\) for \(G\sim N(0,I_r)\). Let \(\mu\) and \(K\) be the mean and
covariance matrix of the probability law with density \(w(g)\varphi_r(g)/p\).
Then \(K\preceq I_r\), and
\(\E[(GG^\top-I_r)w(G)]=p(K-I_r+\mu\mu^\top)\).
\end{fact}

\begin{proof}
The function \((z,g)\mapsto e^{-\norm{g-z}^2/2}w(g)\) is log-concave, since
\(\norm{g-z}^2\) is jointly convex in \((z,g)\), and it is bounded and
continuous. By Fact~\ref{tb:prekopa}(b), \(\ell(z)=\log\int e^{-\norm{g-z}^2/2}
w(g)\,dg\) is concave. Since \(\norm{g-z}^2=\norm g^2-2z\cdot g+\norm z^2\),
\[
 \ell(z)=\text{const}-\tfrac12\norm z^2+\log M(z),\qquad
 M(z)=\int e^{z\cdot g}w(g)\varphi_r(g)\,dg=p\,\E_\nu e^{z\cdot g},
\]
where \(\nu\) is the tilted law. The function \(M\) is finite and smooth,
because \(w\) is bounded, and the Hessian of \(\log M\) at \(0\) is the
covariance \(K\) of \(\nu\). A concave function of class \(C^2\) has a
negative semidefinite Hessian, so \(K-I_r\preceq0\). The identity is the
decomposition of the second moment \(\E_\nu gg^\top=K+\mu\mu^\top\).
\end{proof}

\usedin{Subsection~\PN{sec:outline-late}{2.4} (``Pr\'ekopa's theorem \dots\
has covariance \(K\preceq I_r\)'') and \PR{Lemma}{appx-an:lem:matrix}{C.4},
Step~1, where \(w=w_\alpha\) and \(K=\Sigma_\alpha\).}
\emph{Source.} For indicators of convex cells, the cell covariance bound
\TBc{GZ}, Lemma~11, is the same estimate \(K\preceq I_r\), obtained from the
Brascamp--Lieb inequality~\TBc{BL76}; the weights here are not indicators,
which is why the proof goes through Pr\'ekopa's theorem.

\begin{fact}[Log-concavity of the logistic weights]\label{tb:lc-weights}
With the notation of Fact~\ref{tb:logistic}, for every \(k\ge0\) and every
tuple \((\alpha_1,\ldots,\alpha_k)\), the function
\(F\pi_{\alpha_1}\cdots\pi_{\alpha_k}=\exp\{\sum_{a\le k}t_{\alpha_a}\}
(1+T)^{-k-1}\), where \(t_\alpha=(\nu_\alpha\cdot x-a_\alpha)/h\), is
log-concave on \(\R^n\), and so is its composition
\(g\mapsto F\pi_{\alpha_1}\cdots\pi_{\alpha_k}(v+\sigma g)\) with any affine map.
\end{fact}

\begin{proof}
The logarithm is a sum of the affine functions \(t_{\alpha_a}\) and
\(-(k+1)\log(1+\sum_\alpha e^{t_\alpha})\). The function
\(\log(1+\sum_\alpha e^{t_\alpha})\) is the logarithm of a sum of
exponentials of the affine functions \(0,t_1,\ldots,t_m\) and is convex by
H\"older's inequality: \(\sum_\alpha e^{(1-s)t_\alpha+st'_\alpha}\le
(\sum_\alpha e^{t_\alpha})^{1-s}(\sum_\alpha e^{t'_\alpha})^s\). Composition
with an affine map preserves concavity.
\end{proof}

\usedin{Subsection~\PN{sec:outline-late}{2.4} and Step~1 of the proof of
\PR{Lemma}{appx-an:lem:matrix}{C.4}.}

\subsection{Polyhedra: faces, normal cones, and nearest points}\label{tb:sec:poly}

A \emph{polyhedron} is a set
\[
 Q=\bigcap_{j=1}^m\{x\in\R^n:u_j\cdot x\le a_j\},
\]
with unit vectors \(u_j\), called \emph{outward normals}, and numbers
\(a_j\), called \emph{offsets}. It is closed and convex. It is
\emph{full-dimensional} if it has an interior point, and \emph{proper} if it
is neither empty nor \(\R^n\). The description is \emph{irredundant} if
deleting any one inequality changes the set; deleting redundant or repeated
inequalities one at a time turns every description into an irredundant one
of the same set. For \(j\le m\), put \(H_j=\{x:u_j\cdot x=a_j\}\) and
\(F_j=Q\cap H_j\). A \emph{face} of \(Q\) is \(Q\) itself or a set
\(Q\cap\{u\cdot x=c\}\) with \(u\ne0\) and \(Q\subseteq\{u\cdot x\le c\}\).
A \emph{facet} is a face of dimension \(n-1\), and a \emph{ridge} is a face
of dimension \(n-2\). In the plane the facets are the sides of a polygon
and the ridges its corners, in \(\R^3\) the facets are the faces of a solid
and the ridges its edges, and on the line the facets are the endpoints of an
interval, which has no ridges. We use these standard properties of faces:
the relative interiors of the faces of \(Q\) are pairwise disjoint and cover
\(Q\); a face of a face of \(Q\) is a face of \(Q\); and a face \(G'\)
of a face \(G\) with \(G'\ne G\) lies in the relative boundary
\(G\setminus\relint G\) of \(G\) (\TBc{Roc70}, Section~18; \TBc{Zie95}, Chapter~2).

\begin{fact}[Facets and ridges of a polyhedron]\label{tb:facets}
Let \(Q\) be a proper, full-dimensional polyhedron in \(\R^n\) with an
irredundant description.
\begin{enumerate}
\item[(a)] The normals \(u_1,\ldots,u_m\) are pairwise distinct, each
\(F_j\) is a facet, \(Q\) has no other facets, and
\(\partial Q=F_1\cup\cdots\cup F_m\).
\item[(b)] For \(j\ne k\), \(F_j\cap F_k\subseteq H_j\cap H_k\), and
\(H_j\cap H_k\) is empty if \(u_k=-u_j\), and an affine subspace of
dimension \(n-2\) otherwise. In particular \(F_j\cap F_k\) lies in a
finite union of affine subspaces of dimension at most \(n-2\).
\item[(c)] Each \(F_j\) is a full-dimensional polyhedron in the hyperplane
\(H_j\). Write the points of \(H_j\) as \(a_ju_j+y\) with \(y\perp u_j\), and,
for \(k\ne j\) with \(u_k\ne\pm u_j\), put
\(\alpha_{jk}=\arccos(u_j\cdot u_k)\in(0,\pi)\),
\[
 \nu_{k\mid j}=\frac{u_k-(u_j\cdot u_k)u_j}{\sin\alpha_{jk}},\qquad
 d_{k\mid j}=\frac{a_k-(u_j\cdot u_k)a_j}{\sin\alpha_{jk}} .
\]
Then \(F_j\) is described in the coordinate \(y\) by the inequalities
\(\nu_{k\mid j}\cdot y\le d_{k\mid j}\), and its facets are exactly the
ridges \(E_{jk}=F_j\cap F_k\) of \(Q\) of dimension \(n-2\); the outward
unit normal of \(F_j\) at \(E_{jk}\) is \(\nu_{k\mid j}\).
\item[(d)] Every ridge of \(Q\) lies in exactly two facets, so the ridges
correspond to the unordered pairs \(\{j,k\}\) with
\(\dim(F_j\cap F_k)=n-2\); for such a pair \(u_k\ne\pm u_j\), and the
only inequalities active at a point of \(\relint E_{jk}\) are \(j\) and \(k\).
\end{enumerate}
\end{fact}

\begin{proof}[Proof (sketch)]
(a) Fix \(j\). Irredundancy gives a point \(x_j\) with \(u_j\cdot x_j>a_j\)
and \(u_k\cdot x_j\le a_k\) for \(k\ne j\). Let \(x_0\) be an interior point of
\(Q\), and let \(x_*\) be the point where the segment from \(x_0\) to \(x_j\)
crosses \(H_j\). For \(k\ne j\), \(u_k\cdot x_*<a_k\), because \(x_*\) is a
convex combination of \(x_0\), where the inequality is strict, with positive
weight, and of \(x_j\), where it holds weakly. Hence a neighborhood of
\(x_*\) in \(H_j\) lies in \(F_j\), and \(\dim F_j=n-1\). Two inequalities
with the same normal cannot both be irredundant, since the one with the
larger offset is implied by the other. A point of \(Q\) at
which no inequality is active is interior, and at a point where
inequality \(j\) is active the points \(x+\varepsilon u_j\) lie outside \(Q\),
so \(\partial Q=\bigcup_jF_j\). If \(G=Q\cap\{u\cdot x=c\}\) is a facet and
\(x\in\relint G\), some inequality \(j\) is active at \(x\); on \(G\) the
affine function \(u_j\cdot x-a_j\le0\) attains its maximum at a relative
interior point and is therefore zero, so \(G\subseteq H_j\), and since \(G\) has
dimension \(n-1\), \(G=Q\cap H_j=F_j\).

(b) If \(u_k=-u_j\) and \(H_j\cap H_k\ne\emptyset\), then \(Q\) would lie in
the hyperplane \(H_j\), contradicting full dimension. If \(u_k\ne\pm u_j\),
the two equations are independent.

(c) The set \(F_j\) is the polyhedron in \(H_j\) cut out by the remaining
inequalities, which take the stated form in the coordinate \(y\) because
\(u_k\cdot(a_ju_j+y)=a_ju_j\cdot u_k+(u_k-(u_j\cdot u_k)u_j)\cdot y\); an
inequality with \(u_k=-u_j\) holds automatically on \(H_j\) since
\(F_j\ne\emptyset\). By (a) applied inside \(H_j\), the facets of \(F_j\) are
the faces of \(F_j\) of dimension \(n-2\), hence faces of \(Q\) of that
dimension, that is, ridges of \(Q\) in \(F_j\); each equals
\(F_j\cap H_k=F_j\cap F_k\) for the inequality \(k\) that defines it.

(d) Let \(E\) be a ridge and \(y\in\relint E\), and let \(J\) be the set of
inequalities active at \(y\). On \(E\) each affine function
\(u_i\cdot x-a_i\), \(i\in J\), is at most zero and vanishes at the relative
interior point \(y\), so it vanishes on \(E\); thus \(E\subseteq F_i\) for
\(i\in J\) and \(\dim\operatorname{span}\{u_i:i\in J\}\le2\). A single active
inequality is impossible, since then \(y\) would lie in the relative interior
of \(F_j\) and the only face of \(F_j\) containing it is \(F_j\) itself, of
dimension \(n-1\). Opposite normals \(u_i=-u_j\) are excluded by (b). Hence
the active normals span a plane \(\Pi\). Near \(y\), \(Q\) coincides with
\(y\) plus the product of \(\Pi^\perp\) and a planar cone \(K\subseteq\Pi\)
with nonempty interior, defined by the active inequalities and contained in
a wedge of opening \(\pi-\alpha_{jk}<\pi\); this cone has exactly two
boundary rays. Each active inequality \(i\) has \(\dim F_i=n-1\), so
\(F_i\) meets every neighborhood of \(y\) in a set of dimension \(n-1\), and
the line \(\{u_i\cdot v=0\}\) must therefore carry a boundary ray of \(K\).
With distinct normals this leaves exactly two active inequalities.
\end{proof}

\usedin{\PR{Lemma}{appx-an:lem:surface}{C.5}, in the paragraph on the terms
from the geometry of polyhedra and in Step~1 (``each inequality defines
exactly one facet, \dots\ two facets that share a ridge have normals that
are neither equal nor opposite'') and Step~2 (``every ridge lies in exactly
two facets'').}
\emph{Source.} \TBc{Zie95}, Chapter~2; \TBc{Roc70}, Sections~18 and~19;
the facet decomposition of \(\partial Q\) with the pairing of ridges is
the content of \TBc{GZ}, Lemma~27.

\begin{fact}[Normal cones]\label{tb:normalcone}
Let \(Q\) be a polyhedron as in Fact~\ref{tb:facets}. For \(x\in Q\) let
\(J(x)=\{j:u_j\cdot x=a_j\}\), and let
\[
 N_Q(x)=\{v\in\R^n:v\cdot(q-x)\le0\text{ for all }q\in Q\}
\]
be the \emph{normal cone} of \(Q\) at \(x\). Then
\(N_Q(x)=\{\sum_{j\in J(x)}c_ju_j:c_j\ge0\}\). In particular
\(N_Q(x)=\{0\}\) at interior points, \(N_Q(x)=\{cu_j:c\ge0\}\) for
\(x\in\relint F_j\), and \(N_Q(x)=\{c_ju_j+c_ku_k:c_j,c_k\ge0\}\), a sector of
angle \(\alpha_{jk}\) in the plane spanned by \(u_j\) and \(u_k\), for
\(x\in\relint E_{jk}\).
\end{fact}

\begin{proof}
Put \(T=\{d:u_j\cdot d\le0\text{ for }j\in J(x)\}\). A vector \(v\) lies in
\(N_Q(x)\) if and only if \(v\cdot d\le0\) for all \(d\in T\): if \(d\in T\)
then \(x+\varepsilon d\in Q\) for small \(\varepsilon>0\), because the
inactive inequalities are strict at \(x\), and conversely \(q-x\in T\) for
\(q\in Q\). By Farkas's lemma, which says that \(v\cdot d\le0\) for all
\(d\) with \(u_j\cdot d\le0\) (\(j\in J\)) if and only if \(v\) is a
nonnegative combination of the \(u_j\), \(j\in J\), this is the stated
cone. For \(x\in\relint F_j\), \(J(x)=\{j\}\), since an index \(k\ne j\) in
\(J(x)\) would put \(x\) in \(F_j\cap F_k\), a face of \(F_j\) of dimension
at most \(n-2\), and so in the relative boundary of \(F_j\). For
\(x\in\relint E_{jk}\), \(J(x)=\{j,k\}\) by Fact~\ref{tb:facets}(d).
\end{proof}

\usedin{\PR{Lemma}{appx-an:lem:surface}{C.5}, Step~2 (``the normal cone at a
relative-interior point of the ridge shared by facets \(j\) and \(k\) is the
cone spanned by \(u_j\) and \(u_k\)'') and the definition of the outward
tubes.}
\emph{Source.} \TBc{Roc70}, Sections~22 and~23, for Farkas's lemma and normal
cones; \TBc{RW98}, Chapter~6.

\begin{fact}[Nearest point]\label{tb:projection}
Let \(Q\subseteq\R^n\) be nonempty, closed, and convex. For each
\(y\in\R^n\) there is a unique point \(\prj_Q(y)\in Q\) with
\(\norm{y-\prj_Q(y)}=\dist(y,Q)\), and \(x=\prj_Q(y)\) if and only if
\(x\in Q\) and \(y-x\in N_Q(x)\). Consequently every \(y\in\R^n\) can be
written in exactly one way as \(y=x+v\) with \(x\in Q\) and \(v\in N_Q(x)\),
and then \(\norm v=\dist(y,Q)\). A direction \(v\) lies in \(N_Q(x)\) if and
only if \(x\) is the nearest point of \(x+tv\) in \(Q\) for every \(t>0\).
\end{fact}

\begin{proof}
For \(x,q\in Q\), \(\norm{y-q}^2=\norm{y-x}^2+\norm{x-q}^2-2(y-x)\cdot(q-x)\).
If \(y-x\in N_Q(x)\), the last term is nonnegative, so
\(\norm{y-q}^2\ge\norm{y-x}^2+\norm{x-q}^2>\norm{y-x}^2\) for \(q\ne x\); thus
\(x\) is the unique nearest point. Conversely, if \((y-x)\cdot(q-x)>0\) for
some \(q\in Q\), then moving from \(x\) toward \(q\) decreases the distance
to \(y\) for small steps, by the displayed identity applied to
\(x+\varepsilon(q-x)\in Q\), so \(x\) is not nearest. A nearest point exists
because \(Q\) is closed. The last statement follows because \(N_Q(x)\) is a
cone.
\end{proof}

\usedin{\PR{Lemma}{appx-an:lem:surface}{C.5}, Step~2, to show that the tubes
are disjoint (``their points have distinct metric projections'').}
\emph{Source.} \TBc{RW98}, Chapter~6 (Moreau's decomposition); \TBc{GZ},
proof of Lemma~28.

\subsection{Surface measure, Gaussian surface integrals, and the
divergence theorem}\label{tb:sec:surface}

\begin{fact}[Hausdorff measure on flat pieces]\label{tb:hausdorff}
For \(0\le k\le n\), the \(k\)-dimensional Hausdorff measure \(\Haus^k\) on
\(\R^n\) is the Borel measure
\[
 \Haus^k(S)=\lim_{\epsilon\downarrow0}\inf\Bigl\{\sum_i\omega_k
 \Bigl(\tfrac12\operatorname{diam}C_i\Bigr)^k:S\subseteq\bigcup_iC_i,\
 \operatorname{diam}C_i\le\epsilon\Bigr\},
\]
where \(\omega_k\) is the volume of the unit ball of \(\R^k\). (a) On an
affine subspace of dimension \(k\), \(\Haus^k\) is \(k\)-dimensional Lebesgue
measure; in particular \(\Haus^0\) is counting measure and \(\Haus^n\) is
Lebesgue measure on \(\R^n\). (b) \(\Haus^k\) is invariant under isometries.
(c) A set contained in a finite union of affine subspaces of dimension less
than \(k\) has \(\Haus^k\)-measure zero.
\end{fact}

Part (c) follows from (a), since each such subspace lies in a
\(k\)-dimensional affine subspace in which it is a Lebesgue null set. For a
proper, full-dimensional polyhedron \(Q\) with an irredundant description,
and for a function \(f\ge0\) or \(f\) integrable on \(\partial Q\), the
\emph{surface integral} is
\[
 \int_{\partial Q}f\,d\Haus^{n-1}=\sum_{j=1}^m\int_{F_j}f\,d\Haus^{n-1},
\]
where on each facet \(\Haus^{n-1}\) is Lebesgue measure on the hyperplane
\(H_j\) by part~(a). By Fact~\ref{tb:facets}(b) and part~(c), distinct
facets overlap in sets of \(\Haus^{n-1}\)-measure zero, so the sum is the
integral of \(f\) over the set \(\partial Q\) with respect to
\(\Haus^{n-1}\). It depends on the set \(Q\) and not on its description,
which justifies deleting redundant inequalities before computing it.
The Gaussian surface integral \(\surf(Q)=\int_{\partial Q}\Hh\,\varphi_r\,
d\Haus^{r-1}\) of Subsection~\PN{sec:outline-late}{2.4} is the case
\(n=r\) and \(f=\Hh\varphi_r\), with \(\Hh(g)=\norm g^2-r\).

\usedin{the definition of \(\surf\) in Subsection~\PN{sec:outline-late}{2.4}
and in \PR{Lemma}{appx-an:lem:surface}{C.5}.}
\emph{Source.} \TBc{EG15}, Chapter~2.

\begin{fact}[Gauss--Green formula (divergence theorem) for a polyhedron]\label{tb:divergence}
Let \(Q\subseteq\R^n\) (\(n\ge1\)) be a proper, full-dimensional polyhedron
with an irredundant description, facets \(F_1,\ldots,F_m\), and outward unit
normals \(u_1,\ldots,u_m\). Let \(X:\R^n\to\R^n\) be of class \(C^1\), with
\(\norm{X(x)}+|\operatorname{div}X(x)|\le C(1+\norm x)^k\varphi_n(x)\) for
some \(C\) and \(k\). Then
\[
 \int_Q\operatorname{div}X\,dx=\sum_{j=1}^m\int_{F_j}X\cdot u_j\,d\Haus^{n-1}.
\]
For \(n=1\) the facets are the endpoints of an interval, \(\Haus^0\) is
counting measure, and the identity is the fundamental theorem of calculus.
\end{fact}

\begin{proof}
It suffices to show, for a unit vector \(e\) and a function
\(f\in C^1(\R^n)\) with \(|f|+|\partial_ef|\le C(1+\norm x)^k\varphi_n(x)\),
that
\begin{equation}
 \int_Q\partial_ef\,dx=\sum_j(e\cdot u_j)\int_{F_j}f\,d\Haus^{n-1},
 \tag{\(*\)}
\end{equation}
since the theorem follows by taking \(e=e_i\) and \(f=X_i\) and summing over
\(i\). Write \(x=w+te\) with \(w\in e^\perp\). By Fubini's theorem,
\(\int_Q\partial_ef\,dx=\int_{e^\perp}\int_{I(w)}\partial_tf(w+te)\,dt\,dw\),
where \(I(w)=\{t:w+te\in Q\}\) is a closed interval, possibly empty, a point,
or unbounded, because \(Q\) is convex. Since \(f\) tends to zero along every
line, the inner integral equals \(f(w+t_+(w)e)-f(w+t_-(w)e)\), where
\(t_\pm(w)\) are the endpoints of \(I(w)\) and a term with an infinite
endpoint is read as zero.

Consider the upper endpoints \(x=w+t_+(w)e\). If \(x\in\relint F_j\), then
only the inequality \(j\) is active at \(x\) (Fact~\ref{tb:normalcone}), so
\(x+te\in Q\) for small \(t>0\) unless \(e\cdot u_j>0\), and \(x\) is an upper
endpoint if and only if \(e\cdot u_j>0\). The upper endpoints that lie in no
\(\relint F_j\) lie in \(\bigcup_{j\ne k}F_j\cap F_k\), which is a subset of a
finite union of affine subspaces of dimension at most \(n-2\)
(Fact~\ref{tb:facets}(b)); its projection along \(e\) onto \(e^\perp\) is a
Lebesgue null set of \(e^\perp\). For fixed \(j\) with \(e\cdot u_j>0\), the
map \(x\mapsto x-(x\cdot e)e\) is injective on the set of upper endpoints in
\(\relint F_j\), since a line meets \(Q\) in one interval, and it is the
restriction to \(H_j\) of a linear map whose Jacobian with respect to
\(\Haus^{n-1}\) on \(H_j\) is \(e\cdot u_j\): if \(b_1,\ldots,b_{n-1}\) is an
orthonormal basis of \(u_j^\perp\), the volume of the projection along \(e\)
of the cube they span equals the volume of the parallelepiped spanned by
\(b_1,\ldots,b_{n-1},e\), which is \(|e\cdot u_j|\). Therefore
\(\int f(w+t_+(w)e)\,dw=\sum_{j:\,e\cdot u_j>0}(e\cdot u_j)\int_{F_j}f\,d\Haus^{n-1}\),
the sets \(F_j\setminus\relint F_j\) being \(\Haus^{n-1}\)-null. The same
argument for the lower endpoints gives
\(-\sum_{j:\,e\cdot u_j<0}|e\cdot u_j|\int_{F_j}f\,d\Haus^{n-1}\), and
subtracting proves \((*)\); facets with \(e\cdot u_j=0\) contribute nothing
on either side.
\end{proof}

\usedin{\PR{Lemma}{appx-an:lem:surface}{C.5}, Step~1, applied in the hyperplane
\(H_j\) to the facet \(F_j\), regarded as a polyhedron of dimension
\(n=r-1\), whose facets are the ridges \(E_{jk}\) with outward normals
\(\nu_{k\mid j}\) (Fact~\ref{tb:facets}(c)), and to the vector field
\(X(y)=y\varphi_{r-1}(y)\), for which
\(\operatorname{div}X(y)=\{(r-1)-\norm y^2\}\varphi_{r-1}(y)\).
The decay of \(X\) is what replaces the truncation of unbounded facets.}
\emph{Source.} \TBc{GZ}, Lemma~24, proved there by truncation to a ball and
the Gauss--Green formula for Lipschitz domains~\TBc{EG15}, Chapters~4
and~5.

\subsection{Tubes around a polyhedron}\label{tb:sec:tubes}

Let \(Q\) be a proper, full-dimensional polyhedron in \(\R^n\) with an
irredundant description, facets \(F_j\), and ridges \(E_{jk}\) (\(n\ge2\)). For
a ridge, \(L_{jk}=\operatorname{span}(u_j,u_k)\), and \(z_{jk}\) is the unique
point of \(L_{jk}\) with \(u_j\cdot z_{jk}=a_j\) and \(u_k\cdot z_{jk}=a_k\). Every
point \(x\) of \(E_{jk}\) satisfies these two equations, so
\(E_{jk}\subseteq z_{jk}+L_{jk}^\perp\), and \(z_{jk}\) is the orthogonal
projection of every point of the ridge onto \(L_{jk}\). For \(\varepsilon>0\),
the \emph{outward tubes} of length \(\varepsilon\) are
\begin{align*}
 \mathcal T_j(\varepsilon)&=\{x+tu_j:x\in\relint F_j,\ 0<t<\varepsilon\},\\
 \mathcal T_{jk}(\varepsilon)&=\{x+v:x\in\relint E_{jk},\ v\in N_Q(x),\
 0<\norm v<\varepsilon\}.
\end{align*}
By Fact~\ref{tb:normalcone}, \(v\) ranges over the sector
\(S_{jk}(\varepsilon)=\{c_ju_j+c_ku_k:c_j,c_k\ge0\}\cap
\{\norm v<\varepsilon\}\) of angle \(\alpha_{jk}\) and area
\(\alpha_{jk}\varepsilon^2/2\) in \(L_{jk}\), without its apex. The tubes over the closed facets and ridges that
the proofs describe differ from these by sets of Lebesgue measure zero.
For \(Q(a)=\bigcap_j\{u_j\cdot x\le a_j\}\), write
\(Q(a)^t=Q(a_1+t,\ldots,a_m+t)\), the polyhedron with every offset raised by \(t\).

\begin{fact}[Tubes: disjointness and Gaussian mass]\label{tb:tubes}
(a) The sets \(\mathcal T_j(\varepsilon)\), \(j\le m\), and
\(\mathcal T_{jk}(\varepsilon)\), for the ridges \(E_{jk}\), are pairwise
disjoint, and they are contained in
\(\{y:0<\dist(y,Q)<\varepsilon\}\subseteq Q^\varepsilon\setminus Q\), where
\(Q^\varepsilon=Q(a)^\varepsilon\).

(b) With \(s_j=\int_{F_j}\varphi_n\,d\Haus^{n-1}\) and
\(\rho_{jk}=\int_{E_{jk}}\varphi_n\,d\Haus^{n-2}\),
\[
 \Phi_n(\mathcal T_j(\varepsilon))=\int_{F_j}\int_0^\varepsilon
 \varphi_n(x+tu_j)\,dt\,d\Haus^{n-1}(x),\qquad
 \Phi_n(\mathcal T_{jk}(\varepsilon))=\int_{E_{jk}}\int_{S_{jk}(\varepsilon)}
 \varphi_n(x+v)\,dv\,d\Haus^{n-2}(x).
\]

(c) Let \(R>0\) and \(\varepsilon=1/(1+R)\). If \(|a_j|\le R\), then
\(\Phi_n(\mathcal T_j(\varepsilon))\ge e^{-3/2}\,\varepsilon\,s_j\). If
\(\norm{z_{jk}}\le R\), then
\(\Phi_n(\mathcal T_{jk}(\varepsilon))\ge e^{-3/2}\,(\alpha_{jk}\varepsilon^2/2)\,
\rho_{jk}\).
\end{fact}

\begin{proof}
(a) By Fact~\ref{tb:projection}, a point \(y\notin Q\) is \(x+v\) for exactly
one pair with \(x=\prj_Q(y)\in Q\) and \(v\in N_Q(x)\setminus\{0\}\), and
\(\norm v=\dist(y,Q)\). A point of \(\mathcal T_j(\varepsilon)\) has its nearest
point in \(\relint F_j\), and a point of \(\mathcal T_{jk}(\varepsilon)\) has
its nearest point in \(\relint E_{jk}\), by Fact~\ref{tb:normalcone}. The
sets \(\relint F_j\) and \(\relint E_{jk}\) are pairwise disjoint, being
relative interiors of distinct faces, so the tubes are disjoint. If
\(\dist(y,Q)<\varepsilon\) and \(x=\prj_Q(y)\), then
\(u_j\cdot y\le u_j\cdot x+\norm{y-x}<a_j+\varepsilon\) for all \(j\), since the
\(u_j\) are unit vectors; thus \(y\in Q^\varepsilon\).

(b) The map \((x,t)\mapsto x+tu_j\) from \(H_j\times\R\) to \(\R^n\) is an
isometry, and \((x,v)\mapsto x+v\) is an isometry from
\((z_{jk}+L_{jk}^\perp)\times L_{jk}\) to \(\R^n\). Lebesgue measure on
\(\R^n\) is therefore the product of \(\Haus^{n-1}\) on \(H_j\) and \(dt\), and
of \(\Haus^{n-2}\) on \(z_{jk}+L_{jk}^\perp\) and \(dv\) on \(L_{jk}\)
(Fact~\ref{tb:hausdorff}(a)). Apply Fact~\ref{tb:fubini}.

(c) For \(x\in F_j\), \(x\cdot u_j=a_j\), and by Fact~\ref{tb:gauss-coord}(b)
\(\varphi_n(x+tu_j)=\varphi_n(x)\exp\{-ta_j-t^2/2\}\ge e^{-R\varepsilon-
\varepsilon^2/2}\varphi_n(x)\ge e^{-3/2}\varphi_n(x)\), since
\(R\varepsilon<1\) and \(\varepsilon<1\). Integrating over
\(t\in(0,\varepsilon)\) gives the first bound. For \(x\in E_{jk}\) and
\(v\in L_{jk}\), \(x\cdot v=z_{jk}\cdot v\), because \(x-z_{jk}\perp L_{jk}\),
so \(\varphi_n(x+v)=\varphi_n(x)\exp\{-z_{jk}\cdot v-\norm v^2/2\}
\ge e^{-3/2}\varphi_n(x)\) for \(\norm v<\varepsilon\). Integrating over the
sector, whose area is \(\alpha_{jk}\varepsilon^2/2\), gives the second bound.
\end{proof}

In the plane the tubes over the sides are strips and the tubes over the
corners are circular sectors, and together they tile the outer
\(\varepsilon\)-neighborhood of a polygon. When \(n\ge3\) the faces of
dimension at most \(n-3\) carry further strata of the neighborhood, which the
proofs do not need.

\usedin{\PR{Lemma}{appx-an:lem:surface}{C.5}, Step~2, for the bounds
\(q_j\ge c_*\varepsilon_Rs_j\) and \(q_{jk}\ge c_*\alpha_{jk}\varepsilon_R^2\rho_{jk}\),
and Step~3, for the disjointness that allows the Gaussian masses of the
tubes in a dyadic shell to be added, and for the inclusion
\(\mathcal O_\iota(V_0)\subseteq(Q')^{\varepsilon}\setminus Q'\) from which the
event \(V_0+\sigma G\in Q^{\sigma\varepsilon}\setminus Q\) is read: the
polyhedron \(Q'=(Q-V_0)/\sigma\) has offsets \((a_j-u_j\cdot V_0)/\sigma\),
and raising them by \(\varepsilon\) is raising those of \(Q\) by
\(\sigma\varepsilon\).}
\emph{Source.} \TBc{GZ}, Lemmas~28 and~29.

\subsection{The coarea formula and the first variation of a Gaussian
integral}\label{tb:sec:coarea}

\begin{fact}[Coarea formula]\label{tb:coarea}
Let \(\psi:\R^n\to\R\) be Lipschitz, so that \(\nabla\psi\) exists almost
everywhere. For every measurable \(f\ge0\),
\[
 \int_{\R^n}f\,\norm{\nabla\psi}\,dx=\int_{-\infty}^\infty
 \Bigl(\int_{\{\psi=s\}}f\,d\Haus^{n-1}\Bigr)ds ,
\]
and the same holds for integrable \(f\).
\end{fact}

\emph{Source.} \TBc{EG15}, Section~3.4. The proofs use the formula only for
a maximum of finitely many affine functions with gradients of equal length, and
for these Fact~\ref{tb:variation} gives a direct proof from Fubini's theorem.

\begin{fact}[First variation of a polyhedral sublevel integral]\label{tb:variation}
Let \(u_1,\ldots,u_m\) be unit vectors in \(\R^n\) (\(m\ge1\)), \(\kappa>0\),
\(b_1,\ldots,b_m\in\R\), and
\[
 \psi(g)=\max_{j\le m}\{\kappa\,u_j\cdot g-b_j\},\qquad Q_s=\{g:\psi(g)\le s\}.
\]
(a) For \(f\ge0\) measurable, or \(f\) integrable, and \(s\in\R\),
\[
 \int_{Q_s}f\,dg=\kappa^{-1}\int_{-\infty}^s\Bigl(\int_{\{\psi=s'\}}f\,d\Haus^{n-1}
 \Bigr)ds'.
\]
(b) If \(f\) is integrable, then \(\Lambda(s)=\int_{Q_s}f\,dg\) is absolutely
continuous, and for almost every \(s\)
\[
 \Lambda'(s)=\kappa^{-1}\int_{\{\psi=s\}}f\,d\Haus^{n-1},\qquad
 |\Lambda'(s)|\le\kappa^{-1}\sum_{j=1}^m\int_{\{\kappa u_j\cdot g-b_j=s\}}|f|\,d\Haus^{n-1}.
\]
(c) If \(s>\inf\psi\), then \(Q_s\) is a proper, full-dimensional polyhedron
with normals \(u_j\) and offsets \((s+b_j)/\kappa\), \(\{\psi=s\}=\partial Q_s\),
and \(\int_{\{\psi=s\}}f\,d\Haus^{n-1}\) is the surface integral of \(f\) over
\(\partial Q_s\) of Subsection~\ref{tb:sec:surface}. For \(s<\inf\psi\) the set
\(Q_s\) is empty, and \(s=\inf\psi\) is the only other level at which \(Q_s\)
may fail to be full-dimensional.

(d) If \(\Lambda\) is Lipschitz and \(S\) is a real random variable with a
density, then \(s\mapsto\E\Lambda(s+S)\) is differentiable with derivative
\(\E\Lambda'(s+S)\), where \(\Lambda'\) is defined almost everywhere.
\end{fact}

\begin{proof}
(a) Let \(C_j\) be the set where \(j\) is the smallest index attaining the
maximum defining \(\psi\); the Borel sets \(C_1,\ldots,C_m\) partition
\(\R^n\). On \(C_j\), \(\psi(g)=\kappa t-b_j\) in the coordinates
\(g=tu_j+y\), \(y\perp u_j\). By Fact~\ref{tb:fubini}, in these coordinates,
\[
 \int_{C_j\cap Q_s}f\,dg=\int_{u_j^\perp}\int_{-\infty}^{(s+b_j)/\kappa}
 (f\ind_{C_j})(tu_j+y)\,dt\,dy .
\]
The substitution \(s'=\kappa t-b_j\) turns the inner integral into
\(\kappa^{-1}\int_{-\infty}^s(f\ind_{C_j})(tu_j+y)\,ds'\) with
\(t=(s'+b_j)/\kappa\), and a further application of Fubini's theorem
exchanges the integrals. For fixed \(s'\), the set of \(tu_j+y\) in \(C_j\)
with this \(t\) is the slice \(C_j\cap\{\psi=s'\}\), and on the hyperplane
containing it \(\Haus^{n-1}\) is Lebesgue measure \(dy\)
(Fact~\ref{tb:hausdorff}(a)). Hence
\(\int_{C_j\cap Q_s}f\,dg=\kappa^{-1}\int_{-\infty}^s\int_{C_j\cap\{\psi=s'\}}
f\,d\Haus^{n-1}\,ds'\). Summing over \(j\) gives (a), since the slices
\(C_j\cap\{\psi=s'\}\) partition \(\{\psi=s'\}\).

(b) Applying (a) to \(|f|\) and letting \(s\to\infty\) shows that
\(I(s')=\kappa^{-1}\int_{\{\psi=s'\}}f\,d\Haus^{n-1}\) is integrable on
\(\R\), and (a) says \(\Lambda(s)=\int_{-\infty}^sI(s')\,ds'\). Hence \(\Lambda\) is
absolutely continuous, and \(\Lambda'=I\) almost everywhere by Lebesgue's
differentiation theorem. The bound follows because
\(\{\psi=s\}\subseteq\bigcup_j\{\kappa u_j\cdot g-b_j=s\}\), a union of \(m\)
hyperplanes.

(c) The set \(Q_s\) is the polyhedron \(\bigcap_j\{u_j\cdot g\le(s+b_j)/\kappa\}\),
and \(\psi\) is convex and continuous. If \(\psi(g)=s>\inf\psi\), choose \(g'\)
with \(\psi(g')<s\); for \(\tau>0\), \(g=(g_\tau+\tau g')/(1+\tau)\) with
\(g_\tau=g+\tau(g-g')\), so convexity gives
\(\psi(g_\tau)\ge s+\tau(s-\psi(g'))>s\). Thus \(g\) is a limit of points
outside \(Q_s\) and \(g\in\partial Q_s\); the open set \(\{\psi<s\}\) is
nonempty and lies in \(Q_s\), which is therefore full-dimensional, and it
is a proper subset because \(\psi(g)\ge\kappa u_1\cdot g-b_1\to\infty\) along
\(u_1\). A point with \(\psi<s\) is interior, so \(\partial Q_s=\{\psi=s\}\). After
the redundant and repeated inequalities are deleted, \(\partial Q_s=\{\psi=s\}\)
is the union of the sets \(Q_s\cap\{\kappa u_j\cdot g-b_j=s\}\) of the
remaining inequalities, which are faces of \(Q_s\) or empty. A face of
dimension \(n-1\) is a facet (Fact~\ref{tb:facets}(a)) and any other face is
\(\Haus^{n-1}\)-null, so the integral over \(\{\psi=s\}\) is the sum over the
facets of the irredundant description.

(d) A Lipschitz function is absolutely continuous and differentiable outside a
Lebesgue null set \(\mathcal N\). The random variable \(s+S\) avoids
\(\mathcal N\) almost surely, because \(S\) has a density, so the difference
quotients of \(\Lambda(s+S)\) converge almost surely and are bounded by the
Lipschitz constant. Dominated convergence gives the claim.
\end{proof}

\usedin{the paragraph ``First variation and common shifts'' after
\PR{Lemma}{appx-an:lem:surface}{C.5}, with
\(\psi(g)=\max_j\{u_j\cdot(V_0+\sigma g)-a_j\}\), \(\kappa=\sigma\),
\(b_j=a_j-u_j\cdot V_0\), \(Q_s=(Q(a+s\mathbf 1)-V_0)/\sigma\), and
\(f=\Hh\varphi_r\), which gives the derivative
\(\sigma^{-1}\surf((Q(a+s\mathbf 1)-V_0)/\sigma)\) for almost every \(s\) (parts
(a)--(c)); the Lipschitz bound in (b) is the bound \(Cm(r+1)/\sigma\) of that
paragraph, and part (d) justifies averaging the derivative over the
absolutely continuous common shift \(-h\log T_t\) in Step~4 of the proof of
\PR{Lemma}{appx-an:lem:matrix}{C.4}.}

\subsection{Where each fact is used}\label{tb:sec:index}

The table lists, for each fact of this appendix, the places in the proofs
that use it.

\medskip
\begin{center}
\small
\renewcommand{\arraystretch}{1.15}
\begin{tabular}{@{}p{0.27\textwidth}p{0.68\textwidth}@{}}
Facts~\ref{tb:fubini}, \ref{tb:taylor} &
every conditioning on an independent component in Appendix~\PN{appx-sec:support}{C};
the Taylor expansions in \PR{Lemma}{appx-an:lem:identity}{C.1}; Step~3 of
\PR{Lemma}{appx-an:lem:surface}{C.5}\\
Facts~\ref{tb:faa}, \ref{tb:logistic} &
the derivative formula and the local envelope in Subsection~\PN{appx-sec:mixture}{B.2};
Step~1 of \PR{Lemma}{appx-an:lem:restoration}{C.2}\\
Facts~\ref{tb:gauss-coord}--\ref{tb:heat} &
Subsection~\PN{sec:outline-path}{2.2} and \PR{Lemma}{appx-an:lem:identity}{C.1} (heat equation);
Step~1 of \PR{Lemma}{appx-an:lem:restoration}{C.2} (integration by parts);
\PR{Lemma}{appx-lem:endpoint-gap}{B.3}; Steps~1--2 of \PR{Lemma}{appx-an:lem:surface}{C.5}\\
Facts~\ref{tb:nuclear}, \ref{tb:logsum} &
Subsection~\PN{sec:outline-late}{2.4}; Step~4 of \PR{Lemma}{appx-an:lem:restoration}{C.2};
Steps~1 and~3 of \PR{Lemma}{appx-an:lem:matrix}{C.4}\\
Facts~\ref{tb:prekopa}--\ref{tb:lc-weights} &
Subsection~\PN{sec:outline-late}{2.4}; Step~1 of \PR{Lemma}{appx-an:lem:matrix}{C.4}\\
Facts~\ref{tb:facets}--\ref{tb:projection} &
the geometric paragraph and Steps~1--2 of \PR{Lemma}{appx-an:lem:surface}{C.5}\\
Facts~\ref{tb:hausdorff}, \ref{tb:divergence} &
the definition of \(\surf\); Step~1 of \PR{Lemma}{appx-an:lem:surface}{C.5}\\
Fact~\ref{tb:tubes} &
Steps~2 and~3 of \PR{Lemma}{appx-an:lem:surface}{C.5}\\
Facts~\ref{tb:coarea}, \ref{tb:variation} &
the paragraph on first variation after \PR{Lemma}{appx-an:lem:surface}{C.5}; Step~4 of
\PR{Lemma}{appx-an:lem:matrix}{C.4}\\
\end{tabular}
\end{center}

\subsection{Sources}\label{tb:sec:sources}

\begin{description}\small\setlength{\itemsep}{0pt}
\item[{[BL76]}] H.~J. Brascamp and E.~H. Lieb (1976). On extensions of the
Brunn--Minkowski and Pr\'ekopa--Leindler theorems, including inequalities
for log concave functions, and with an application to the diffusion
equation. \emph{Journal of Functional Analysis} \textbf{22}, 366--389.
\item[{[Com74]}] L. Comtet (1974). \emph{Advanced Combinatorics}. Reidel,
Dordrecht.
\item[{[CT06]}] T.~M. Cover and J.~A. Thomas (2006). \emph{Elements of
Information Theory}, 2nd edition. Wiley.
\item[{[EG15]}] L.~C. Evans and R.~F. Gariepy (2015). \emph{Measure Theory
and Fine Properties of Functions}, revised edition. CRC Press.
\item[{[Fol99]}] G.~B. Folland (1999). \emph{Real Analysis: Modern
Techniques and Their Applications}, 2nd edition. Wiley.
\item[{[Gar02]}] R.~J. Gardner (2002). The Brunn--Minkowski inequality.
\emph{Bulletin of the American Mathematical Society} \textbf{39}, 355--405.
\item[{[GZ]}] Z. Gao and W. Zhang (2026). Cubic-root Gaussian approximation
under unrestricted covariance. arXiv:2608.30221v1.
\item[{[Lei72]}] L. Leindler (1972). On a certain converse of H\"older's
inequality II. \emph{Acta Scientiarum Mathematicarum (Szeged)} \textbf{33},
217--223.
\item[{[Pre71]}] A. Pr\'ekopa (1971). Logarithmic concave measures with
application to stochastic programming. \emph{Acta Scientiarum Mathematicarum
(Szeged)} \textbf{32}, 301--316.
\item[{[Pre73]}] A. Pr\'ekopa (1973). On logarithmic concave measures and
functions. \emph{Acta Scientiarum Mathematicarum (Szeged)} \textbf{34},
335--343.
\item[{[Roc70]}] R.~T. Rockafellar (1970). \emph{Convex Analysis}. Princeton
University Press.
\item[{[RW98]}] R.~T. Rockafellar and R.~J.-B. Wets (1998). \emph{Variational
Analysis}. Springer, Berlin.
\item[{[Ste81]}] C.~M. Stein (1981). Estimation of the mean of a multivariate
normal distribution. \emph{The Annals of Statistics} \textbf{9}, 1135--1151.
\item[{[Zie95]}] G.~M. Ziegler (1995). \emph{Lectures on Polytopes}.
Springer, New York.
\end{description}
